% Options for packages loaded elsewhere
\PassOptionsToPackage{unicode}{hyperref}
\PassOptionsToPackage{hyphens}{url}
\documentclass[
]{article}
\usepackage{xcolor}
\usepackage{amsmath,amssymb}
\setcounter{secnumdepth}{-\maxdimen} % remove section numbering
\usepackage{iftex}
\ifPDFTeX
  \usepackage[T1]{fontenc}
  \usepackage[utf8]{inputenc}
  \usepackage{textcomp} % provide euro and other symbols
\else % if luatex or xetex
  \usepackage{unicode-math} % this also loads fontspec
  \defaultfontfeatures{Scale=MatchLowercase}
  \defaultfontfeatures[\rmfamily]{Ligatures=TeX,Scale=1}
\fi
\usepackage{lmodern}
\ifPDFTeX\else
  % xetex/luatex font selection
\fi
% Use upquote if available, for straight quotes in verbatim environments
\IfFileExists{upquote.sty}{\usepackage{upquote}}{}
\IfFileExists{microtype.sty}{% use microtype if available
  \usepackage[]{microtype}
  \UseMicrotypeSet[protrusion]{basicmath} % disable protrusion for tt fonts
}{}
\makeatletter
\@ifundefined{KOMAClassName}{% if non-KOMA class
  \IfFileExists{parskip.sty}{%
    \usepackage{parskip}
  }{% else
    \setlength{\parindent}{0pt}
    \setlength{\parskip}{6pt plus 2pt minus 1pt}}
}{% if KOMA class
  \KOMAoptions{parskip=half}}
\makeatother
\usepackage{color}
\usepackage{fancyvrb}
\newcommand{\VerbBar}{|}
\newcommand{\VERB}{\Verb[commandchars=\\\{\}]}
\DefineVerbatimEnvironment{Highlighting}{Verbatim}{commandchars=\\\{\}}
% Add ',fontsize=\small' for more characters per line
\newenvironment{Shaded}{}{}
\newcommand{\AlertTok}[1]{\textcolor[rgb]{1.00,0.00,0.00}{\textbf{#1}}}
\newcommand{\AnnotationTok}[1]{\textcolor[rgb]{0.38,0.63,0.69}{\textbf{\textit{#1}}}}
\newcommand{\AttributeTok}[1]{\textcolor[rgb]{0.49,0.56,0.16}{#1}}
\newcommand{\BaseNTok}[1]{\textcolor[rgb]{0.25,0.63,0.44}{#1}}
\newcommand{\BuiltInTok}[1]{\textcolor[rgb]{0.00,0.50,0.00}{#1}}
\newcommand{\CharTok}[1]{\textcolor[rgb]{0.25,0.44,0.63}{#1}}
\newcommand{\CommentTok}[1]{\textcolor[rgb]{0.38,0.63,0.69}{\textit{#1}}}
\newcommand{\CommentVarTok}[1]{\textcolor[rgb]{0.38,0.63,0.69}{\textbf{\textit{#1}}}}
\newcommand{\ConstantTok}[1]{\textcolor[rgb]{0.53,0.00,0.00}{#1}}
\newcommand{\ControlFlowTok}[1]{\textcolor[rgb]{0.00,0.44,0.13}{\textbf{#1}}}
\newcommand{\DataTypeTok}[1]{\textcolor[rgb]{0.56,0.13,0.00}{#1}}
\newcommand{\DecValTok}[1]{\textcolor[rgb]{0.25,0.63,0.44}{#1}}
\newcommand{\DocumentationTok}[1]{\textcolor[rgb]{0.73,0.13,0.13}{\textit{#1}}}
\newcommand{\ErrorTok}[1]{\textcolor[rgb]{1.00,0.00,0.00}{\textbf{#1}}}
\newcommand{\ExtensionTok}[1]{#1}
\newcommand{\FloatTok}[1]{\textcolor[rgb]{0.25,0.63,0.44}{#1}}
\newcommand{\FunctionTok}[1]{\textcolor[rgb]{0.02,0.16,0.49}{#1}}
\newcommand{\ImportTok}[1]{\textcolor[rgb]{0.00,0.50,0.00}{\textbf{#1}}}
\newcommand{\InformationTok}[1]{\textcolor[rgb]{0.38,0.63,0.69}{\textbf{\textit{#1}}}}
\newcommand{\KeywordTok}[1]{\textcolor[rgb]{0.00,0.44,0.13}{\textbf{#1}}}
\newcommand{\NormalTok}[1]{#1}
\newcommand{\OperatorTok}[1]{\textcolor[rgb]{0.40,0.40,0.40}{#1}}
\newcommand{\OtherTok}[1]{\textcolor[rgb]{0.00,0.44,0.13}{#1}}
\newcommand{\PreprocessorTok}[1]{\textcolor[rgb]{0.74,0.48,0.00}{#1}}
\newcommand{\RegionMarkerTok}[1]{#1}
\newcommand{\SpecialCharTok}[1]{\textcolor[rgb]{0.25,0.44,0.63}{#1}}
\newcommand{\SpecialStringTok}[1]{\textcolor[rgb]{0.73,0.40,0.53}{#1}}
\newcommand{\StringTok}[1]{\textcolor[rgb]{0.25,0.44,0.63}{#1}}
\newcommand{\VariableTok}[1]{\textcolor[rgb]{0.10,0.09,0.49}{#1}}
\newcommand{\VerbatimStringTok}[1]{\textcolor[rgb]{0.25,0.44,0.63}{#1}}
\newcommand{\WarningTok}[1]{\textcolor[rgb]{0.38,0.63,0.69}{\textbf{\textit{#1}}}}
\usepackage{longtable,booktabs,array}
\newcounter{none} % for unnumbered tables
\usepackage{calc} % for calculating minipage widths
% Correct order of tables after \paragraph or \subparagraph
\usepackage{etoolbox}
\makeatletter
\patchcmd\longtable{\par}{\if@noskipsec\mbox{}\fi\par}{}{}
\makeatother
% Allow footnotes in longtable head/foot
\IfFileExists{footnotehyper.sty}{\usepackage{footnotehyper}}{\usepackage{footnote}}
\makesavenoteenv{longtable}
\setlength{\emergencystretch}{3em} % prevent overfull lines
\providecommand{\tightlist}{%
  \setlength{\itemsep}{0pt}\setlength{\parskip}{0pt}}
\usepackage{bookmark}
\IfFileExists{xurl.sty}{\usepackage{xurl}}{} % add URL line breaks if available
\urlstyle{same}
\hypersetup{
  hidelinks,
  pdfcreator={LaTeX via pandoc}}

\author{}
\date{}

\begin{document}

\section{Decomposed Prompting: A MODULAR APPROACH FOR SOLVING COMPLEX
TASKS}\label{decomposed-prompting-a-modular-approach-for-solving-complex-tasks}

Tushar Khot♣, Harsh Trivedi♥, Matthew Finlayson♣, Yao Fu♠, Kyle
Richardson♣, Peter Clark♣, Ashish Sabharwal♣

♣Allen Institute for AI ♥Stony Brook University ♠University of Edinburgh
tushark@allenai.org, hjtrivedi@cs.stonybrook.edu, matthewf@allenai.org,
yao.fu@ed.ac.uk, kyler@allenai.org, peterc@allenai.org,
ashishs@allenai.org

\section{ABSTRACT}\label{abstract}

Few-shot prompting is a surprisingly powerful way to use Large Language
Models (LLMs) to solve various tasks. However, this approach struggles
as the task complexity increases or when the individual reasoning steps
of the task themselves are hard to learn, especially when embedded in
more complex tasks. To address this, we propose Decomposed Prompting, a
new approach to solve complex tasks by decomposing them (via prompting)
into simpler sub-tasks that can be delegated to a shared library of
prompting-based LLMs dedicated to these sub-tasks. This modular
structure allows each prompt to be optimized for its specific sub-task,
further decomposed if necessary, and even easily replaced with more
effective prompts, trained models, or symbolic functions if desired.

We show that the flexibility and modularity of Decomposed Prompting
allows it to outperform prior work on few-shot prompting using GPT-3. On
symbolic reasoning tasks, we can further decompose sub-tasks that are
hard for LLMs into even simpler solvable sub-tasks. When the complexity
comes from the input length, we can recursively decompose the task into
the same task but with smaller inputs. We also evaluate our approach on
textual multi-step reasoning tasks: on long-context multi-hop QA, we can
more effectively teach the sub-tasks via our separate sub-tasks prompts;
and on open-domain multi-hop QA, we can easily incorporate a symbolic
information retrieval module within our decomposition framework, leading
to improved performance on both tasks. \(^{1}\)

\section{1 INTRODUCTION}\label{introduction}

Large Language Models (LLMs) such as GPT-3 (Brown et al., 2020) have
been shown to solve various tasks given only a few examples as prompts,
also referred to as in-context learning. These models can even perform
more complex reasoning tasks when shown the sequence of simple reasoning
steps needed to perform the complex task as a prompt (Wei et al., 2022;
Nye et al., 2021). In essence, the sequence of reasoning steps, such as
in Chains-of-Thought (CoT) prompting (Wei et al., 2022), demonstrates
how to decompose the complex task as well as how each reasoning step
should be performed. However, as tasks become more complex, few
demonstrations of the complex task aren't sufficient for current models
to learn to perform all necessary reasoning steps. E.g., few-shot
demonstrations of concatenating the \(k^{\text{th}}\) letter of words in
a string is insufficient for GPT-3 to learn to extract the
\(k^{\text{th}}\) letter, or learn to answer hard single-hop questions
when only provided a few demonstrations of multi-hop questions.
Additionally, it is unclear whether tasks such as document retrieval and
integration, for knowledge-intensive tasks, can even be done by few-shot
prompts.

To address these limitations, we propose Decomposed Prompting (DECOMP),
a new approach to solve complex tasks by instead decomposing them into
simpler sub-tasks and delegating these to sub-task specific LLMs, with
both the decomposer and the sub-task LLMs (henceforth, sub-task
handlers) having their own few-shot prompts. Fig 1 illustrates our
approach. The decomposer

\textit{[Image omitted: images/f59ab9c0e75ba4de7adc97a458a797da26b4b372208f5c784cd5dcc2e9b7e5f2.jpg]}\\
Standard Prompting

\textit{[Image omitted: images/af9aaffa56475bbb1682cb38dd416c96590f4e4f195488469760710363776a9c.jpg]}\\
Chain-of-Thought Prompting

\textit{[Image omitted: images/6bad0e3c9bbfd0ae3030c0806dc2fb5d14a7c4d9aa5867dc72235b76e28e1db8.jpg]}\\
Decomposer Prompt

\textit{[Image omitted: images/c4416a59bab388c6522161ebe12bd5e0b9813349cc90812d7e411a6261857942.jpg]}

\textit{[Image omitted: images/2295b63d0c7f91b81c0a9d25fc5158c4712802d43a0a3c50ea2bae7c16319545.jpg]}

\textit{[Image omitted: images/b72ee535fb5254d1d8280985b0004ef51a6ac0617cfc3596a4f1f57bcadb110d.jpg]}\\
Decomposed Prompting\\
Figure 1: While standard approaches only provide labeled examples (shown
as a grey input box with green label box), Chain-of-Thought prompting
also describes the reasoning steps to arrive at the answer for every
example in the prompt. Decomposed Prompting, on the other hand, uses the
decomposer prompt to only describe the procedure to solve the complex
tasks using certain subtasks. Each sub-task, indicated here with A, B
and C is handled by sub-task specific handlers which can vary from a
standard prompt (sub-task A), a further decomposed prompt (sub-task B)
or a symbolic function such as retrieval (sub-task C)

prompt only describes a sequence of sub-tasks (A, B, and C) needed to
solve the complex tasks, indicated with the dashed lines. Each sub-task
is then delegated to the corresponding sub-task handler shown on the
right.

Using a software engineering analogy, the decomposer defines the
top-level program for the complex task using interfaces to simpler,
sub-task functions. The sub-task handlers serve as modular, debuggable,
and upgradable implementations of these simpler functions, akin to a
software library. If a particular sub-task handler, say the one for
identifying the \(k^{th}\) letter or retrieving a document, is not
performing well enough, we can debug this handler in isolation, explore
alternative prompts or implementations, and seamlessly plug the improved
module back into the overall system, as a systematic way to try to
improve performance on the complex end-task.

This approach has several advantages over prior work (as also shown in
the figure). The sub-task handlers can be shown a broader and richer set
of examples (of the simpler task) than the specific ones needed for the
complex task prompt (task A). If a sub-task is too complex, it can be
further decomposed into simpler sub-tasks (task B). Similar to software
libraries, these sub-task handlers can be shared across multiple tasks;
e.g., here tasks A and C are reused in the model for task B. As noted
above, a sub-task handler can be easily swapped with an improved
implementation without any change to the rest of the system. Few-shot
prompt based LLMs can be even replaced with a symbolic system for tasks
more suited for non-neural methods; e.g., task C uses a symbolic
retrieval system such as Elasticsearch that can handle very large-scale
corpora. Lastly, we can even improve upon prior work by simply adding an
error-correcting sub-task handler as a post-processing step.

To illustrate these advantages of DECOMP, we empirically evaluate it
against prior work on eight challenging datasets using GPT3 models: (1)
On a task of concatenating the \(k^{th}\) letter, we show that our
approach of factoring out each sub-task allows us to more effectively
teach the sub-problem of extracting the \(k^{th}\) letter(specifically,
by decomposing it into even easier sub-tasks). (2) On a task of
reversing a list, we show that DECOMP allows us to extend the
capabilities of a weaker model and build a scale-invariant system by
recursively decomposing the task into reversal of smaller and smaller
lists. (3) On a task of long-context QA (Khot et al., 2022), our
approach allows each sub-task handler to accommodate more examples than
feasible with CoT prompting leading to better QA performance. (4) On
three multi-hop open-domain QA datasets (Yang et al., 2018; Ho et al.,
2020; Trivedi et al., 2022), we can incorporate a symbolic retrieval
(ElasticSearch) API as the handler for the retrieval sub-task leading to
better results than CoT. (5) On two Math QA datasets (Cobbe et al.,
2021; Roy \& Roth, 2015), we can post-process CoT to easily fix frequent
formatting errors, resulting in a surprisingly high improvement of 14-17
pts.

\section{2 RELATED WORK}\label{related-work}

Few-shot Prompts for Multi-Step Reasoning Large-scale Language models
(LLMs) have been shown to learn various NLP tasks given just few
examples as prompts (Brown et al., 2020). Recently, they have also been
successfully applied to various multi-step reasoning tasks by providing
the intermediate reasoning steps, i.e.~Chain-of-Thought (Wei et al.,
2022; Chowdhery et al., 2022), needed to arrive at the answer. An
alternate approach has been to compose multiple LLMs or LLMs with
symbolic functions to perform multi-step reasoning (Jung et al., 2022;
Creswell et al., 2023;

Press et al., 2022; Parisi et al., 2022; Gao et al., 2022; Schick et
al., 2023, inter alia). We view these prior works as specialized systems
with a pre-defined decomposition structure.

The closest works to our approach are the ideas of least-to-most
prompting (Zhou et al., 2023) and successive prompting (Dua et al.,
2022) where one prompt/model is used to generate the sub-questions
needed to answer a complex question and a second prompt/model
sequentially answers these sub-questions. In contrast, our approach
allows for diverse decomposition structures including recursion and
other non-linear decomposition structures. E.g., by definition,
least-to-most asks questions from easiest to the hardest and requires an
LLM to eventually answer the complete question (``most'' in
least-to-most) whereas we have no such restriction. Additionally, we
iteratively generate new questions based on previous answers (similar to
successive prompting) and can explicitly assign different prompts or
symbolic systems to answer each sub-question.

Modular Approaches for Multi-Step Reasoning Our work follows a long
literature in NLP on neural modular modeling architectures (Andreas et
al., 2016; Talmor \& Berant, 2018; Min et al., 2019; Jiang \& Bansal,
2019; Gupta et al., 2020; Perez et al., 2020; Khot et al., 2021; Levine
et al., 2022) for question-answering and other tasks. We take particular
inspiration from the Text Modular Networks approach of Khot et
al.~(2021), whereby problem decomposition consists of a learned next
question generator trained to generate questions in the language of a
collection of textual and symbolic agents. Best-first search strategy
was used to explore the space of possible decompositions during
inference. In contrast to this work, which largely centered around
supervised training of the next-question generator given existing
agents, we leverage the power and recent successes of few-shot LLMs to
build both the decomposer and the sub-task agents that best fit the
ideal decomposition. This has the advantage of obviating the need for
specialized supervised training data that may not always be available
for all sub-tasks -- a key bottleneck of this prior work.

\section{3 DECOMPOSED PROMPTING}\label{decomposed-prompting}

As with conventional few-shot prompting, the goal is to teach an LLM to
find an answer A to a query Q using a small set of in-context examples
\(D = \{E_{1}, \ldots, E_{|D|}\}\) . The answer A is obtained from the
underlying distribution \(p(A \mid Q, D, \theta)\) (Dohan et al., 2022).
In the most basic few-shot setup, examples take the form
\(E_{j} = (Q_{j}, A_{j})\) . In the case of CoT-style prompting, the
goal is to obtain answers by first generating a sequence or chain of
intermediate reasoning steps or ``thoughts'' T, and then deriving the
final answer based on T. To teach this ability, one uses more
sophisticated in-context examples that take the form
\(E_{j} = (Q_{j}, (T_{j,1}, \ldots, T_{j,k}), A_{j})\) .

In DECOMP, the core is a decomposer LLM that tries to solve a complex
task by generating a prompting program P for it. Each step of P directs
a simpler sub-query to a function in an auxiliary set of sub-task
functions F available to the system. Given a query Q whose answer is A,
the program P is a sequence of the form
\(\left((f_{1}, Q_{1}, A_{1}), \ldots, (f_{k}, Q_{k}, A_{k})\right)\)
where \(A_{k}\) is the final answer predicted by P and \(Q_{i}\) is a
sub-query directed to the sub-task function \(f_{i} \in F\) . P is
executed by a high-level imperative controller, which passes the inputs
and outputs between the decomposer and sub-task handler until a stopping
condition in P is met and the final output obtained.

To teach the decomposer LLM in a few-shot prompting manner, we use
in-context examples that take the form
\(E_{j} = \left(\left(Q_{j},\left(f_{j,1},Q_{j,1},A_{j,1}\right),\dots,\left(f_{j,k_{j}},Q_{j,k_{j}},A_{j,k_{j}}\right)\right)\right)\)
where \(A_{j,k_j} = A_j\) is the final answer for \(Q_{j}\) and
\((Q_{j,1},\dots ,Q_{j,k_j})\) is a decomposition of \(Q_{j}\) . Each
sub-task function \(f\) , in turn, is operationalized via a sub-task
handler as an in-context prompting LLM (e.g., a separate CoT-style
prompt or a additional prompting program dedicated to that sub-task), or
any other symbolic or learned function (e.g., a calculator or
specialized supervised trained model).

\section{3.1 DECOMPOSED PROMPTS}\label{decomposed-prompts}

To illustrate this with an example, consider a multi-step task such as
``Concatenate the first letter of every word in str using a space''. We
can solve this task by decomposing it into a sequence of three simple
sub-tasks: 1) Collect the list of words in the str; 2) For each word,
extract the third letter; 3) Concatenate the extracted letters using
space as the separator. Fig. 2 shows an example decomposition prompt for
this task. Much like a conventional structured program, the top-level
decomp prompt provides an example program \(E_{j}\) using three sub-task
functions: \(f_{1}\) :split that

\textit{[Image omitted: images/4bd726eb7e9041e87573cd38463c1d84da5028c50ef3dcf8c3ea9617e5743ee3.jpg]}

text\_image

QC: Concatenate the first letter of every word in ``Jack Ryan'' using
spaces Q1: \hyperref[split]{split} What are the words in ``Jack Ryan''?
\#1: {[}``Jack''. ``Ryan''{]} Q2: (foreach) {[}str\_pos{]} What is the
first letter of \#1? \#2: {[}``J'', ``R''{]} Q3: \hyperref[merge]{merge}
Concatenate \#2 with spaces \#3: ``J R'' Q4: {[}EOQ{]} \ldots{}

\textit{[Image omitted: images/116a9bed630e749a4bfb8c19dcd0c45026ed906b833fae10545d0e3e94510e47.jpg]}

text\_image

Q: What are the words in ``Elon Musk Tesla''? A: {[}``Elon'', ``Musk'',
``Tesla''{]} Q: What are the letters in ``C++''? A: {[}``C'', ``+'',
``+''{]} \ldots{} split

Q: Concatenate {[}``n'', ``i'', ``e''{]} A: ``nie'' Q: Concatenate
{[}``n'', ``i'', ``c'', ``e''{]} using spaces A: ``n i c e'' \ldots{}
merge

Figure 2: Prompts for the decomposer and the split and merge sub-tasks
used by the decomposer. The decomposer specifies the sequence of
questions and corresponding sub-tasks (within square braces). The
sub-task prompts can be written independent of the complex task examples
and can even capture generalizations, e.g., letters in word (split) and
no delimiter (merge).

splits words in an input string, \(f_{2}\) :str\_pos that finds
character positions in strings and \(f_{3}\) :merge that concatenates
characters. In this case, we operationalize each sub-task function as a
separate in-context prompt (e.g., using a standard prompting approach
for split and merge on the right side), each containing a set of
in-context examples that are independent of the original complex task.

In addition to the three functions described above, additional control
structure is included, such as the symbolic function foreach, which
iterates over arrays and references to previous answers such as \#1. We
note that such a helper function is not strictly necessary (e.g., we
could directly generate ``Q2'': What is the first letter of Jack?'' and
``Q3'': What is the first letter of Ryan?'' instead of Q2 in the figure)
and is added to reduce the manual effort needed to specify the
decomposition and also reduce potential errors during decomposition. In
our experiments we use two of the compositional operators defined by
Khot et al.~(2022) (see appendix for details), although it is capable of
using all their operators (which also capture the QDMR operators from
Wolfson et al.~(2020)).

\textit{[Image omitted: images/c34a9140c553612ff75691ece7a456b9d6a7566baa9186fe6af685de23e471f8.jpg]}

flowchart

\begin{Shaded}
\begin{Highlighting}[]
\NormalTok{graph TD}
\NormalTok{    A["QC: Concatenate the second letter of every word in \&quot;John Smith\&quot; using spaces"] {-}{-}\textgreater{} B["Q1: [split"] What are the words in "John Smith"? \#1: ["John", "Smith"]]}
\NormalTok{    A {-}{-}\textgreater{} C["Q2: [foreach"] [str\_pos] What is the second letter in \#1? \#2: ["o", "m"]]}
\NormalTok{    A {-}{-}\textgreater{} D["Q3: [merge"] Concatenate \#2 with spaces \#3: "o m"]}
\NormalTok{    A {-}{-}\textgreater{} E["Q4: [EOQ"] A: "o m"]}
\NormalTok{    B {-}{-}\textgreater{} F["Q: What are the words in \&quot;John Smith\&quot;? A: [\&quot;John\&quot;, \&quot;Smith\&quot;"]]}
\NormalTok{    C {-}{-}\textgreater{} G["Q: What is the second letter in \&quot;John\&quot;? A: \&quot;o\&quot;"]}
\NormalTok{    C {-}{-}\textgreater{} H["Q: What is the second letter in \&quot;Smith\&quot;? A: \&quot;m\&quot;"]}
\NormalTok{    D {-}{-}\textgreater{} I["Q: Concatenate [\&quot;o\&quot;, \&quot;m\&quot;"] with spaces A: "o m"]}
\NormalTok{    E {-}{-}\textgreater{} J["merge"]}
\NormalTok{    F {-}{-}\textgreater{} K["split"]}
\NormalTok{    G {-}{-}\textgreater{} L["str\_pos"]}
\NormalTok{    H {-}{-}\textgreater{} M["merge"]}
\NormalTok{    I {-}{-}\textgreater{} N["merge"]}
\end{Highlighting}
\end{Shaded}

Figure 3: The inference procedure in DECOMP iteratively calls the
decomposer prompt to generate the next question and sub-task at each
step, given the current history of question and answers. The generated
question is then routed to the assigned sub-task handler (with some
handling of special operators, when needed). When the special
end-of-questions {[}EOQ{]} marker is generated, the previous answer is
returned as the final prediction.

\section{3.2 PROMPT EXECUTION AND
INFERENCE}\label{prompt-execution-and-inference}

Given a new question and a set of background in-context examples D, the
inference (i.e., the program construction and execution) process is
illustrated in Fig. 3. The new complex question is fed to the decomposer
prompt to get the first sub-question to be asked to the split prompt.
With the help of our symbolic controller, the answer generated from this
prompt is then appended to the decomposer prompt to get the second
sub-question, Q2. Due to the foreach operator in the generated question,
Q2 results in two questions (one for each word in \#1) to be fed to the
str\_pos prompt. The answers are combined into an array to get the
answer \#2. The entire decomposition history is used to generate Q3 and
passed to the merge prompt to get the final answer. Since the task has
been solved, the decomposition prompt produces the special
end-of-sequence marker({[}EOQ{]}) and the last answer is returned as the
final answer. Formally, performing inference involves finding the best
answer A to a new query Q, which in the simplest form involves computing
the MAP answer using the

LLMs predictive distribution for \(A\) , i.e.,
\(\hat{A} = \arg \max_A p(A \mid D, Q, \theta)\) (Dohan et al., 2022).
For practicality, such computations are approximated using greedy search
in our experiments.

\section{3.3 DECOMP CAPABILITIES}\label{decomp-capabilities}

Hierarchical Decomposition Certain sub-tasks, even when given many
examples, are not solvable with few-shot prompting. E.g., we found
identifying the \(k^{th}\) letter of a string to be challenging for the
GPT3 text-davinci-002 model. In such a scenario, we can decompose the
sub-task prompt further, to first identify the letters and their
position and then select the \(k^{th}\) element of this array (see Fig.
4). We can also re-use existing sub-task prompts in our framework. E.g.,
the split prompt can be reused since it was developed for the general
task of splitting strings. \(^{2}\)

\textit{[Image omitted: images/3a3af31a5a009623fae81500bed2706e48e91332b0ac0fe3b8607f5b9e1e844f.jpg]}

text\_image

QC: What is the letter at position 1 of the word ``Google''? QS:
\hyperref[split]{split} What are the letters and their positions in
``Google''? A: ``{[}(G, 1), (o, 2), (o, 3), (g, 4), (l, 5), (e, 6){]}''
QS: {[}arr\_pos{]} What is at position 1 in \#1? A: ``G'' QS: {[}EOQ{]}
\ldots.

Q: What is at position 1 in "{[}("Colorless", 1), ("green", 2),
("ideas", 3), ("sleep", 4), ("furiously", 5){]}"? A: "Colorless"

Q: What is at position 2 in "{[}(G, 1), (o, 2), (o, 3), (g, 4), (l, 5),
(e, 6){]}"? A: "o" \ldots. arr\_pos

Figure 4: Since identifying the \(k^{th}\) character is challenging for
GPT3 davinci-002 model, we further decompose it into two simpler
sub-tasks: split the word into its letters (using the shared sub-task
split) and then return the \(k^{th}\) item of this list using the
arr\_pos prompt.

Recursive Decomposition Some problems can be naturally broken down into
one or more smaller problems of the same form. Recursive algorithms such
as merge sort use this idea to solve large problems efficiently, using a
succinctly described method. We apply this same principle in DECOMP by
allowing the decomposer prompt to recursively call itself, as shown in
Fig. 5 for the task of list reversal. By using recursion, we are able to
generalize any base prompting approach (CoT in this figure) to much
longer lists by breaking the input into smaller and smaller lists till
we reach a list length where the model is highly accurate. Such
recursive approaches can not be described by current methods such as CoT
and standard prompting. Least-to-most prompting (Zhou et al., 2023) also
proposes a similar solution but differs in two key aspects (a) it has to
identify all the subproblems in one-shot instead of our iterative
top-down decomposition (b) it has to learn to identify the relevant
answers from the previous solutions which we get for free from our
decomposition.

\textit{[Image omitted: images/b782931e1b5882a388129d5415b6911e65c4159e08937637a086bf390b549a30.jpg]}\\
Figure 5: Sample prompt for recursive decomposition for reversing lists.
Each list is split into two halves and each half is reversed and
concatenated in the reverse order. We can recursively split a list till
we hit the base case (lists of length 3 here) where existing approaches
such as CoT are accurate.

External API Calls In certain cases, the sub-tasks may not be feasible
to solve using only a LLM. E.g., retrieving knowledge from a KB or large
corpus. Such sub-tasks, however, can be easily solved using existing
systems such as retrieving documents using an Elasticsearch index or
webpages using Google search (Lazaridou et al., 2022). Fig. 6 shows how
DECOMP can easily use such a system to retrieve the relevant documents
and answer a single-hop open-domain question.

\textit{[Image omitted: images/8f38237251d8319827af746b857e0fe932477877955769f1ed0cee7b8d92aab0.jpg]}\\
Figure 6: A Decomposed Prompt to answer open-domain questions using
Elasticsearch-based retrieval. Full usage of this prompt for open-domain
multihop questions is given in Fig. 11.

\section{4 CASE STUDIES}\label{case-studies}

We showcase DECOMP's strengths through four tasks; two symbolic
manipulation tasks similar to those investigated by Wei et al.~(2022)
and two existing textual multi-hop reasoning tasks. Unless specified, we
use text-davinci-002 InstructGPT3 model (Ouyang et al., 2022) as the LLM
and report the Exact Match (EM) numbers, following prior work. For
order-independent list answers, we evaluate set equality as EM. We
compare our approach to CoT rather than each specific decomposition
structure used in prior work. See App. G for the complete prompts for
all our tasks.

\section{\texorpdfstring{4.1 \(k^{th}\) LETTER CONCATENATION
(HIERARCHICAL
DECOMPOSITION)}{4.1 k\^{}\{th\} LETTER CONCATENATION (HIERARCHICAL DECOMPOSITION)}}\label{kth-letter-concatenation-hierarchical-decomposition}

We compare DECOMP to CoT prompting for concatenating letters at the
\(k^{th}\) position. All prompts contain examples of concatenating
letters in position 1, 4, and last position of strings with 3 words. We
create three different prompts for all our baselines and present the
average to account for variance due to the choice of examples following
Perez et al.~(2021). We use the decomp, split, str\_pos (further
decomposed as shown in Fig. 4), and merge prompts for decomposition
prompting. We adapt the CoT for last letter concatenation from prior
work (Wei et al., 2022) for this task as shown below. In addition, we
consider a rolled out version of our decomposition prompts in terms of a
CoT, i.e., we describe the entire decomposition process (identify words,
split each word into letters, take \(k^{th}\) letter and concatenate) as
a single CoT. e.g., for the question ``Take the letters at position 4 of
the words in ``Herbert Alexander Simon'' and concatenate them using a
space.'', we use the CoT:

\textit{[Image omitted: images/ba9fd76f8c734cf1e78a5d3c92e7beec9c83b0a154240f88acd4a431ce4f3001.jpg]}

We similarly adapt the least-to-most prompt (Zhou et al., 2023) to
include rollout. (see App. G). We compare these four prompting
techniques on 4 datasets to evaluate generalization along 3 axes: (1)
new letter position \(k = 3\) ; \(^{3}\) (2) longer inputs, \#words=4
and 5; (3) new delimiter ``;''. The words in the test examples come from
a list of most popular first and last names. \(^{4}\) All evaluation
datasets have 100 examples. We present results on space as a delimiter
averaged across three prompts in Fig. 7. \(^{5}\)

DECOMP outperforms chain-of-thought and least-to-most prompting, even
when the prompt uses the same reasoning procedure as the rolled out
decomposition. This shows that the separate prompts are more effective
at teaching hard sub-tasks than a single CoT prompt.

DECOMP generalizes perfectly to longer sequences. As the length of the
input sequence increases, our approach continues to achieve close to
100\% accuracy on this task. \(^{6}\) The CoT-based approaches drop
noticeably in their scores with longer input lengths, widening the
performance gap.

\textit{[Image omitted: images/1d9f01b79c6e77e43470c3d652cc6324fb9c054a4779c49cc594ef7279ee18a8.jpg]}

bar

{\def\LTcaptype{none} % do not increment counter
\begin{longtable}[]{@{}lllll@{}}
\toprule\noalign{}
Category & CoT & CoT w/ rollout & L2M w/ rollout & DecomP \\
\midrule\noalign{}
\endhead
\bottomrule\noalign{}
\endlastfoot
N=3(IID) & 23.7 & 71.3 & 74.7 & 98.0 \\
N=4(OOD) & 12.0 & 53.7 & 70.5 & 96.0 \\
N=5(OOD) & 6.0 & 53.3 & 66.0 & 97.0 \\
\end{longtable}
}

Figure 7: EM Results on the \(k^{th}\) letter concatenation task (k=3)
using space as delimiter with different number of words in the input.
DECOMP outperforms and generalizes better than CoT as well as
Least-to-most prompting.

\textit{[Image omitted: images/123046f5348edf40fef2a923c0cc84b95a2fc274c720f17f3a457a3ca373a150.jpg]}

bar

{\def\LTcaptype{none} % do not increment counter
\begin{longtable}[]{@{}llll@{}}
\toprule\noalign{}
Category & CoT & CoT w/ rollout & DecomP \\
\midrule\noalign{}
\endhead
\bottomrule\noalign{}
\endlastfoot
N=4(IID) & 5.5 & 76 & 86 \\
N=6(OOD) & 3.0 & 53 & 59 \\
N=8(OOD) & 6.5 & 16 & 64 \\
N=10(OOD) & 4.5 & 1 & 42 \\
\end{longtable}
}

Figure 8: EM results on reversing sequences. Incorporating CoT in DECOMP
greatly increases the ability of the model to generalize to new sequence
lengths.

\section{4.2 LIST REVERSAL (RECURSIVE
DECOMPOSITION)}\label{list-reversal-recursive-decomposition}

We use the task of reversing lists of words \(^{7}\) to show how
recursive DECOMP enables length generalization. We adapt the relevant
CoT prompt from Wei et al.~(2022), and integrate it in a decomposed
prompt. As a control, we also compare to a CoT version w/ rollout of our
decomposed prompt. All prompts contain the same 3 examples of reversing
word sequences with 3-5 items. We evaluate all prompts for
generalization to 4, 6, 8, and 10-item sequences. Here we use
davinci-001 to show that DECOMP enables a weaker model approach
davinci-002's performance (which does solve this task). We use the
strategy from Fig. 5 and provide our prompts in App. G. Fig. 8 shows the
results of the prompting strategies on different input lengths.

DECOMP improves the length generalization of few-shot prompting. While
our base CoT prompt does not generalize at all to longer sequences, our
approach can recursively decompose the problem and achieve better length
generalization. Moreover, the CoT version of our decomposition strategy
fails because the unrolled prompt becomes too long and convoluted
without the ability to abstract away sub-modules.

\section{4.3 LONG-CONTEXT QUESTION
ANSWERING}\label{long-context-question-answering}

We next evaluate on the CommaQA-E dataset (Khot et al., 2022) under the
reading comprehension setting. The dataset consists of synthetically
generated entities (e.g.~Erowid award), facts (``Wetherality was an
actor in the movie Dewbar.'') and multi-hop questions (e.g., ``What
awards have the actors of the Erowid winning movies received?''). Due to
the presence of many distractors and, as a result, longer context, this
dataset has been shown to be hard for standard LMs even when fine-tuned.

What awards have movies produced by people born in 1910 won?QS:
\hyperref[qa]{qa} Who were born in the year 1910?A: {[}"Teeplemole",
"Muntaril"{]}QS: (foreach\_merge) \hyperref[qa]{qa} For which movies was
\#1 the producer?A: {[}"Featsaw", "Zalate", "Premercy"{]}QS:
(foreach\_merge) \hyperref[qa]{qa} Which awards were given to \#2?A:
{[}"Zorgion", "Chowwurst", "Hallowcock"{]}QS: {[}EOQ{]}\ldots{}

What awards have movies produced by people born in 1910 won?QS:
{[}simp\_qa{]} Who were born in the year 1910?A: {[}"Teeplemole",
"Muntaril"{]}QS: (foreach\_merge) {[}pos\_qa{]} For which movies was \#1
the producer?A: {[}"Featsaw", "Zalate", "Premercy"{]}QS:
(foreach\_merge) {[}aw\_qa{]} Which awards were given to \#2?A:
{[}"Zorgion", "Chowwurst", "Hallowcock"{]}QS: {[}EOQ{]}\ldots{}

fine

Figure 9: Sample prompts used for the CommaQA dataset. On the left, the
coarse-grained decomposition defines a single QA sub-task with all
single-hop questions being delegated to a single sub-task handler. On
the right, the fine-grained decomposition assigns questions to three
different sub-tasks (see App. G for their prompts) depending on the
question type. This allows us to provide more examples for each question
type allowing the model to learn the sub-task more effectively.

To fit these questions within GPT3's context limit (2049 tokens), we
generate a smaller version of the CommaQA-E dataset and of the
compositional generalization split such that we can fit at least

four examples in the context for CoT prompts. The CoT prompts describe
the sequence of facts needed to arrive at the answer (see App. G for all
the prompts).

For DECOMP, we can separate the task of decomposition (independent of
the context) from the sub-tasks of single-hop question answering. As
shown in Fig. 9, we provide examples of the context-independent
decomposition in the decomposer prompt and use the separate sub-task
prompts to teach the QA skill over the given context. Additionally, we
can choose the granularity of decomposition to trade off human effort
for increased accuracy. For example, we could have single QA prompt to
handle all the questions or create QA prompts for different classes of
questions. In our experiments, each sub-task prompt contains 8 QA
examples (2 questions/para). We evaluate three different prompts and
report the average results in Fig. 10.

We make three observations on CommaQA. DECOMP is more accurate than CoT
irrespective of the granularity of decomposition or the evaluation
split. Finer grained decomposition can help im-

prove task performance by providing more examples for each class of
questions, which in turn increases single-hop QA accuracy. DECOMP
generalizes to new compositions such as the compositional generalization
split of CommaQA, which tests models on unseen compositions of relations
observed in the training set. While CoT has a drop in score, both
decomposition-based approaches actually get a small bump (the subset of
relations used in this split are easier for our QA models).

\textit{[Image omitted: images/78d77b7dae3eb4693ac8b5f774816e45009c838720c51e2a134702b0637cc0e1.jpg]}

bar

GPT3 (text-davinci-002) \textbar{} Dataset \textbar{} CoT \textbar{}
DecomP(coarse) \textbar{} DecomP(fine) \textbar{} \textbar{} :---
\textbar{} :--- \textbar{} :--- \textbar{} :--- \textbar{} \textbar{}
IID \textbar{} 42.0 \textbar{} 55.0 \textbar{} 59.7 \textbar{}
\textbar{} Comp. Gen \textbar{} 33.8 \textbar{} 59.7 \textbar{} 64.2
\textbar{}

Figure 10: EM results on the CommaQA-E datasets. DECOMP always
outperforms CoT, with fine-grained marginally out-performing
coarse-grained decomposition.

\section{4.4 OPEN-DOMAIN QUESTION
ANSWERING}\label{open-domain-question-answering}

Next, we demonstrate the ability of our approach to integrate external
API calls on the task of open-domain multihop question answering. We
evaluate our approach on three datasets: (1) 2Wiki-MultihopQA (Ho et
al., 2020) (2) MuSiQue (Trivedi et al., 2022) (3) HotpotQA (Yang et al.,
2018). We describe the open-domain versions of these datasets in more
detail in App. A We use the Codex (code-davinci-002) model here since it
can fit the much longer contexts needed. We also evaluate the impact of
model scale on DECOMP by using models from the Flan-T5 family:
Flan-T5-Large (0.7B), Flan-T5-XL (3B), and Flan-T5-XXL (11B). \(^{8}\)

\textit{[Image omitted: images/6b66619d3530a2daea60f1a53eb69c0ae265d6b6782925d37e3725a4709cd1ce.jpg]}\\
Figure 11: The prompt used to answer open-domain multihop questions
using Elasticsearch-based retrieval. The retrieve\_odqa prompt is given
in Fig. 6.

Fig. 11 shows the decomposition prompt we use. The decomposer generates
(singlehop) sub-questions and delegates them to retrieve\_odqa
(described in Fig. 6). As we showed earlier, this module retrieves
relevant documents then uses an RC model to answer. retrieve\_odqa
returns both the answer and the documents, allowing subsequent
sub-questions to use the answers (e.g.~``Mack Rides'') and the
multihop\_rcqa model to use the documents. The final multihop\_rcqa
model is prompted to produce the answer directly or using CoT given K
paragraphs.

We compare our approach against two baselines: A. No Context (No-Ctxt),
A closed-book setting baseline where the model must rely only on its
parametric knowledge. B. NoDecomp Context (NoDecomp-Ctxt), A simple
retrieval baseline where we retrieve K paragraphs using the multi-hop
question as the input and use that as context. For both NoDecomp-Ctxt
and Decomp-Ctxt, K is selected by hyperparameter tuning (App. A). We
manually annotate CoTs and decompositions for 20 training set questions,
and sample 3 prompts of 15 questions each for all approaches. The
detailed prompts are given in the Appendix G. We evaluate on 300
held-out dev questions in each dataset.

\textit{[Image omitted: images/9cdebacdbe729fd2b9a0676d758bbd98364cd6ceaf228e0b9ea59183be87ffda.jpg]}

bar

{\def\LTcaptype{none} % do not increment counter
\begin{longtable}[]{@{}llll@{}}
\toprule\noalign{}
Dataset & No-Ctxt QA & NoDecomp-Ctxt QA & Decomp-Ctxt QA \\
\midrule\noalign{}
\endhead
\bottomrule\noalign{}
\endlastfoot
MuSiQue & 18.3 & 21.4 & 25.4 \\
HotpotQA & 41.2 & 50.2 & 49.9 \\
2WikiMultihopQA & 38.1 & 47.1 & 64.1 \\
\end{longtable}
}

code-davinci-002

\textit{[Image omitted: images/2f6fd8f08fda2191ca9aab0ae5ec71a104ca5366c881a0d8c5462f989723d58d.jpg]}

bar

Flan-T5-XXL \textbar{} Dataset \textbar{} No-Ctxt QA \textbar{}
NoDecomp-Ctxt QA \textbar{} Decomp-Ctxt QA \textbar{} \textbar{} :---
\textbar{} :--- \textbar{} :--- \textbar{} :--- \textbar{} \textbar{}
MuSiQue \textbar{} 15.1 \textbar{} 23.3 \textbar{} 29.3 \textbar{}
\textbar{} HotpotQA \textbar{} 24.8 \textbar{} 49.0 \textbar{} 56.4
\textbar{} \textbar{} 2WikiMultihopQA \textbar{} 32.3 \textbar{} 49.9
\textbar{} 63.3 \textbar{}

Figure 12: Answer F1 \(^{9}\) on three open-domain QA datasets using two
base LMs: Codex (left) and Flan-T5-XXL (right) with direct prompting.
Decomp-Ctxt models (ours) significantly outperforms the No-Ctxt models
(no retrieval) in all settings and also outperforms our strong retrieval
baseline (NoDecomp-Ctxt QA), with the exception of Codex on HotpotQA
where it is comparable. See App. A.3 for results on smaller Flan-T5
models and CoT prompting.

We present results on all three datasets with direct QA prompts in Fig.
12 with other results in App. A. The Decomp-Ctxt models perform
significantly better than No-Ctxt models in all the settings showing
that external knowledge can be leveraged to improve few-shot models on
open-domain mulithop QA. Furthermore, we show that our Decomp-Ctxt
models outperform the strong retrieval baseline (NoDecomp-Ctxt) in all
settings except one (Codex with HotpotQA). Finally, we show that even
with the much smaller Flan-T5-XXL model, Decomp-Ctxt outperforms all the
baselines and can even achieve scores comparable to the Codex-only
systems.

\section{4.5 ADDITIONAL RESULTS}\label{additional-results}

Post-processing CoT for error correction DECOMP also allows us to create
a targeted sub-task handler to focus on the source of error in any
system. For example, CoT for arithmetic reasoning often rely on patterns
(answer is .*) to extract answers but the CoT does not always fit this
pattern. Instead, we can assign the answer extraction to a better
sub-task handler (GPT3) and reduce these types of errors. This results
in a 17 pt improvement on MultiArith (78 → 95) and 14 pt improvement on
GSM8K (36 → 50.6) compared to CoT prompting (details in App. B).

While DECOMP outperforms the baselines in aggregate, we also see the
gains of DECOMP are consistent across prompt choices (see App. D) and
decomposition schemes (see App. E).

\section{5 CONCLUSION}\label{conclusion}

We proposed a new approach, Decomposed Prompting, to solve complex tasks
using few-shot prompts, by decomposing them into a prompting program
built out of simpler sub-tasks. Drawing inspiration from software
libraries, our decomposer and shared sub-tasks are designed in a modular
fashion: they use their own few-shot prompts, allowing one to
independently optimize each prompt, decompose a sub-task further if
necessary, or even seamlessly replace it with a symbolic system. We show
that Decomposed Prompting outperforms prior work on four different tasks
and generalization settings, establishing it as an effective few-shot
paradigm for solving complex tasks.

\section{ACKNOWLEDGEMENTS}\label{acknowledgements}

We thank members of the Aristo team at the Allen Institute for AI (AI2)
for their constructive feedback and the reviewers for their invaluable
suggestions. This work was supported in part by the National Science
Foundation under grants IIS2007290.

\section{REFERENCES}\label{references}

Jacob Andreas, Marcus Rohrbach, Trevor Darrell, and Dan Klein. Neural
module networks. In CVPR, 2016.\\
Tom Brown, Benjamin Mann, Nick Ryder, Melanie Subbiah, Jared D Kaplan,
Prafulla Dhariwal, Arvind Neelakantan, Pranav Shyam, Girish Sastry,
Amanda Askell, Sandhini Agarwal, Ariel Herbert-Voss, Gretchen Krueger,
Tom Henighan, Rewon Child, Aditya Ramesh, Daniel Ziegler, Jeffrey Wu,
Clemens Winter, Chris Hesse, Mark Chen, Eric Sigler, Mateusz Litwin,
Scott Gray, Benjamin Chess, Jack Clark, Christopher Berner, Sam
McCandlish, Alec Radford, Ilya Sutskever, and Dario Amodei. Language
models are few-shot learners. In H. Larochelle, M. Ranzato, R. Hadsell,
M.F. Balcan, and H. Lin (eds.), Advances in Neural Information
Processing Systems, volume 33, pp.~1877--1901. Curran Associates, Inc.,
2020. URL
https://proceedings.neurips.cc/paper/2020/file/1457c0d6bfcb4967418bfb8ac142f64a-Paper.pdf.\\
Aakanksha Chowdhery, Sharan Narang, Jacob Devlin, Maarten Bosma, Gaurav
Mishra, Adam Roberts, Paul Barham, Hyung Won Chung, Charles Sutton,
Sebastian Gehrmann, et al.~PaLM: Scaling language modeling with
pathways. arXiv preprint arXiv:2204.02311, 2022.\\
Karl Cobbe, Vineet Kosaraju, Mohammad Bavarian, Mark Chen, Heewoo Jun,
Lukasz Kaiser, Matthias Plappert, Jerry Tworek, Jacob Hilton, Reiichiro
Nakano, Christopher Hesse, and John Schulman. Training verifiers to
solve math word problems. arXiv preprint arXiv:2110.14168, 2021.\\
Antonia Creswell, Murray Shanahan, and Irina Higgins.
Selection-inference: Exploiting large language models for interpretable
logical reasoning. In ICLR, 2023. URL
https://openreview.net/forum?id=3Pf3Wg6o-A4.\\
David Dohan, Winnie Xu, Aitor Lewkowycz, Jacob Austin, David Bieber,
Raphael Gontijo Lopes, Yuhuai Wu, Henryk Michalewski, Rif A Saurous,
Jascha Sohl-dickstein, et al.~Language model cascades. arXiv preprint
arXiv:2207.10342, 2022.\\
Dheeru Dua, Shivanshu Gupta, Sameer Singh, and Matt Gardner. Successive
prompting for decomposing complex questions. In Proceedings of the 2022
Conference on Empirical Methods in Natural Language Processing,
pp.~1251--1265, Abu Dhabi, United Arab Emirates, December 2022.
Association for Computational Linguistics. URL
https://aclanthology.org/2022.emnlp-main.81.\\
Luyu Gao, Aman Madaan, Shuyan Zhou, Uri Alon, Pengfei Liu, Yiming Yang,
Jamie Callan, and Graham Neubig. PAL: Program-aided language models.
ArXiv, abs/2211.10435, 2022.\\
Nitish Gupta, Kevin Lin, Dan Roth, Sameer Singh, and Matt Gardner.
Neural module networks for reasoning over text. In ICLR, 2020.\\
Xanh Ho, A. Nguyen, Saku Sugawara, and Akiko Aizawa. Constructing a
multi-hop qa dataset for comprehensive evaluation of reasoning steps. In
COLING, 2020.\\
Yichen Jiang and Mohit Bansal. Self-assembling modular networks for
interpretable multi-hop reasoning. In EMNLP, 2019.\\
Jaehun Jung, Lianhui Qin, Sean Welleck, Faeze Brahman, Chandra
Bhagavatula, Ronan Le Bras, and Yejin Choi. Maieutic prompting:
Logically consistent reasoning with recursive explanations. In EMNLP,
2022.\\
Tushar Khot, Daniel Khashabi, Kyle Richardson, Peter Clark, and Ashish
Sabharwal. Text modular networks: Learning to decompose tasks in the
language of existing models. In NAACL, 2021.

Tushar Khot, Kyle Richardson, Daniel Khashabi, and Ashish Sabharwal. Hey
AI, can you solve complex tasks by talking to agents? In Findings of
ACL, 2022.\\
Angeliki Lazaridou, Elena Gribovskaya, Wojciech Stokowiec, and Nikolai
Grigorev. Internet-augmented language models through few-shot prompting
for open-domain question answering. ArXiv, abs/2203.05115, 2022.\\
Yoav Levine, Itay Dalmedigos, Ori Ram, Yoel Zeldes, Daniel Jannai, Dor
Muhlgay, Yoni Osin, Opher Lieber, Barak Lenz, Shai Shalev-Shwartz, et
al.~Standing on the shoulders of giant frozen language models. arXiv
preprint arXiv:2204.10019, 2022.\\
Sewon Min, Victor Zhong, Luke Zettlemoyer, and Hannaneh Hajishirzi.
Multi-hop reading comprehension through question decomposition and
rescoring. In ACL, 2019.\\
Maxwell Nye, Anders Johan Andreassen, Guy Gur-Ari, Henryk Michalewski,
Jacob Austin, David Bieber, David Dohan, Aitor Lewkowycz, Maarten Bosma,
David Luan, Charles Sutton, and Augustus Odena. Show your work:
Scratchpads for intermediate computation with language models. ArXiv,
abs/2112.00114, 2021.\\
Long Ouyang, Jeff Wu, Xu Jiang, Diogo Almeida, Carroll L. Wainwright,
Pamela Mishkin, Chong Zhang, Sandhini Agarwal, Katarina Slama, Alex Ray,
John Schulman, Jacob Hilton, Fraser Kelton, Luke E. Miller, Maddie
Simens, Amanda Askell, Peter Welinder, Paul Francis Christiano, Jan
Leike, and Ryan J. Lowe. Training language models to follow instructions
with human feedback. In NeurIPS, 2022.\\
Aaron Parisi, Yao Zhao, and Noah Fiedel. TALM: Tool augmented language
models. arXiv preprint arXiv:2205.12255, 2022.\\
Ethan Perez, Patrick Lewis, Wen-tau Yih, Kyunghyun Cho, and Douwe Kiela.
Unsupervised question decomposition for question answering. In EMNLP,
2020.\\
Ethan Perez, Douwe Kiela, and Kyunghyun Cho. True few-shot learning with
language models. In NeurIPS, 2021.\\
Ofir Press, Muru Zhang, Sewon Min, Ludwig Schmidt, Noah A Smith, and
Mike Lewis. Measuring and narrowing the compositionality gap in language
models. arXiv preprint arXiv:2210.03350, 2022.\\
Pranav Rajpurkar, Jian Zhang, Konstantin Lopyrev, and Percy Liang.
SQuAD: 100, 000+ questions for machine comprehension of text. In EMNLP,
2016.\\
Subhro Roy and Dan Roth. Solving general arithmetic word problems. In
Proceedings of the 2015 Conference on Empirical Methods in Natural
Language Processing, pp.~1743--1752, 2015. doi: 10.18653/v1/D15-1202.
URL https://aclanthology.org/D15-1202.\\
Timo Schick, Jane Dwivedi-Yu, Roberto Dessì, Roberta Raileanu, Maria
Lomeli, Luke Zettlemoyer, Nicola Cancedda, and Thomas Scialom.
Toolformer: Language models can teach themselves to use tools. ArXiv,
abs/2302.04761, 2023.\\
Alon Talmor and Jonathan Berant. The web as a knowledge-base for
answering complex questions. In NAACL, 2018.\\
Harsh Trivedi, Niranjan Balasubramanian, Tushar Khot, and Ashish
Sabharwal. MuSiQue: Multi-hop questions via single-hop question
composition. TACL, 2022.\\
Xuezhi Wang, Jason Wei, Dale Schuurmans, Quoc Le, Ed Chi, and Denny
Zhou. Self-consistency improves chain of thought reasoning in language
models. In ICLR, 2023.\\
Jason Wei, Xuezhi Wang, Dale Schuurmans, Maarten Bosma, Ed Chi, Quoc Le,
and Denny Zhou. Chain of thought prompting elicits reasoning in large
language models. In NeurIPS, 2022.\\
Tomer Wolfson, Mor Geva, Ankit Gupta, Matt Gardner, Yoav Goldberg,
Daniel Deutch, and Jonathan Berant. Break it down: A question
understanding benchmark. TACL, 2020.

Zhilin Yang, Peng Qi, Saizheng Zhang, Yoshua Bengio, William W. Cohen,
Ruslan Salakhutdinov, and Christopher D. Manning. HotpotQA: A dataset
for diverse, explainable multi-hop question answering. In EMNLP, 2018.\\
Denny Zhou, Nathanael Scharli, Le Hou, Jason Wei, Nathan Scales, Xuezhi
Wang, Dale Schuurmans, Olivier Bousquet, Quoc Le, and Ed Chi.
Least-to-most prompting enables complex reasoning in large language
models. In ICLR, 2023.

\section{A OPEN DOMAIN QA DETAILS}\label{a-open-domain-qa-details}

\section{A.1 RETRIEVAL CORPUSES FOR OPEN DOMAIN
QA}\label{a.1-retrieval-corpuses-for-open-domain-qa}

We use HotpotQA in the fullwiki setting where it comes with the
associated Wikipedia corpus for open-domain QA. 2WikiMultihopQA and
MuSiQue, however, are originally reading comprehension datasets.
Questions in 2WikiMultihopQA and MuSiQue are associated with 10 and 20
paragraphs respectively. To turn these datasets into open-domain QA
datasets, we create a corpora for each dataset by combining all the
paragraphs in the train, dev and test questions. As a result we get a
corpus size of 430,225 paragraphs for 2WikiMultihopQA and 139,416 for
MuSiQue.

\section{A.2 HYPERPARAMETER TUNING FOR OPEN DOMAIN
QA}\label{a.2-hyperparameter-tuning-for-open-domain-qa}

We treat the number of paragraphs to retrieve \((K)\) in NoDecomp-Ctxt
and Decomp-Ctxt models as a hyperparameter. We select it based on a grid
search on a set of values to maximize performance on a held out set of
100 questions for each dataset. For NoDecomp-Ctxt, we search
\(K \in \{6, 8, 10\}\) for GPT3 models and \(K \in 2, 4, 6, 8\) for
Flan-T5-* models. For Decomp-Ctxt, we search \(K \in \{2, 4, 6\}\) for
GPT3 and Flan-T5-* models. Note that the ranges are different between
GPT3 and Flan-T5-* as GPT3 can fit in more number of tokens. The ranges
are different for NoDecomp-Ctxt and Decomp-Ctxt as K refers to number of
paragraphs retrieved in each round of retrieval, and NoDecomp-Ctxt has
only one step of retrieval whereas Decomp-Ctxt usually has multiple
retrieval steps.

A.3 ADDITIONAL RESULTS\\
\textit{[Image omitted: images/5964bdfae9692537590f09e33d55115bf17a1981fab15d5293db039c1c502269.jpg]}

bar

MuSiQue \textbar{} Model \textbar{} No-Ctxt QA \textbar{} NoDecomp-Ctxt
QA \textbar{} Decomp-Ctxt QA \textbar{}
\textbar---\textbar---\textbar---\textbar---\textbar{} \textbar{}
Flan-T5-Large \textbar{} 8.4 \textbar{} 19.0 \textbar{} 19.6 \textbar{}
\textbar{} Flan-T5-XL \textbar{} 10.6 \textbar{} 21.7 \textbar{} 25.3
\textbar{} \textbar{} Flan-T5-XXL \textbar{} 10.3 \textbar{} 22.0
\textbar{} 26.8 \textbar{} \textbar{} code-davinci-002 \textbar{} 24.4
\textbar{} 29.7 \textbar{} 30.9 \textbar{}

\begin{enumerate}
\def\labelenumi{(\alph{enumi})}
\tightlist
\item
  CoT QA Prompt
\end{enumerate}

\textit{[Image omitted: images/0d5c12e36411eb0104f01da5ebc497b09f26a782f1ab15ab612fadc4bbc5c647.jpg]}

bar

MuSiQue \textbar{} Model \textbar{} No-Ctxt QA \textbar{} NoDecomp-Ctxt
QA \textbar{} Decomp-Ctxt QA \textbar{}
\textbar---\textbar---\textbar---\textbar---\textbar{} \textbar{}
Flan-T5-Large \textbar{} 7.9 \textbar{} 15.4 \textbar{} 18.8 \textbar{}
\textbar{} Flan-T5-XL \textbar{} 11.1 \textbar{} 22.3 \textbar{} 25.6
\textbar{} \textbar{} Flan-T5-XXL \textbar{} 15.1 \textbar{} 23.3
\textbar{} 29.3 \textbar{} \textbar{} code-davinci-002 \textbar{} 18.3
\textbar{} 21.4 \textbar{} 25.4 \textbar{}

\begin{enumerate}
\def\labelenumi{(\alph{enumi})}
\setcounter{enumi}{1}
\tightlist
\item
  Direct QA Prompt\\
  Figure 13: Results on MuSiQue dataset
\end{enumerate}

\textit{[Image omitted: images/45dba6e52f4f2aa115f3929de1d2963123ba131ec11b16264fd8537cabbf199b.jpg]}

bar

{\def\LTcaptype{none} % do not increment counter
\begin{longtable}[]{@{}llll@{}}
\toprule\noalign{}
Model & No-Ctxt QA & NoDecomp-Ctxt QA & Decomp-Ctxt QA \\
\midrule\noalign{}
\endhead
\bottomrule\noalign{}
\endlastfoot
Flan-T5-Large & 16.8 & 41.6 & 47.8 \\
Flan-T5-XL & 18.9 & 49.4 & 51.6 \\
Flan-T5-XXL & 21.3 & 49.1 & 51.8 \\
code-davinci-002 & 49.0 & 54.9 & 53.5 \\
\end{longtable}
}

HotpotQA

\begin{enumerate}
\def\labelenumi{(\alph{enumi})}
\tightlist
\item
  CoT QA Prompt
\end{enumerate}

\textit{[Image omitted: images/d4504b8e7a82f4890eea8d325a19d5630f073a2303aa4933e8ceaaa732aa5477.jpg]}

bar

{\def\LTcaptype{none} % do not increment counter
\begin{longtable}[]{@{}llll@{}}
\toprule\noalign{}
Model & No-Ctxt QA & NoDecomp-Ctxt QA & Decomp-Ctxt QA \\
\midrule\noalign{}
\endhead
\bottomrule\noalign{}
\endlastfoot
Flan-T5-Large & 19.3 & 36.6 & 42.7 \\
Flan-T5-XL & 22.6 & 47.9 & 52.7 \\
Flan-T5-XXL & 24.8 & 49.0 & 56.4 \\
code-davinci-002 & 41.2 & 50.2 & 49.9 \\
\end{longtable}
}

HotpotQA

\begin{enumerate}
\def\labelenumi{(\alph{enumi})}
\setcounter{enumi}{1}
\tightlist
\item
  Direct QA Prompt\\
  Figure 14: Results on HotpotQA dataset
\end{enumerate}

\textit{[Image omitted: images/2609a08967550a1743bf591908b659ba434f4202dcb6d48618953ca7a3f6c1db.jpg]}

bar

2WikiMultihopQA \textbar{} Model \textbar{} No-Ctxt QA \textbar{}
NoDecomp-Ctxt QA \textbar{} Decomp-Ctxt QA \textbar{}
\textbar---\textbar---\textbar---\textbar---\textbar{} \textbar{}
Flan-T5-Large \textbar{} 30.4 \textbar{} 38.9 \textbar{} 47.8 \textbar{}
\textbar{} Flan-T5-XL \textbar{} 31.4 \textbar{} 44.7 \textbar{} 58.2
\textbar{} \textbar{} Flan-T5-XXL \textbar{} 32.1 \textbar{} 48.8
\textbar{} 63.0 \textbar{} \textbar{} code-davinci-002 \textbar{} 40.9
\textbar{} 58.4 \textbar{} 70.8 \textbar{}

\begin{enumerate}
\def\labelenumi{(\alph{enumi})}
\tightlist
\item
  CoT QA Prompt
\end{enumerate}

\textit{[Image omitted: images/2bc371a43b220927690a10f2c4899ff40627372f338fbc0f43ed7422f79b7334.jpg]}

bar

2WikiMultihopQA \textbar{} Model \textbar{} No-Ctxt QA \textbar{}
NoDecomp-Ctxt QA \textbar{} Decomp-Ctxt QA \textbar{}
\textbar---\textbar---\textbar---\textbar---\textbar{} \textbar{}
Flan-T5-Large \textbar{} 27.3 \textbar{} 35.0 \textbar{} 48.9 \textbar{}
\textbar{} Flan-T5-XL \textbar{} 32.4 \textbar{} 44.4 \textbar{} 56.7
\textbar{} \textbar{} Flan-T5-XXL \textbar{} 32.3 \textbar{} 49.9
\textbar{} 63.3 \textbar{} \textbar{} code-davinci-002 \textbar{} 38.1
\textbar{} 47.1 \textbar{} 64.1 \textbar{}

\begin{enumerate}
\def\labelenumi{(\alph{enumi})}
\setcounter{enumi}{1}
\tightlist
\item
  Direct QA Prompt\\
  Figure 15: Results on 2WikiMultihopQA dataset
\end{enumerate}

\section{A.3.1 MUSIQUE}\label{a.3.1-musique}

We present all the results on the MuSiQue dataset in Fig. 13. Across all
settings, we can see that retrieval helps substantially (large gains
over No-Ctxt QA) with further improvements achieved by our DecomP-based
Decomp-Ctxt QA model.

\section{A.3.2 HOTPOTQA}\label{a.3.2-hotpotqa}

We present all the results on the HotpotQA dataset in Fig. 14. On this
dataset too, we can see large gains by incorporating retrieval but the
gains from using DecomP are mostly seen in the smaller models.

\section{A.3.3 2WIKIMULTIHOPQA}\label{a.3.3-2wikimultihopqa}

We present all the results on the 2WikiMultihopQA dataset in Fig. 15. On
this dataset, we can see large gains by incorporating retrieval and also
observe substantial gains by incorporating DecomP (as compared to
NoDecomp-Ctxt).

\section{B MATH QA}\label{b-math-qa}

We apply Decomposed Prompting to two math QA datasets: GSM8K Cobbe et
al.~(2021) and MultiArith Roy \& Roth (2015). For Chain-of-thought, we
used the original prompts for math reasoning Wei et al.~(2022). For
example:

Q: There are 15 trees in the grove. Grove workers will plant trees in
the grove today. After they are done, there will be 21 trees. How many
trees did the grove workers plant today? A: There are 15 trees
originally. Then there were 21 trees after some more were planted. So
there must have been 21 - 15 = 6. The answer is 6.

Most CoT systems Wei et al.~(2022); Wang et al.~(2023) rely on
extracting the answer by finding the number following ``answer is''.
However, this may not always be accurate. For example, the following CoT
would be unanswerable by relying on simple patterns.

Parker chews 4 pieces of gum a day. There are 15 pieces of gum in a
pack. So he will need 4 * 30 / 15 = 8 packs of gum to last him 30 days.

Rather than relying on patterns with limited generalization, we can use
a language model to extract the answer more reliably. Specifically, we
use Decomposed Prompting to decompose the task into first identifying
the chain-of-thought reasoning and then using a second GPT3-based
sub-module to extract the answer from the CoT. We show examples of our
prompts here (full prompt in App. G):

\section{Example from the Decomposition
Prompt}\label{example-from-the-decomposition-prompt}

QC: There are 15 trees in the grove. Grove workers will plant trees in
the grove today. After they are done, there will be 21 trees. How many
trees did the grove workers plant today?\\
QS: \hyperref[cot]{cot} There are 15 trees in the grove. Grove workers
will plant trees in the grove today. After they are done, there will be
21 trees. How many trees did the grove workers plant today?\\
A: There are 15 trees originally. Then there were 21 trees after some
more were planted. So there must have been 21 - 15 = 6 trees planted.\\
QS: \hyperref[gpt_ans]{gpt\_ans} There are 15 trees originally. Then
there were 21 trees after some more were planted. So there must have
been 21 - 15 = 6 trees planted.\\
A: 6\\
QS: {[}EOQ{]}

\section{Example from the gpt\_ans
prompt}\label{example-from-the-gpt_ans-prompt}

Q: There are 15 trees originally. Then there were 21 trees after some
more were planted. So there must have been 21 - 15 = 6 trees planted.

A: 6

\textit{[Image omitted: images/8683c0b832601df55712ac33cc6f79e8ed303cf3861607b6275b394d98be14ea.jpg]}

bar

{\def\LTcaptype{none} % do not increment counter
\begin{longtable}[]{@{}lll@{}}
\toprule\noalign{}
Model & CoT & DecomP \\
\midrule\noalign{}
\endhead
\bottomrule\noalign{}
\endlastfoot
GSM8K & 36.0 & 50.7 \\
MultiArith & 78.0 & 95.0 \\
\end{longtable}
}

Figure 16: Our simple decomposition results in 14-17 pts on two MathQA
datasets: GSM8k and MultiArith.

\textit{[Image omitted: images/c02fb75b74c69b439ad6b8f1f5d84a2d79536331ff33897a453410ba70ca7401.jpg]}

bar

{\def\LTcaptype{none} % do not increment counter
\begin{longtable}[]{@{}llll@{}}
\toprule\noalign{}
Category & CoT & DecomP(coarse) & DecomP(fine) \\
\midrule\noalign{}
\endhead
\bottomrule\noalign{}
\endlastfoot
text-curie-001 & 7 & 9 & 8 \\
text-davinci-001 & 10 & 14 & 19 \\
text-davinci-002 & 33 & 49 & 56 \\
\end{longtable}
}

Figure 17: As the models become weaker (davinci-001) and smaller
(curie-001), the performance of all the models drop. DECOMP still
outperforms CoT till the performance reaches close to zero with curie.

We present our results in Fig. 16. On the GSM8K data set \(^{10}\) , we
outperform CoT by 14 points. On the MultiArith dataset \(^{11}\) , we
achieve a 17 pt improvement compare to CoT. While this is a simple
change, it illustrates the possibility of using DECOMP for other complex
answer types, e.g.~non-extractive answer generation from
chain-of-thoughts.

\section{C EFFECT OF SCALE ON
COMMAQA}\label{c-effect-of-scale-on-commaqa}

We evaluate text-curie-001, text-davinci-001 and text-davinci-002 on the
CommAQA dataset. Since the curie-001 and davinci-001 have a smaller
context window size, we further reduced our prompts to fit within their
context windows (2048 tokens). As shown in Fig. 17, both CoT and DECOMP
are effected by the model size.

\section{D RESULTS ON ALL PROMPTS}\label{d-results-on-all-prompts}

\section{D.1 PER-PROMPT RESULT ON LETTER
CONCATENATION}\label{d.1-per-prompt-result-on-letter-concatenation}

We present the results of the letter concatenation task (with space
delimiter) for different values of N in Fig. 18. Our results are stable
across the different prompts (P1, P2 and P3) and always outperform CoT
and Least-to-Most prompting.

\textit{[Image omitted: images/b33cd81fed1f7af4cd3456d64ef299f3235500f54cb02fe648927df1b978710a.jpg]}

bar

{\def\LTcaptype{none} % do not increment counter
\begin{longtable}[]{@{}lllll@{}}
\toprule\noalign{}
Category & COT & COT w/ rollout & L2M w/ rollout & Decomp \\
\midrule\noalign{}
\endhead
\bottomrule\noalign{}
\endlastfoot
P1 & 27 & 82 & 83 & 98 \\
P2 & 22 & 74 & 75 & 98 \\
P3 & 19 & 58 & 66 & 98 \\
\end{longtable}
}

N=3

\textit{[Image omitted: images/8cde6d65d7deef962251ea0a1a82f436a14e340d559c0108023a6edb07f3537c.jpg]}

bar

{\def\LTcaptype{none} % do not increment counter
\begin{longtable}[]{@{}lllll@{}}
\toprule\noalign{}
Category & COT & COT w/ rollout & L2M w/ rollout & Decomp \\
\midrule\noalign{}
\endhead
\bottomrule\noalign{}
\endlastfoot
P1 & 16 & 68 & 64 & 96 \\
P2 & 11 & 45 & 45 & 96 \\
P3 & 9 & 48 & 55 & 96 \\
\end{longtable}
}

N=4

\textit{[Image omitted: images/d1329050481c53e831724e2a0e0f6136930e2b39c3c038e89a4f2394053824b4.jpg]}

bar

{\def\LTcaptype{none} % do not increment counter
\begin{longtable}[]{@{}lllll@{}}
\toprule\noalign{}
Category & COT & COT w/ rollout & L2M w/ rollout & Decomp \\
\midrule\noalign{}
\endhead
\bottomrule\noalign{}
\endlastfoot
P1 & 6 & 62 & 63 & 97 \\
P2 & 6 & 48 & 40 & 97 \\
P3 & 6 & 50 & 51 & 97 \\
\end{longtable}
}

N=5

Figure 18: Across all values of N and different prompts (P1, P2 and P3),
DECOMP outperform chain-of-thought reasoning and even least-to-most
prompting.

\section{D.2 PER-PROMPT RESULTS ON
COMMAQA}\label{d.2-per-prompt-results-on-commaqa}

We also present the results of all the prompts on the CommAQA dataset in
Fig. 19. Here too, we can observe that DECOMP outperforms CoT on each
prompt set.

\textit{[Image omitted: images/aebc4be5b0c4b5e0dd15a9c1f88ef3128a07a4e9c131f755a874724403e79cd0.jpg]}

bar

{\def\LTcaptype{none} % do not increment counter
\begin{longtable}[]{@{}llll@{}}
\toprule\noalign{}
Category & COT & Decomp (coarse) & Decomp (fine) \\
\midrule\noalign{}
\endhead
\bottomrule\noalign{}
\endlastfoot
P1 & 41 & 54 & 58 \\
P2 & 46 & 52 & 60 \\
P3 & 39 & 59 & 61 \\
\end{longtable}
}

\begin{enumerate}
\def\labelenumi{(\alph{enumi})}
\tightlist
\item
  Test Set
\end{enumerate}

\textit{[Image omitted: images/7a5b8894af4714c78c9d787d5d66f4f5253e9eecee1df3080ee6a7bb4c400f9c.jpg]}

bar

{\def\LTcaptype{none} % do not increment counter
\begin{longtable}[]{@{}llll@{}}
\toprule\noalign{}
Category & COT & Decomposition (coarse) & Decomposition (fine) \\
\midrule\noalign{}
\endhead
\bottomrule\noalign{}
\endlastfoot
P1 & 37 & 60 & 62 \\
P2 & 34 & 58 & 61 \\
P3 & 30 & 61 & 70 \\
\end{longtable}
}

\begin{enumerate}
\def\labelenumi{(\alph{enumi})}
\setcounter{enumi}{1}
\tightlist
\item
  Comp. Gen.~Set\\
  Figure 19: Results of different prompts on the CommAQA dataset.
\end{enumerate}

\section{E EFFECT OF DECOMPOSITION
SCHEME}\label{e-effect-of-decomposition-scheme}

To evaluate the effect of the decomposition scheme, we experiment with
two other simple decomposition structures for the letter concatenation
and reversal tasks.

Letter Concatenation For letter concatenation, we consider an alternate
scheme where we use GPT3 to generate each question rather than loop over
the answers, e.g.,

QC: Take the last letters of the words in ``Augusta Ada King'' and
concatenate them using a space.\\
QS: \hyperref[split]{split} What are the words in ``Augusta Ada
King''?\\
A: {[}``Augusta'', ``Ada'', ``King''{]}\\
QS: \hyperref[str_position]{str\_position} What is the last letter in
``Augusta''?\\
A: ``a''\\
QS: \hyperref[str_position]{str\_position} What is the last letter in
``Ada''?\\
A: ``a''\\
QS: \hyperref[str_position]{str\_position} What is the last letter in
``King''?\\
A: ``g''\\
QS: \hyperref[merge]{merge} Concatenate {[}``a'', ``a'', ``g''{]} using
a space.\\
A: ``a a g''\\
QS: {[}EOQ{]}

By using the decomposer prompt model to generate the sub-questions, we
can be more robust to formatting issues in the output answers, e.g., we
can expect GPT3 to still generate the appropriate sub-questions even if
the first answer is not a valid array. However, the generated
sub-questions may not correctly use all the elements of the list (change
in order, missed element, repeated elements, etc).

List Reversal For list reversal, instead of splitting into halves, we
take the tail of the list, reverse it and then concatenate it to the
head. i.e.~reverse(list) = reverse(list{[}1:{]}) + list{[}0{]}. This
requires more GPT3 calls (O(n)) compared to the original approach of
splitting the list into halves (O(log(n))).

In both these cases, we noticed that the performance did not drop as
shown in Fig. 20 and Fig. 21. On the letter concatenation task, the
results were exactly the same. The new reversal decomposition schema was
actually stronger on longer inputs at the cost of more calls to GPT3
(O(ln(n)) using binary splits vs O(n) one element at a time). Both these
decomposition schemes are still better than CoT.

\textit{[Image omitted: images/fbf1dbe12dcaa944a9ae8ddd96c406b6ab01669763970962e77998467e091469.jpg]}

bar

{\def\LTcaptype{none} % do not increment counter
\begin{longtable}[]{@{}lll@{}}
\toprule\noalign{}
Group & Decomp (Loop) & Decomp (Generate) \\
\midrule\noalign{}
\endhead
\bottomrule\noalign{}
\endlastfoot
N=3 & 98 & 98 \\
N=4 & 96 & 96 \\
N=5 & 97 & 97 \\
\end{longtable}
}

Figure 20: Both decomposition schemes for the letter concatenation task
have the same scores.

\textit{[Image omitted: images/48b2ec5d64843772f77b28627425f3dd2ac0d410635f8b595fc0f97fe97128fe.jpg]}

bar

{\def\LTcaptype{none} % do not increment counter
\begin{longtable}[]{@{}lll@{}}
\toprule\noalign{}
N & Decomp (Mid Split) & Decomp (Tail Split) \\
\midrule\noalign{}
\endhead
\bottomrule\noalign{}
\endlastfoot
4 & 86 & 79 \\
6 & 59 & 80 \\
8 & 64 & 74 \\
10 & 42 & 67 \\
\end{longtable}
}

Figure 21: Recursively reversing the tail of a list is more stable at
longer lengths but comes at the cost of more calls to GPT3.

\section{F ERROR ANALYSIS}\label{f-error-analysis}

\section{F.1 LETTER CONCATENATION}\label{f.1-letter-concatenation}

\section{F.1.1 DECOMP}\label{f.1.1-decomp}

We analyzed the errors in DECOMP on the letter concatenation task and
only found errors in the sub-task execution.

\textit{[Image omitted: images/3373651a862fa23f1acd37b0425d8ec01261c2f9db5069eb93c5c25b054a542f.jpg]}

bar

{\def\LTcaptype{none} % do not increment counter
\begin{longtable}[]{@{}llll@{}}
\toprule\noalign{}
Category & CoT & CoT w/ rollout & DecomP \\
\midrule\noalign{}
\endhead
\bottomrule\noalign{}
\endlastfoot
N=3(IID) & 31 & 82 & 97 \\
N=4(OOD) & 12 & 70 & 94 \\
N=5(OOD) & 7 & 51 & 98 \\
\end{longtable}
}

Figure 22: EM Results on the \(k^{th}\) letter concatenation task (k=3)
using semi-colon as delimiter with different values for N, the number of
words in the input. DECOMP always outperforms and generalizes better
than CoT.

Q: Take the letters at position 3 of the words in ``Nancy Samina Abbas
Caudhari Bano'' and concatenate them using a space.\\
A: n m b u n\\
Prediction: c m b u n\\
Error: Incorrect letter extraction (Sub-task)\\
What is at position 3 in ``{[} \(N, 1\) , \(a, 2\) , \(n, 3\) , \(c, 4\)
, \(y, 5\) {]}''? ⇒ ``c''\\
Q: Take the letters at position 3 of the words in ``Orlando Stephen Cho
Teixeira Pierre'' and concatenate them using a space.\\
A: l e o i e\\
Prediction: leoie\\
Error: Incorrect concatenation (Sub-task)\\
Concatenate {[}``l'', ``e'', ``o'', ``i'', ``e''{]} using a space.
\(\Rightarrow\) ``leoie''

\section{F.1.2 CoT w/ ROLLOUT}\label{f.1.2-cot-w-rollout}

We analyzed the errors in CoT on the letter concatenation task and found
similar errors during the generation of CoT. But the frequency of these
errors was higher than DECOMP, as it is not possible to effectively
teach each sub-task with CoT.

Q: Take the letters at position 3 of the words in ``Sheila Nicolas Verma
Sha Sousa'' and concatenate them using a space.\\
A: e c r a u\\
Pred: icra u\\
Error: Incorrect letter extraction\\
\ldots The letters and their positions in ``Sheila'' are ``{[}(S, 1),
(h, 2), (e, 3), (i, 4), (l, 5), (a, 6){]}''. The letter at position 3 in
this sequence is ``i''\ldots{}\\
Q: Take the letters at position 3 of the words in ``Shobha Kailash
Nakamura Peter Benitez'' and concatenate them using a space.\\
A: oiktn\\
Pred: o l k t i\\
Error: Incorrect letter extraction\\
\ldots''Benitez'' are ''{[}(B, 1), (e, 2), (n, 3), (i, 4), (t, 5), (e,
6), (z, 7){]}''. The letter at position 3 in this sequence is
''i''\ldots{}

\section{F.2 COMMAQA}\label{f.2-commaqa}

Similarly in CommaQA, the errors are mostly due to sub-task errors,
which in this dataset correspond to answering single-hop questions. CoT
also makes the same types of errors but they are more frequent since
this QA sub-task can not be delegated to a specialized prompt in CoT.
Since all errors are of this type, we show only one example here.

\begin{Shaded}
\begin{Highlighting}[]
\FunctionTok{Q}\KeywordTok{:}\AttributeTok{ What awards have movies written by people born in 1933 won?}
\FunctionTok{A}\KeywordTok{:}\AttributeTok{ }\KeywordTok{[}\StringTok{"Hydrallium"}\KeywordTok{,}\AttributeTok{ }\StringTok{"Pompasole"}\KeywordTok{]}
\FunctionTok{Pred}\KeywordTok{:}\AttributeTok{ }\KeywordTok{[}\StringTok{"Pompasole"}\KeywordTok{]}
\FunctionTok{Error}\KeywordTok{:}\AttributeTok{ Incorrect sub{-}question answer}
\FunctionTok{Sub{-}Q}\KeywordTok{:}\AttributeTok{ What movies has Haldron written?}
\FunctionTok{Sub{-}A}\KeywordTok{:}\AttributeTok{ }\KeywordTok{[}\StringTok{"Polytetrafluoromethane"}\KeywordTok{,}\AttributeTok{ }\StringTok{"Skia"}\KeywordTok{,}\AttributeTok{ }\StringTok{"Warpstone"}\KeywordTok{]}
\FunctionTok{Pred}\KeywordTok{:}\AttributeTok{ }\KeywordTok{[}\StringTok{"Skia"}\KeywordTok{,}\AttributeTok{ }\StringTok{"Warpstone"}\KeywordTok{]}\AttributeTok{ }
\end{Highlighting}
\end{Shaded}

\section{G TASK PROMPTS}\label{g-task-prompts}

We have provided the task prompts for all the datasets for COT and our
Decomposed Prompting approach.

CoT Since CoT methods also perform 2-step reasoning: first generate the
chain-of-thought and second extract the answer from the CoT, we use the
same decomposition-based framework for COT baselines too. For example,
consider the following example in our COT prompt:

\begin{Shaded}
\begin{Highlighting}[]
\NormalTok{QC: Take the letters at position 1 of the words in "Alan Mathison Turing" and concatenate them using a space.}
\NormalTok{QS: [extract] The letter at position 1 of "Alan" is "A". The letter at position 1 of "Mathison" is "M". The letter at position 1 of "Turing" is "T". Concatenating "A", "M", "T" using a space leads to "A M T". So, "Alan Mathison Turing" outputs "A M T".}
\NormalTok{A: "A M T"}
\NormalTok{QS: [EOQ] }
\end{Highlighting}
\end{Shaded}

GPT3 generates the chain-of-thought during the ``decomposition'' step
and a regex-based answer extractor extract (`.* outputs ``(.*)''.' )
then takes this CoT and generates the answer. In some cases, the module
name is skipped in the prompt (the CoT is sent to the extractor by
default).

Operators In this work, we use the same operators as defined by Khot et
al.. Their select operator is just the basic operator that replaces
references to an answer index with its answer. When not specified,
select is assumed to be the default operator. In addition, we consider
two operators in our experiments: project\_values and
project\_values\_flat\_unique.

\begin{itemize}
\tightlist
\item
  project\_values: This operator takes a list answer \#i = X and
  iterates over it to generate new questions by replacing mentions of
  \#i i.e.~Q = {[}q.replace(\#i, x) for x ∈ X{]}. The answer to each
  question is simply concatenated to get the final answer i.e.~A =
  {[}model(q) for q ∈ Q{]}. We refer to this as foreach for simplicity
  in the main text.\\
\item
  project\_values\_flat\_unique: This operator performs the same steps
  as project\_values but then additionally flattens the list and only
  returns the unique entities in the flattened list. We refer to this as
  foreach\_merge in the main text for simplicity.
\end{itemize}

\section{G.1 LETTER CONCATENTATION}\label{g.1-letter-concatentation}

We show one of the prompts used for experiments here. The entire set of
prompts is provided as supplementary material.

\section{G.1.1 DECOMPOSED PROMPTING}\label{g.1.1-decomposed-prompting}

\section{decomp}\label{decomp}

QC: Take the last letters of the words in ``Augusta Ada King'' and
concatenate them using a space.

QS: \hyperref[split]{split} What are the words in ``Augusta Ada King''?

A: {[}``Augusta'', ``Ada'', ``King''{]}

QS: (project\_values) \hyperref[str_position]{str\_position} What is the
last letter in ``\#1''?

A: {[}``a'', ``a'', ``g''{]}

QS: \hyperref[merge]{merge} Concatenate \#2 using a space.

A: ``a a g''

QS: {[}EOQ{]}

QC: Take the letters at position 1 of the words in ``Alan Mathison
Turing'' and concatenate them using a space.

QS: \hyperref[split]{split} What are the words in ``Alan Mathison
Turing''?

A: {[}``Alan'', ``Mathison'', ``Turing''{]}

QS: (project\_values) \hyperref[str_position]{str\_position} What is the
letter at position 1 in ``\#1''?

A: {[}``A'', ``M'', ``T''{]}

QS: \hyperref[merge]{merge} Concatenate \#2 using a space.

A: ``A M T''

QS: {[}EOQ{]}

QC: Take the letters at position 4 of the words in ``Herbert Alexander
Simon'' and concatenate them using a space.

QS: \hyperref[split]{split} What are the words in ``Herbert Alexander
Simon''?

A: {[}``Herbert'', ``Alexander'', ``Simon''{]}

QS: (project\_values) \hyperref[str_position]{str\_position} What is the
letter at position 4 in ``\#1''?

A: {[}``b'', ``x'', ``o''{]}

QS: \hyperref[merge]{merge} Concatenate \#2 using a space.

A: ``b x o''

QS: {[}EOQ{]}

\section{split}\label{split}

Q: What are the words in ``Alan Mathison Turing''?

A: {[}``Alan'', ``Mathison'', ``Turing''{]}

Q: What are the letters in ``Alan''?

A: {[}``A'', ``l'', ``a'', ``n''{]}

Q: What are the letters and their positions in ``Mathison''?

A: ''{[}(M, 1), (a, 2), (t, 3), (h, 4), (i, 5), (s, 6), (o, 7), (n,
8){]}''

Q: What are the words and their positions in ``Herbert Alexander
Simon''?

A: ''{[}(Herbert, 1), (Alexander, 2), (Simon, 3){]}''

\section{str\_position}\label{str_position}

QC: What is the letter at position 1 of the word ``Augusta''?

QS: (select) \hyperref[split]{split} What are the letters and their
positions in ``Augusta''?

A: ''{[}(A, 1), (u, 2), (g, 3), (u, 4), (s, 5), (t, 6), (a, 7){]}''

QS: (select) \hyperref[arr_position]{arr\_position} What is at position
1 in \#1?

A: ``A''

QS: {[}EOQ{]}

QC: What is the last letter of the word ``Mathison''?

QS: (select) \hyperref[split]{split} What are the letters and their
positions in ``Mathison''?

A: ''{[}(M, 1), (a, 2), (t, 3), (h, 4), (i, 5), (s, 6), (o, 7), (n,
8){]}''

QS: (select) \hyperref[arr_position]{arr\_position} What is the last
letter in \#1?

A: ``n''

QS: {[}EOQ{]}

QC: What is the word at the position 4 in ``Colorless green ideas sleep
furiously''?

QS: (select) \hyperref[split]{split} What are the words and their
positions in ``Colorless green ideas sleep furiously''?

A: ''{[}(Colorless, 1), (green, 2), (ideas, 3), (sleep, 4), (furiously,
5){]}''

QS: (select) \hyperref[arr_position]{arr\_position} What is at the
position 4 in \#1?

A: ``sleep''

QS: {[}EOQ{]}

\section{merge}\label{merge}

Q: Concatenate {[}``A'', ``l'', ``a'', ``n''{]}.

A: ``Alan''

Q: Concatenate {[}``b'', ``x'', ``o''{]} using a space.

A: ``b x o''

Q: Concatenate {[}``a'', ``a'', ``g''{]} using a comma.

A: ``a,a,g''

Q: Concatenate {[}``Alan'', ``Mathison'', ``Turing''{]} using a space.

A: ``Alan Mathison Turing''

Q: Concatenate {[}``Allen'', ``Institute''{]}.

A: ``AllenInstitute''

\section{arr\_position}\label{arr_position}

Q: What is at position 4 in ''{[}(``Colorless'', 1), (``green'', 2),
(``ideas'', 3), (``sleep'', 4), (``furiously'', 5){]}''?

A: ``sleep''

Q: What is at position 1 in ''{[}(M, 1), (a, 2), (t, 3), (h, 4), (i, 5),
(s, 6), (o, 7), (n, 8){]}''?

A: ``M''

Q: What is at the last position in ''{[}(A, 1), (u, 2), (g, 3), (u, 4),
(s, 5), (t, 6), (a, 7){]}''?

A: ``a''

Q: What is at position 1 in ''{[}(Herbert, 1), (Alexander, 2), (Simon,
3){]}''?

A: ``Herbert''

Q: What is at last position in ''{[}(Allen, 1), (Institute, 2), (for,
3), (Artificial, 4), (Intelligence, 5){]}''?

A: ``Intelligence''

Q: What is at position 4 in ''{[}(A, 1), (1, 2), (e, 3), (x, 4), (a, 5),
(n, 6), (d, 7), (e, 8), (r, 9){]}''?

A: ``x''

\section{G.1.2 COT WITH ROLLOUT}\label{g.1.2-cot-with-rollout}

\section{COT w/ rollout}\label{cot-w-rollout}

QC: Take the last letters of the words in ``Augusta Ada King'' and
concatenate them using a space.

QS: The words in ``Augusta Ada King'' are ``Augusta'', ``Ada'' and
``King''. The letters and their positions in ``Augusta'' are ``{[}(A,
1), (u, 2), (g, 3), (u, 4), (s, 5), (t, 6), (a, 7){]}''. The last letter
in this sequence is ``a''. The letters and their positions in ``Ada''
are ``{[}(A, 1), (d, 2), (a, 3){]}''. The last letter in this sequence
is ``a''. The letters and their positions in ``King'' are ``{[}(K, 1),
(i, 2), (n, 3), (g, 4){]}''. The last letter in this sequence is ``g''.
Concatenating ``a'', ``a'', ``g'' using a space leads to ``a a g''. So,
``Augusta Ada King'' outputs ``a a g''.

A: ``a a g''

QS: {[}EOQ{]}

QC: Take the letters at position 1 of the words in ``Alan Mathison
Turing'' and concatenate them using a space.

QS: The words in ``Alan Mathison Turing'' are ``Alan'', ``Mathison'',
and ``Turing''. The letters and their positions in ``Alan'' are ``{[}(A,
1), (l, 2), (a, 3), (n, 4){]}''. The letter at position 1 in this
sequence is ``A''. The letters and their positions in ``Mathison'' are
{[}(M, 1), (a, 2), (t, 3), (h, 4), (i, 5), (s, 6), (o, 7), (n, 8){]}``.
The letter at position 1 in this sequence is''M''. The letters and their
positions in ``Turing'' are ``{[}(T, 1), (u, 2), (r, 3), (i, 4), (n, 5),
(g, 6){]}''. The letter at position 1 in this sequence is ``T''.
Concatenating ``A'', ``M'', ``T'' using a space leads to ``A M T''. So,
``Alan Mathison Turing'' outputs ``A M T''.

A: ``A M T''

QS: {[}EOQ{]}

QC: Take the letters at position 4 of the words in ``Herbert Alexander
Simon'' and concatenate them using a space.

QS: The words in ``Herbert Alexander Simon'' are ``Herbert'',
``Alexander'', and ``Simon''. The letters and their positions in
``Herbert'' are ``{[}(H, 1), (e, 2), (r, 3), (b, 4), (e, 5), (r, 6), (t,
7){]}''. The letter at position 4 in this sequence is ``b''. The letters
and their positions in ``Alexander'' are ``{[}(A, 1), (l, 2), (e, 3),
(x, 4), (a, 5), (n, 6), (d, 7), (e, 8), (r, 9){]}''. The letter at
position 4 in this sequence is ``x''. The letters and their positions in
``Simon'' are ``{[}(S, 1), (i, 2), (m, 3), (o, 4), (n, 5){]}''. The
letter at position 4 in this sequence is ``o''. Concatenating ``b'',
``x'', ``o'' using a space leads to ``b x o''. So, ``Herbert Alexander
Simon'' outputs ``b x o''.

A: ``b x o''

QS: {[}EOQ{]}

\section{G.1.3 COT}\label{g.1.3-cot}

\section{COT}\label{cot}

QC: Take the last letters of the words in ``Augusta Ada King'' and
concatenate them using a space.

QS: The last letter of ``Augusta'' is ``a''. The last letter of ``Ada''
is ``a''. The last letter of ``King'' is ``g''. Concatenating ``a'',
``a'', ``g'' using a space leads to ``a a g''. So, ``Augusta Ada King''
outputs ``a a g''.

A: ``a a g''

QS: {[}EOQ{]}

QC: Take the letters at position 1 of the words in ``Alan Mathison
Turing'' and concatenate them using a space.

QS: The letter at position 1 of ``Alan'' is ``A''. The letter at
position 1 of ``Mathison'' is ``M''. The letter at position 1 of
``Turing'' is ``T''. Concatenating ``A'', ``M'', ``T'' using a space
leads to ``A M T''. So, ``Alan Mathison Turing'' outputs ``A M T''.

A: ``A M T''

QS: {[}EOQ{]}

QC: Take the letters at position 4 of the words in ``Herbert Alexander
Simon'' and concatenate them using a space.

QS: The letter at position 4 of ``Herbert'' is ``b''. The letter at
position 4 of ``Alexander'' is ``x''. The letter at position 4 of
``Simon'' is ``o''. Concatenating ``b'', ``x'', ``o'' using a space
leads to ``b x o''. So, ``Herbert Alexander Simon'' outputs ``b x o''.

A: ``b x o''

QS: {[}EOQ{]}

\section{G.1.4 LEAST-TO-MOST W/
ROLLOUT}\label{g.1.4-least-to-most-w-rollout}

\section{Least-to-most Decomp}\label{least-to-most-decomp}

QC: Take the last letters of the words in ``Augusta Ada King'' and
concatenate them using a space.

QS: {[}12m{]} Take the last letters of the words in ``Augusta Ada'' and
concatenate them using a space.

A: The words in ``Augusta Ada'' are ``Augusta'' and ``Ada''. The letters
and their positions in ``Augusta'' are ``{[}(A, 1), (u, 2), (g, 3), (u,
4), (s, 5), (t, 6), (a, 7){]}''. The last letter in this sequence

is ``a''. The letters and their positions in ``Ada'' are ``{[}(A, 1),
(d, 2), (a, 3){]}''. The last letter in this sequence is ``a''.
Concatenating ``a'', ``a'' using a space leads to ``a a''. So, ``Augusta
Ada'' outputs ``a a''.\\
QS: {[}l2m{]} Take the last letters of the words in ``Augusta Ada King''
and concatenate them using a space.\\
A: ``Augusta Ada'' outputs ``a a''. The letters and their positions in
``King'' are ``{[}(K, 1), (i, 2), (n, 3), (g, 4){]}''. The last letter
in this sequence is ``g''. Concatenating ``a a'', ``g'' using a space
leads to ``a a g''. So, ``Augusta Ada King'' outputs ``a a g''.\\
QS: {[}ans\_ext{]} So, ``Augusta Ada King'' outputs ``a a g''.\\
A: ``a a g''\\
QS: {[}EOQ{]}\\
QC: Take the letters at position 1 of the words in ``Alan Mathison
Turing'' and concatenate them using a space.\\
QS: {[}l2m{]} Take the letters at position 1 of the words in ``Alan
Mathison'' and concatenate them using a space.\\
A: The words in ``Alan Mathison'' are ``Alan'' and ``Mathison''. The
letters and their positions in ``Alan'' are ``{[}(A, 1), (l, 2), (a, 3),
(n, 4){]}''. The letter at position 1 in this sequence is ``A''. The
letters and their positions in ``Mathison'' are {[}(M, 1), (a, 2), (t,
3), (h, 4), (i, 5), (s, 6), (o, 7), (n, 8){]}``. The letter at position
1 in this sequence is''M''. Concatenating ``A'', ``M'' using a space
leads to ``A M''. So, ``Alan Mathison'' outputs ``A M''.\\
QS: {[}12m{]} Take the letters at position 1 of the words in ``Alan
Mathison Turing'' and concatenate them using a space.\\
A: ``Alan Mathison'' outputs ``A M''. The letters and their positions in
``Turing'' are ``{[}(T, 1), (u, 2), (r, 3), (i, 4), (n, 5), (g, 6){]}''.
The letter at position 1 in this sequence is ``T''. Concatenating ``A
M'', ``T'' using a space leads to ``A M T''. So, ``Alan Mathison
Turing'' outputs ``A M T''.\\
QS: {[}ans\_ext{]} So, ``Alan Mathison Turing'' outputs ``A M T''.\\
A: ``A M T''\\
QS: {[}EOQ{]}\\
QC: Take the letters at position 4 of the words in ``Herbert Alexander
Simon'' and concatenate them using a space.\\
QS: {[}12m{]} Take the letters at position 4 of the words in ``Herbert
Alexander'' and concatenate them using a space.\\
A: The words in ``Herbert Alexander'' are ``Herbert'' and ``Alexander''.
The letters and their positions in ``Herbert'' are ``{[}(H, 1), (e, 2),
(r, 3), (b, 4), (e, 5), (r, 6), (t, 7){]}''. The letter at position 4 in
this sequence is ``b''. The letters and their positions in ``Alexander''
are ``{[}(A, 1), (l, 2), (e, 3), (x, 4), (a, 5), (n, 6), (d, 7), (e, 8),
(r, 9){]}''. The letter at position 4 in this sequence is ``x''.
Concatenating ``b'', ``x'' using a space leads to ``b x''. So, ``Herbert
Alexander'' outputs ``b x''.\\
QS: {[}12m{]} Take the letters at position 4 of the words in ``Herbert
Alexander Simon'' and concatenate them using a space.\\
A: ``Herbert Alexander'' outputs ``b x''. The letters and their
positions in ``Simon'' are ``{[}(S, 1), (i, 2), (m, 3), (o, 4), (n,
5){]}''. The letter at position 4 in this sequence is ``o''.
Concatenating ``b x'', ``o'' using a space leads to ``b x o''. So,
``Herbert Alexander Simon'' outputs ``b x o''.\\
QS: {[}ans\_ext{]} So, ``Herbert Alexander Simon'' outputs ``b x o''.\\
A: ``b x o''\\
QS: {[}EOQ{]}

\section{Least-to-most COT(12m)}\label{least-to-most-cot12m}

Q: Take the last letters of the words in ``Augusta Ada'' and concatenate
them using a space.

A: The words in ``Augusta Ada'' are ``Augusta'' and ``Ada''. The letters
and their positions in ``Augusta'' are ``{[}(A, 1), (u, 2), (g, 3), (u,
4), (s, 5), (t, 6), (a, 7){]}''. The last letter in this sequence is
``a''. The letters and their positions in ``Ada'' are ``{[}(A, 1), (d,
2), (a, 3){]}''. The last letter in this sequence is ``a''.
Concatenating ``a'', ``a'' using a space leads to ``a a''. So, ``Augusta
Ada'' outputs ``a a''.

Q: Take the last letters of the words in ``Augusta Ada King'' and
concatenate them using a space.

A: ``Augusta Ada'' outputs ``a a''. The letters and their positions in
``King'' are ``{[}(K, 1), (i, 2), (n, 3), (g, 4){]}''. The last letter
in this sequence is ``g''. Concatenating ``a a'', ``g'' using a space
leads to ``a a g''. So, ``Augusta Ada King'' outputs ``a a g''.\\
Q: Take the letters at position 1 of the words in ``Alan Mathison'' and
concatenate them using a space.\\
A: The words in ``Alan Mathison'' are ``Alan'' and ``Mathison''. The
letters and their positions in ``Alan'' are ``{[}(A, 1), (l, 2), (a, 3),
(n, 4){]}''. The letter at position 1 in this sequence is ``A''. The
letters and their positions in ``Mathison'' are {[}(M, 1), (a, 2), (t,
3), (h, 4), (i, 5), (s, 6), (o, 7), (n, 8){]}``. The letter at position
1 in this sequence is''M''. Concatenating ``A'', ``M'' using a space
leads to ``A M''. So, ``Alan Mathison'' outputs ``A M''.\\
Q: Take the letters at position 1 of the words in ``Alan Mathison
Turing'' and concatenate them using a space.\\
A: ``Alan Mathison'' outputs ``A M''. The letters and their positions in
``Turing'' are ``{[}(T, 1), (u, 2), (r, 3), (i, 4), (n, 5), (g, 6){]}''.
The letter at position 1 in this sequence is ``T''. Concatenating ``A
M'', ``T'' using a space leads to ``A M T''. So, ``Alan Mathison
Turing'' outputs ``A M T''.\\
Q: Take the letters at position 4 of the words in ``Herbert Alexander''
and concatenate them using a space.\\
A: The words in ``Herbert Alexander'' are ``Herbert'' and ``Alexander''.
The letters and their positions in ``Herbert'' are ``{[}(H, 1), (e, 2),
(r, 3), (b, 4), (e, 5), (r, 6), (t, 7){]}''. The letter at position 4 in
this sequence is ``b''. The letters and their positions in ``Alexander''
are ``{[}(A, 1), (l, 2), (e, 3), (x, 4), (a, 5), (n, 6), (d, 7), (e, 8),
(r, 9){]}''. The letter at position 4 in this sequence is ``x''.
Concatenating ``b'', ``x'' using a space leads to ``b x''. So, ``Herbert
Alexander'' outputs ``b x''.\\
Q: Take the letters at position 4 of the words in ``Herbert Alexander
Simon'' and concatenate them using a space.\\
A: ``Herbert Alexander'' outputs ``b x''. The letters and their
positions in ``Simon'' are ``{[}(S, 1), (i, 2), (m, 3), (o, 4), (n,
5){]}''. The letter at position 4 in this sequence is ``o''.
Concatenating ``b x'', ``o'' using a space leads to ``b x o''. So,
``Herbert Alexander Simon'' outputs ``b x o''.

\section{G.1.5 ALT DECOMP SCHEMA (GENERATE EACH
SUB-QUESTION)}\label{g.1.5-alt-decomp-schema-generate-each-sub-question}

\section{decomp}\label{decomp-1}

QC: Take the last letters of the words in ``Augusta Ada King'' and
concatenate them using a space.

QS: \hyperref[split]{split} What are the words in ``Augusta Ada King''?

A: {[}``Augusta'', ``Ada'', ``King''{]}

QS: \hyperref[str_position]{str\_position} What is the last letter in
``Augusta''?

A: ``a''

QS: \hyperref[str_position]{str\_position} What is the last letter in
``Ada''?

A: ``a''

QS: \hyperref[str_position]{str\_position} What is the last letter in
``King''?

A: ``g''

QS: \hyperref[merge]{merge} Concatenate {[}``a'', ``a'', ``g''{]} using
a space.

A: ``a a g''

QS: {[}EOQ{]}

QC: Take the letters at position 1 of the words in ``Alan Mathison
Turing'' and concatenate them using a space.

QS: \hyperref[split]{split} What are the words in ``Alan Mathison
Turing''?

A: {[}``Alan'', ``Mathison'', ``Turing''{]}

QS: \hyperref[str_position]{str\_position} What is the letter at
position 1 in ``Alan''?

A: ``A''

QS: \hyperref[str_position]{str\_position} What is the letter at
position 1 in ``Mathison''?

A: ``M''

QS: \hyperref[str_position]{str\_position} What is the letter at
position 1 in ``Turing''?

A: ``T''

QS: \hyperref[merge]{merge} Concatenate {[}``A'', ``M'', ``T''{]} using
a space.

A: ``A M T''

QS: {[}EOQ{]}

QC: Take the letters at position 4 of the words in ``Herbert Alexander
Simon'' and concatenate them using a space.

QS: \hyperref[split]{split} What are the words in ``Herbert Alexander
Simon''?

A: {[}``Herbert'', ``Alexander'', ``Simon''{]}

QS: \hyperref[str_position]{str\_position} What is the letter at
position 4 in ``Herbert''?

A: ``b''

QS: \hyperref[str_position]{str\_position} What is the letter at
position 4 in ``Alexander''?

A: ``x''

QS: \hyperref[str_position]{str\_position} What is the letter at
position 4 in ``Simon''?

A: ``o''

QS: \hyperref[merge]{merge} Concatenate {[}``b'', ``x'', ``o''{]} using
a space.

A: ``b x o''

QS: {[}EOQ{]}

\section{G.2 SEQUENCE REVERSAL}\label{g.2-sequence-reversal}

\section{G.2.1 SPLIT REVERSAL}\label{g.2.1-split-reversal}

The prompts in this section implement Algorithm 1.

Algorithm 1 A recursive reversal strategy that splits the sequence in
half, reverses each half, and concatenates them. Runs in \(O(\log n)\)
calls to the LM where n is the number of items in the sequence.\\
procedure SPLITREVERSE(x) ▷ The reversal of x if \textbar x\textbar{}
\textless{} 4 then ▷ The base case return x\textbar x\textbar, \ldots,
x1 ▷ The reversed sequence else ▷ The inductive case n ←
\textbar x\textbar/2 ▷ Half the length of x ℓ ← x1, \ldots, xn ▷ The
first half of x ℓR ← SPLITREVERSE(ℓ) ▷ The reversed first half r ← xn+1,
\ldots, x\textbar x\textbar{} ▷ The second half of x rR ←
SPLITREVERSE(r) ▷ The reversed second half return r1R, \ldots, rnR,
ℓn+1, \ldots, ℓR\textbar x\textbar{} ▷ The concatenated reversed halves
end if end procedure

\section{reverse}\label{reverse}

QC: Reverse the sequence ``driving license, button, packet, identity
card, shoe''.

QS: {[}extract{]} The sequence is ``1. driving license, 2. button, 3.
packet, 4. identity card, 5. shoe''. The sequence is 5 items long, which
is more than the minimum length of 4, so we split it. Half of 5 is 5 / 2
= 2.5. Dropping the decimal, we get that the first half will be 2 items
long, ending in ``2. button''. The first half (2 items) is ``1. driving
license, 2. button''.

A: ``1. driving license, 2. button''

QS: {[}extract{]} The first half of the sequence ends with ``2.
button'', so the second half starts after ``2. button'' with ``3.
packet''. The full sequence is 5 items long, and the first half is 2
items long, so the second half will be 5 - 2 = 3 items long. The second
half of the sequence (3 items) is ``3. packet, 4. identity card, 5.
shoe''.

A: ``3. packet, 4. identity card, 5. shoe''

QS: \hyperref[remove_numbers]{remove\_numbers} Remove the numbers from
\#1.

A: ``driving license, button''

QS: \hyperref[remove_numbers]{remove\_numbers} Remove the numbers from
\#2.

A: ``packet, identity card, shoe''

QS: \hyperref[reverse]{reverse} Reverse the sequence \#3.

A: ``button, driving license''

QS: \hyperref[reverse]{reverse} Reverse the sequence \#4.

A: ``shoe, identity card, packet''

QS: \hyperref[join]{join} \#6 \#5

A: ``shoe, identity card, packet, button, driving license''

QS: {[}EOQ{]}

QC: Reverse the sequence ``laptop, photo, clip''.

QS: {[}extract{]} The sequence is ``1. laptop, 2. photo, 3. clip''. The
sequence is 3 items long, which is less than the minimum length of 4, so
we don't need to split it. All we need to do is reverse ``laptop, photo,
clip''.

A: ``laptop, photo, clip''

QS: \hyperref[cot]{cot} Reverse the sequence \#1.

A: ``clip, photo, laptop''

QS: {[}EOQ{]}

QC: Reverse the sequence ``newspaper, glasses, laptop, bottle''.

QS: {[}extract{]} The sequence is ``1. newspaper, 2. glasses, 3. laptop,
4. bottle''. The sequence is 4 items long, which is equal to the minimum
length of 4, so we split it. Half of 4 is 4 / 2 = 2.0. Dropping the
decimal, we get that the first half will be 2 items long. The first half
(2 items) of the sequence is ``1. newspaper, 2. glasses''.

A: ``1. newspaper, 2. glasses''

QS: {[}extract{]} The first half of the sequence ends with ``2.
glasses'', so the second half starts after ``2. glasses'' with ``3.
laptop''. The full sequence is 4 items long and the first half is 2
items long, so the second half will be 4 - 2 = 2 items long, ending in
``2. glasses''. The second half of the sequence (2 items) is ``3.
laptop, 4. bottle''.

A: ``3. laptop, 4. bottle''

QS: \hyperref[remove_numbers]{remove\_numbers} Remove the numbers from
\#1.

A: ``newspaper, glasses''

QS: \hyperref[remove_numbers]{remove\_numbers} Remove the numbers from
\#2.

A: ``laptop, bottle''

QS: \hyperref[reverse]{reverse} Reverse the sequence \#3.

A: ``glasses, newspaper''

QS: \hyperref[reverse]{reverse} Reverse the sequence \#4.

A: ``bottle, laptop''

QS: \hyperref[join]{join} \#6 \#5

A: ``bottle, laptop, glasses, newspaper''

QS: {[}EOQ{]}

\section{remove\_numbers}\label{remove_numbers}

Q: Remove the numbers from ``4. bottle, 3. laptop, 2. glasses, 1.
newspaper''.

A: ``bottle, laptop, glasses, newspaper''

Q: Remove the numbers from ``1. identity card, 2. packet, 3. button''.

A: ``identity card, packet, button''

Q: Remove the numbers from ``1. player, 2. passport, 3. umbrella, 4.
radio''.

A: ``player, passport, umbrella, radio''

\section{join}\label{join}

Q: ``bottle, laptop'' ``glasses, newspaper''.

A: ``bottle, laptop, glasses, newspaper''

Q: ``identity card, packet, button'' ``magazine, notebook, glasses''.

A: ``identity card, packet, button, magazine, notebook, glasses''

Q: ``passport, umbrella, radio, mobile phone, photo'' ``player''.

A: ``passport, umbrella, radio, mobile phone, photo, player''

Q: ``mirror, case'' ``toothbrush, alarm clock''.

A: ``mirror, case, toothbrush, alarm clock''

Q: ``light bulb, clip, umbrella'' ``driving licence, watch''.

A: ``light bulb, clip, umbrella, driving licence, watch''

\section{cot}\label{cot-1}

QC: Reverse the sequence ``newspaper, glasses, laptop, bottle''.

QS: First is newspaper. Second is glasses. Third is laptop. Fourth is
bottle. Now to reverse, change the order to: Fourth is bottle. Third is
laptop. Second is glasses. First is newspaper. So the answer is
``bottle, laptop, glasses, newspaper''.

A: ``bottle, laptop, glasses, newspaper''

QC: Reverse the sequence ``laptop, photo, clip''.

QS: First is laptop. Second is photo. Third is clip. Now to reverse,
change the order to: Third is clip. Second is photo. First is laptop. So
the answer is ``clip, photo, laptop''.

A: ``clip, photo, laptop''

QC: Reverse the sequence ``driving license, button, packet, identity
card, pineapple''.

QS: First is driving license. Second is button. Third is packet. Fourth
is identity card. Fifth is pineapple. Now to reverse, change the order
to: Fifth is pineapple. Fourth is identity card. Third is packet. Second
is button. First is driving license. So the answer is ``pineapple,
identity card, packet, button, driving license''.

A: ``pineapple, identity card, packet, button, driving license''

\section{unrolled\_decomp}\label{unrolled_decomp}

QC: Reverse the sequence ``driving license, button, packet, identity
card, shoe''.

QS: The sequence is ``1. driving license, 2. button, 3. packet, 4.
identity card, 5. shoe''. The sequence is 5 items long, which is more
than the minimum length of 4, so we split it. Half of 5 is 5 / 2 = 2.5.
Dropping the decimal, we get that the first half will be 2 items long,
ending in ``2. button''. The first half (2 items) is ``1. driving
license, 2. button''. The first half of the sequence ends with ``2.
button'', so the second half starts after ``2. button'' with ``3.
packet''. The full sequence is 5 items long, and the first half is 2
items long, so the second half will be 5 - 2 = 3 items long. The second
half of the sequence (3 items) is ``3. packet, 4. identity card, 5.
shoe''. Removing the numbers from ``1. driving license, 2. button'', we
get ``driving license, button''. Removing the numbers from ``3. packet,
4. identity card, 5. shoe'', we get ``packet, identity card, shoe''.
Reversing the sequence ``driving license, button'', First is driving
license. Second is button. Now to reverse, change the order to: Second
is button. First is driving license. So the answer is ``button, driving
license''. Reversing the sequence ``packet, identity card, shoe'', First
is packet. Second is identity card. Third is shoe. Now to reverse,
change the order to: Third is shoe. Second is identity card. First is
packet. So the answer is ``shoe, identity card, packet''. Joining
``shoe, identity card, packet'' and ``button, driving license'', the
answer is ``shoe, identity card, packet, button, driving license''.

A: ``shoe, identity card, packet, button, driving license''

QS: {[}EOQ{]}

QC: Reverse the sequence ``laptop, photo, clip''.

QS: The sequence is ``1. laptop, 2. photo, 3. clip''. The sequence is 3
items long, which is less than the minimum length of 4, so we don't need
to split it. All we need to do is reverse ``laptop, photo, clip''. First
is laptop. Second is photo. Third is clip. Now to reverse, change the
order to : Third is clip. Second is photo. First is laptop. So the
answer is ``clip, photo, laptop''.

A: ``clip, photo, laptop''

QS: {[}EOQ{]}

QC: Reverse the sequence ``newspaper, glasses, laptop, bottle''.

QS: The sequence is ``1. newspaper, 2. glasses, 3. laptop, 4. bottle''.
The sequence is 4 items long, which is equal to the minimum length of 4,
so we split it. Half of 4 is 4 / 2 = 2.0. Dropping the decimal, we get
that the first half will be 2 items long, ending in ``2. glasses''. The
first half

(2 items) is ``1. newspaper, 2. glasses''. The first half of the
sequence ends with ``2. glasses'', so the second half starts after ``2.
glasses'' with ``3. laptop''. The full sequence is 4 items long, and the
first half is 2 items long, so the second half will be 4 - 2 = 2 items
long. The second half of the sequence (2 items) is ``3. laptop, 4.
bottle''. Removing the numbers from ``1. newspaper, 2. glasses'', we get
``newspaper, glasses''. Removing the numbers from ``3. laptop, 4.
bottle'', we get ``laptop, bottle''. Reversing the sequence ``newspaper,
glasses'', First is newspaper. Second is glasses. Now to reverse, change
the order to: Second is glasses. First is newspaper. So the answer is
``glasses, newspaper''. Reversing the sequence ``laptop, bottle'', First
is laptop. Second is bottle. Now to reverse, change the order to: Second
is bottle. First is laptop. So the answer is ``bottle, laptop''. Joining
``bottle, laptop'' and ``glasses, newspaper'', the answer is ``bottle,
laptop, glasses, newspaper''.

A: ``bottle, laptop, glasses, newspaper''

QS: {[}EOQ{]}

\section{reverse (tail)}\label{reverse-tail}

QC: Reverse the sequence ``driving license, button, packet, identity
card, shoe''.

Q1: {[}extract{]} The sequence is ``1. driving license, 2. button, 3.
packet, 4. identity card, 5. shoe''. The sequence is 5 items long, which
is more than the minimum length of 4, so we split it. We take the first
element in the sequence which is ``1. driving license''.

\#1: ``1. driving license''

Q2: {[}extract{]} The full sequence is 5 items long, so the remaining
sequence will be 5 - 1 = 4 items long. The tail of the sequence with 4
items is ``2. button, 3. packet, 4. identity card, 5. shoe''.

\#2: ``2. button, 3. packet, 4. identity card, 5. shoe''

Q3: \hyperref[remove_numbers]{remove\_numbers} Remove the numbers from
\#1.

\#3: ``driving license''

Q4: \hyperref[remove_numbers]{remove\_numbers} Remove the numbers from
\#2.

\#4: ``button, packet, identity card, shoe''

Q5: \hyperref[reverse]{reverse} Reverse the sequence \#4.

\#5: ``shoe, identity card, packet, button''

Q6: \hyperref[join]{join} \#5 \#3

\#6: ``shoe, identity card, packet, button, driving license''

Q7: {[}EOQ{]}

QC: Reverse the sequence ``laptop, photo, clip''.

Q1: {[}extract{]} The sequence is ``1. laptop, 2. photo, 3. clip''. The
sequence is 3 items long, which is less than the minimum length of 4, so
we don't need to split it. All we need to do is reverse ``laptop, photo,
clip''.

\#1: ``laptop, photo, clip''

Q2: \hyperref[cot]{cot} Reverse the sequence \#1.

\#2: ``clip, photo, laptop''

Q3: {[}EOQ{]}

QC: Reverse the sequence ``newspaper, glasses, laptop, bottle''.

Q1: {[}extract{]} The sequence is ``1. newspaper, 2. glasses, 3. laptop,
4. bottle''. The sequence is 4 items long, which is more than the
minimum length of 4, so we split it. We take the first element in the
sequence which is ``1. newspaper''.

\#1: ``1. newspaper''

Q2: {[}extract{]} The full sequence is 4 items long, so the remaining
sequence will be 4 - 1 = 3 items long. The tail of the sequence with 3
items is ``2. glasses, 3. laptop, 4. bottle''.

\#2: ``2. glasses, 3. laptop, 4. bottle''

Q3: \hyperref[remove_numbers]{remove\_numbers} Remove the numbers from
\#1.

\#3: ``newspaper''

Q4: \hyperref[remove_numbers]{remove\_numbers} Remove the numbers from
\#2.

\#4: ``glasses, laptop, bottle''

Q5: \hyperref[reverse]{reverse} Reverse the sequence \#4.

\#5: ``bottle, laptop, glasses''

Q6: \hyperref[join]{join} \#5 \#3

\#6: ``bottle, laptop, glasses, newspaper''

Q7: {[}EOQ{]}

\section{G.3 LONG-DOCUMENT QA}\label{g.3-long-document-qa}

We show one of the prompts used for CommaQA experiments here. The entire
set of prompts is provided as supplementary material.

\section{G.3.1 DECOMPOSED PROMPTING:
(COARSE)}\label{g.3.1-decomposed-prompting-coarse}

\section{decomp}\label{decomp-2}

What awards have movies produced by people born in 1910 won?

QS: (select) \hyperref[qa]{qa} Who were born in the year 1910?

A: {[}``Teeplemole'', ``Muntaril''{]}

QS: (project\_values\_flat\_unique) \hyperref[qa]{qa} For which movies
was \#1 the producer?

A: {[}``Featsaw'', ``Zalate'', ``Premercy''{]}

QS: (project\_values\_flat\_unique) \hyperref[qa]{qa} Which awards were
given to \#2?

A: {[}``Zorgion'', ``Chowwurst'', ``Hallowcock''{]}

QS: {[}EOQ{]}

QC: What movies have people from the country Stridery acted in?

QS: (select) \hyperref[qa]{qa} Who is from the country Stridery?

A: {[}``Gastrat''{]}

QS: (project\_values\_flat\_unique) \hyperref[qa]{qa} Which movies has
\#1 been an actor in?

A: {[}``Partnershipmaker'', ``Nilitude'', ``Warpstone''{]}

QS: {[}EOQ{]}

QC: What awards have the actors of the Erowid winning movies received?

QS: (select) \hyperref[qa]{qa} Which movies were given the Erowid award?

A: {[}``Dewbar'', ``Caudacite''{]}

QS: (project\_values\_flat\_unique) \hyperref[qa]{qa} Who are the actors
in the movie \#1?

A: {[}``Wetherality'', ``Lougerière'', ``Gigabut''{]}

QS: (project\_values\_flat\_unique) \hyperref[qa]{qa} Which awards were
given to \#2?

A: {[}``Aniconder'', ``Trifogation''{]}

QS: {[}EOQ{]}

QC: What awards did the movies directed by the Modiparity winners
receive?

QS: (select) \hyperref[qa]{qa} Who has been awarded the Modiparity
award?

A: {[}``Bioperatology''{]}

QS: (project\_values\_flat\_unique) \hyperref[qa]{qa} Which movies has
\#1 directed?

A: {[}``Pestok'', ``Vitrilateral''{]}

QS: (project\_values\_flat\_unique) \hyperref[qa]{qa} Which awards were
given to \#2?

A: {[}``Gutney'', ``Antipositive''{]}

QS: {[}EOQ{]}

QC: What awards have movies written by people born in 1935 won?

QS: (select) \hyperref[qa]{qa} Who were born in the year 1935?

A: {[}``Sclerocybin'', ``Zayage''{]}

QS: (project\_values\_flat\_unique) \hyperref[qa]{qa} What movies has
\#1 written?

A: {[}``Noenometer'', ``Tayenne'', ``Pneumodendron''{]}

QS: (project\_values\_flat\_unique) \hyperref[qa]{qa} Which awards were
given to \#2?

A: {[}``Brownbeard'', ``Goosehead'', ``Handt''{]}

QS: {[}EOQ{]}

QC: What movies have the directors from Legault directed?

QS: (select) \hyperref[qa]{qa} Who is from the country Legault?

A: {[}``Metatoun'', ``Sapien''{]}

QS: (project\_values\_flat\_unique) \hyperref[qa]{qa} What movies has
\#1 been the director of?

A: {[}``Coacheship'', ``Misapportionment''{]}

QS: {[}EOQ{]}

\section{qa}\label{qa}

movie: Premercy ; directed by: Muntaril. movie: Skirtsicine ; director:
Teeplemole. movie: Featsaw ; directed by: Monsterscar. movie: Zalate ;
director: Monsterscar. movie: Zalate ; awarded: Hallowcock. movie:
Featsaw ; awarded: Zorgion. movie: Premercy ; award: Chowwurst. movie:
Skirtsicine ; award: Hallowcock. award: Goatfly ; winner: Teeplemole.
person: Monsterscar ; award: Glodome. person: Muntaril ; award: Goatfly.
movie: Featsaw ; release year: 1973. movie: Zalate ; release year: 1964.
movie: Skirtsicine ; release year: 1973. movie: Premercy ; year: 1961.
Teeplemole was an actor in the movie Skirtsicine. Muntaril was an actor
in the movie Skirtsicine. Monsterscar was an actor in the movie
Premercy. Muntaril was an actor in the movie Featsaw. Teeplemole was an
actor in the movie Zalate. Muntaril was born in the year 1910.
Teeplemole was born in 1910. Monsterscar was born in 1942. Teeplemole is
from the country of Piperfish. Monsterscar is from the country of
Piperfish. Muntaril is from the country of Clony. Muntaril produced the
movie Skirtsicine with others. Monsterscar was one of the producers of
the movie Featsaw. Monsterscar produced the movie Premercy with others.
Monsterscar produced the movie Zalate with others. Teeplemole was one of
the producers of the movie Featsaw. Teeplemole produced the movie Zalate
with others. Muntaril produced the movie Premercy with others.
Monsterscar wrote for the movie Premercy. Muntaril was one of the
writers for the movie Zalate. Muntaril wrote for the movie Featsaw.
Teeplemole wrote for the movie Featsaw. Monsterscar was one of the
writers for the movie Zalate. Teeplemole was one of the writers for the
movie Skirtsicine.

Q: For which movies was Teeplemole the producer?

A: {[}``Featsaw'', ``Zalate''{]}

Q: Which awards were given to Featsaw?

A: {[}``Zorgion''{]}

movie: Misgendery ; directed by: Wetherality. movie: Dewbar ; director:
Gigabut. movie: Caudacite ; director: Lougerière. movie: Tayenne ;
directed by: Lougerière. movie: Misgendery ; awarded: Microsouenesis.
movie: Dewbar ; awarded: Erowid. movie: Tayenne ; awarded: Cockspit.
movie: Caudacite ; award: Erowid. award: Aniconder ; winner:
Wetherality. award: Aniconder ; winner: Lougerière. person: Gigabut ;
award: Trifogation. movie: Dewbar ; release year: 1991. movie: Tayenne ;
year: 2013. movie: Caudacite ; release year: 2008. movie: Misgendery ;
year: 1991. Wetherality was an actor in the movie Dewbar. Gigabut was an
actor in the movie Tayenne. Lougerière was an actor in the movie
Tayenne. Lougerière acted in the movie Caudacite. Lougerière acted in
the movie Misgendery. Gigabut was an actor in the movie Caudacite.
Wetherality was an actor in the movie Misgendery. Wetherality was born
in the year 1917. Lougerière was born in 1926. Gigabut was born in the
year 1917. Gigabut grew up in the nation of Triclops. Lougerière is from
the country of Tatkin . Wetherality grew up in the nation of Tatkin.
Lougerière produced the movie Dewbar with others. Gigabut produced the
movie Tayenne with others. Gigabut produced the movie Dewbar with
others. Lougerière was one of the producers of the movie Misgendery.
Wetherality was one of the producers of the movie Caudacite. Gigabut was
one of the producers of the movie Caudacite. Wetherality produced the
movie Misgendery with others. Wetherality produced the movie Tayenne
with others. Wetherality wrote for the movie Tayenne. Gigabut wrote for
the movie Misgendery. Lougerière was one of the writers for the movie
Caudacite. Wetherality wrote for the movie Misgendery. Gigabut wrote for
the movie Tayenne. Gigabut wrote for the movie Dewbar. Lougerière wrote
for the movie Dewbar. Wetherality wrote for the movie Caudacite.

Q: Who are the actors in the movie Caudacite?

A: {[}``Lougerière,''Gigabut''{]}

Q: Which movies were given the Erowid award?

A: {[}``Dewbar'', ``Caudacite''{]}

movie: Pastillobox ; directed by: Firmline. movie: Clenestration ;
directed by: Carblock. movie: Pestok ; directed by: Bioperatology.
movie: Vitrilateral ; director: Bioperatology. movie: Vitrilateral ;
award: Antipositive. movie: Clenestration ; awarded: Handt. movie:
Pastillobox ; awarded: Handt. movie: Pestok ; awarded: Gutney. movie:
Pestok ; writer: Firmline. movie:

Clenestration ; written by: Carblock. movie: Pastillobox ; written by:
Bioperatology. movie: Pestok ; writer: Bioperatology. movie:
Clenestration ; written by: Firmline. movie: Vitrilateral ; writer:
Bioperatology. movie: Pastillobox ; writer: Carblock. movie:
Vitrilateral ; written by : Carblock. movie: Pestok ; release year:
1986. movie: Clenestration ; year: 1986. movie: Vitrilateral ; year:
1999. movie: Pastillobox ; release year: 1984. Carblock was an actor in
the movie Pastillobox. Firmline was an actor in the movie Vitrilateral.
Bioperatology was an actor in the movie Clenestration. Firmline acted in
the movie Pastillobox. Carblock was an actor in the movie Clenestration.
Bioperatology was an actor in the movie Pestok. Firmline was born in the
year 1904. Bioperatology was born in the year 1935. Carblock was born in
1935. Carblock grew up in the nation of Knoppock. Firmline grew up in
the nation of Tatkin. Bioperatology grew up in the nation of Tatkin.
Bioperatology won the Modiparity award. Halfbill was awarded to
Firmline. Halfbill was awarded to Carblock. Bioperatology was one of the
producers of the movie Pestok. Bioperatology produced the movie
Vitrilateral with others. Firmline produced the movie Pastillobox with
others. Firmline produced the movie Clenestration with others. Carblock
was one of the producers of the movie Pastillobox. Carblock produced the
movie Vitrilateral with others. Carblock produced the movie
Clenestration with others. Firmline was one of the producers of the
movie Pestok.

Q: Who has been awarded the Modiparity award?

A: {[}``Bioperatology''{]}

Q: Which movies has Bioperatology directed?

A: {[}``Pestok'', ``Vitrilateral''{]}

movie: Nohit ; director: Mimicocycle. movie: Noenometer ; director:
Mimicocycle. movie: Tayenne ; directed by: Zayage. movie: Pneumodendron
; director: Sclerocybin. movie: Tayenne ; awarded: Goosehead. movie:
Nohit ; awarded: Handt. movie: Pneumodendron ; award: Handt. movie:
Noenometer ; awarded: Brownbeard. movie: Nohit ; writer: Mimicocycle.
movie: Noenometer ; written by: Sclerocybin. movie: Tayenne ; writer:
Sclerocybin. movie: Pneumodendron ; written by: Zayage. movie: Tayenne ;
writer: Zayage. movie: Pneumodendron ; written by: Mimicocycle. movie:
Noenometer ; release year: 1991. movie: Tayenne ; year: 2013. movie:
Nohit ; year: 2005. movie: Pneumodendron ; year: 2005. Mimicocycle was
an actor in the movie Tayenne. Zayage acted in the movie Pneumodendron.
Zayage was an actor in the movie Nohit. Sclerocybin was an actor in the
movie Nohit. Sclerocybin was an actor in the movie Tayenne. Mimicocycle
was an actor in the movie Noenometer. Zayage was born in 1935.
Sclerocybin was born in 1935. Mimicocycle was born in 1930. Mimicocycle
is from the country of Calderita. Sclerocybin grew up in the nation of
Calderita. Zayage is from the country of Obility. Quinion was awarded to
Zayage. Fannyxist was awarded to Sclerocybin. Fannyxist was awarded to
Mimicocycle. Mimicocycle produced the movie Nohit with others. Zayage
was one of the producers of the movie Nohit. Sclerocybin was one of the
producers of the movie Tayenne. Sclerocybin produced the movie
Pneumodendron with others. Zayage produced the movie Pneumodendron with
others. Mimicocycle was one of the producers of the movie Tayenne.
Sclerocybin was one of the producers of the movie Noenometer. Zayage
produced the movie Noenometer with others

Q: What movies has Sclerocybin written?

A: {[}``Noenometer'', ``Tayenne''{]}

Q: Who were born in the year 1935?

A: {[}``Sclerocybin'', ``Zayage''{]}

\section{G.3.2 DECOMPOSED PROMPTING:
(FINE)}\label{g.3.2-decomposed-prompting-fine}

\section{decomp}\label{decomp-3}

What awards have movies produced by people born in 1910 won?

QS: (select) \hyperref[simp_qa]{simp\_qa} Who were born in the year
1910?

A: {[}``Teeplemole'', ``Muntaril''{]}

QS: (project\_values\_flat\_unique) \hyperref[pos_qa]{pos\_qa} For which
movies was \#1 the producer?

A: {[}``Featsaw'', ``Zalate'', ``Premercy''{]}

QS: (project\_values\_flat\_unique) {[}aw\_qa{]} Which awards were given
to \#2?

A: {[}``Zorgion'', ``Chowwurst'', ``Hallowcock''{]}

QS: {[}EOQ{]}

QC: What movies have people from the country Stridery acted in?

QS: (select) \hyperref[simp_qa]{simp\_qa} Who is from the country
Stridery?

A: {[}``Gastrat''{]}

QS: (project\_values\_flat\_unique) \hyperref[pos_qa]{pos\_qa} Which
movies has \#1 been an actor in?

A: {[}``Partnershipmaker'', ``Nilitude'', ``Warpstone''{]}

QS: {[}EOQ{]}

QC: What awards have the actors of the Erowid winning movies received?

QS: (select) {[}aw\_qa{]} Which movies were given the Erowid award?

A: {[}``Dewbar'', ``Caudacite''{]}

QS: (project\_values\_flat\_unique) \hyperref[pos_qa]{pos\_qa} Who are
the actors in the movie \#1?

A: {[}``Wetherality'', ``Lougerière'', ``Gigabut''{]}

QS: (project\_values\_flat\_unique) {[}aw\_qa{]} Which awards were given
to \#2?

A: {[}``Aniconder'', ``Trifogation''{]}

QS: {[}EOQ{]}

QC: What awards did the movies directed by the Modiparity winners
receive?

QS: (select) {[}aw\_qa{]} Who has been awarded the Modiparity award?

A: {[}``Bioperatology''{]}

QS: (project\_values\_flat\_unique) \hyperref[pos_qa]{pos\_qa} Which
movies has \#1 directed?

A: {[}``Pestok'', ``Vitrilateral''{]}

QS: (project\_values\_flat\_unique) {[}aw\_qa{]} Which awards were given
to \#2?

A: {[}``Gutney'', ``Antipositive''{]}

QS: {[}EOQ{]}

QC: What awards have movies written by people born in 1935 won?

QS: (select) \hyperref[simp_qa]{simp\_qa} Who were born in the year
1935?

A: {[}``Sclerocybin'', ``Zayage''{]}

QS: (project\_values\_flat\_unique) \hyperref[pos_qa]{pos\_qa} What
movies has \#1 written?

A: {[}``Noenometer'', ``Tayenne'', ``Pneumodendron''{]}

QS: (project\_values\_flat\_unique) {[}aw\_qa{]} Which awards were given
to \#2?

A: {[}``Brownbeard'', ``Goosehead'', ``Handt''{]}

QS: {[}EOQ{]}

QC: What movies have the directors from Legault directed?

QS: (select) \hyperref[simp_qa]{simp\_qa} Who is from the country
Legault?

A: {[}``Metatoun'', ``Sapien''{]}

QS: (project\_values\_flat\_unique) \hyperref[pos_qa]{pos\_qa} What
movies has \#1 been the director of?

A: {[}``Coacheship'', ``Misapportionment''{]}

QS: {[}EOQ{]}

aw\_qa

movie: Premercy ; directed by: Muntaril. movie: Skirtsicine ; director:
Teeplemole. movie: Featsaw

; directed by: Monsterscar. movie: Zalate ; director: Monsterscar.
movie: Zalate ; awarded: Hallowcock. movie: Featsaw ; awarded: Zorgion.
movie: Premercy ; award: Chowwurst. movie: Skirtsicine ; award:
Hallowcock. award: Goatfly ; winner: Teeplemole. person: Monsterscar ;
award: Glodome. person: Muntaril ; award: Goatfly. movie: Featsaw ;
release year: 1973. movie: Zalate ; release year: 1964. movie:
Skirtsicine ; release year: 1973. movie: Premercy ; year: 1961.
Teeplemole was an actor in the movie Skirtsicine. Muntaril was an actor
in the movie Skirtsicine. Monsterscar was an actor in the movie
Premercy. Muntaril was an actor in the movie Featsaw. Teeplemole was an
actor in the movie Zalate. Muntaril was born in the year 1910.
Teeplemole was born in 1910. Monsterscar was born in 1942. Teeplemole is
from the country of Piperfish. Monsterscar is from the country of
Piperfish. Muntaril is from the country of Clony. Muntaril produced the
movie Skirtsicine with others. Monsterscar was one of the producers of
the movie Featsaw. Monsterscar produced the movie Premercy with others.
Monsterscar produced the movie Zalate with others. Teeplemole was one of
the producers of the movie Featsaw. Teeplemole produced the movie Zalate
with

others. Muntaril produced the movie Premercy with others. Monsterscar
wrote for the movie Premercy. Muntaril was one of the writers for the
movie Zalate. Muntaril wrote for the movie Featsaw. Teeplemole wrote for
the movie Featsaw. Monsterscar was one of the writers for the movie
Zalate. Teeplemole was one of the writers for the movie Skirtsicine.

Q: Which awards were given to Zalate?

A: {[}``Hallowcock''{]}

Q: Which awards were given to Featsaw?

A: {[}``Zorgion''{]}

movie: Misgendery ; directed by: Wetherality. movie: Dewbar ; director:
Gigabut. movie: Caudacite ; director: Lougerière. movie: Tayenne ;
directed by: Lougerière. movie: Misgendery ; awarded: Microsouenesis.
movie: Dewbar ; awarded: Erowid. movie: Tayenne ; awarded: Cockspit.
movie: Caudacite ; award: Erowid. award: Aniconder ; winner:
Wetherality. award: Aniconder ; winner: Lougerière. person: Gigabut ;
award: Trifogation. movie: Dewbar ; release year: 1991. movie: Tayenne ;
year: 2013. movie: Caudacite ; release year: 2008. movie: Misgendery ;
year: 1991. Wetherality was an actor in the movie Dewbar. Gigabut was an
actor in the movie Tayenne. Lougerière was an actor in the movie
Tayenne. Lougerière acted in the movie Caudacite. Lougerière acted in
the movie Misgendery. Gigabut was an actor in the movie Caudacite.
Wetherality was an actor in the movie Misgendery. Wetherality was born
in the year 1917. Lougerière was born in 1926. Gigabut was born in the
year 1917. Gigabut grew up in the nation of Triclops. Lougerière is from
the country of Tatkin . Wetherality grew up in the nation of Tatkin.
Lougerière produced the movie Dewbar with others. Gigabut produced the
movie Tayenne with others. Gigabut produced the movie Dewbar with
others. Lougerière was one of the producers of the movie Misgendery.
Wetherality was one of the producers of the movie Caudacite. Gigabut was
one of the producers of the movie Caudacite. Wetherality produced the
movie Misgendery with others. Wetherality produced the movie Tayenne
with others. Wetherality wrote for the movie Tayenne. Gigabut wrote for
the movie Misgendery. Lougerière was one of the writers for the movie
Caudacite. Wetherality wrote for the movie Misgendery. Gigabut wrote for
the movie Tayenne. Gigabut wrote for the movie Dewbar. Lougerière wrote
for the movie Dewbar. Wetherality wrote for the movie Caudacite.

Q: Which movies were given the Erowid award?

A: {[}``Dewbar'', ``Caudacite''{]}

Q: Which awards were given to Wetherality?

A: {[}``Aniconder''{]}

movie: Pastillobox ; directed by: Firmline. movie: Clenestration ;
directed by: Carblock. movie: Pestok ; directed by: Bioperatology.
movie: Vitrilateral ; director: Bioperatology. movie: Vitrilateral ;
award: Antipositive. movie: Clenestration ; awarded: Handt. movie:
Pastillobox ; awarded: Handt. movie: Pestok ; awarded: Gutney. movie:
Pestok ; writer: Firmline. movie: Clenestration ; written by: Carblock.
movie: Pastillobox ; written by: Bioperatology. movie: Pestok ; writer:
Bioperatology. movie: Clenestration ; written by: Firmline. movie:
Vitrilateral ; writer: Bioperatology. movie: Pastillobox ; writer:
Carblock. movie: Vitrilateral ; written by : Carblock. movie: Pestok ;
release year: 1986. movie: Clenestration ; year: 1986. movie:
Vitrilateral ; year: 1999. movie: Pastillobox ; release year: 1984.
Carblock was an actor in the movie Pastillobox. Firmline was an actor in
the movie Vitrilateral. Bioperatology was an actor in the movie
Clenestration. Firmline acted in the movie Pastillobox. Carblock was an
actor in the movie Clenestration. Bioperatology was an actor in the
movie Pestok. Firmline was born in the year 1904. Bioperatology was born
in the year 1935. Carblock was born in 1935. Carblock grew up in the
nation of Knoppock. Firmline grew up in the nation of Tatkin.
Bioperatology grew up in the nation of Tatkin. Bioperatology won the
Modiparity award. Halfbill was awarded to Firmline. Halfbill was awarded
to Carblock. Bioperatology was one of the producers of the movie Pestok.
Bioperatology produced the movie Vitrilateral with others. Firmline
produced the movie Pastillobox with others. Firmline produced the movie
Clenestration with others. Carblock was one of the producers of the
movie Pastillobox. Carblock produced the movie Vitrilateral with others.
Carblock produced the movie Clenestration with others. Firmline was one
of the producers of the movie Pestok.

Q: Who has been awarded the Modiparity award?

A: {[}``Bioperatology''{]}

Q: Which awards were given to Pestok?

A: {[}``Gutney''{]}

movie: Nohit ; director: Mimicocycle. movie: Noenometer ; director:
Mimicocycle. movie: Tayenne ; directed by: Zayage. movie: Pneumodendron
; director: Sclerocybin. movie: Tayenne ; awarded: Goosehead. movie:
Nohit ; awarded: Handt. movie: Pneumodendron ; award: Handt. movie:
Noenometer ; awarded: Brownbeard. movie: Nohit ; writer: Mimicocycle.
movie: Noenometer ; written by: Sclerocybin. movie: Tayenne ; writer:
Sclerocybin. movie: Pneumodendron ; written by: Zayage. movie: Tayenne ;
writer: Zayage. movie: Pneumodendron ; written by: Mimicocycle. movie:
Noenometer ; release year: 1991. movie: Tayenne ; year: 2013. movie:
Nohit ; year: 2005. movie: Pneumodendron ; year: 2005. Mimicocycle was
an actor in the movie Tayenne. Zayage acted in the movie Pneumodendron.
Zayage was an actor in the movie Nohit. Sclerocybin was an actor in the
movie Nohit. Sclerocybin was an actor in the movie Tayenne. Mimicocycle
was an actor in the movie Noenometer. Zayage was born in 1935.
Sclerocybin was born in 1935. Mimicocycle was born in 1930. Mimicocycle
is from the country of Calderita. Sclerocybin grew up in the nation of
Calderita. Zayage is from the country of Obility. Quinion was awarded to
Zayage. Fannyxist was awarded to Sclerocybin. Fannyxist was awarded to
Mimicocycle. Mimicocycle produced the movie Nohit with others. Zayage
was one of the producers of the movie Nohit. Sclerocybin was one of the
producers of the movie Tayenne. Sclerocybin produced the movie
Pneumodendron with others. Zayage produced the movie Pneumodendron with
others. Mimicocycle was one of the producers of the movie Tayenne.
Sclerocybin was one of the producers of the movie Noenometer. Zayage
produced the movie Noenometer with others

Q: Which awards were given to Noenometer?

A: {[}``Brownbeard''{]}

Q: Which awards were given to Pneumodendron?

A: {[}''Handt''{]}

\section{pos\_qa}\label{pos_qa}

movie: Premercy ; directed by: Muntaril. movie: Skirtsicine ; director:
Teeplemole. movie: Featsaw ; directed by: Monsterscar. movie: Zalate ;
director: Monsterscar. movie: Zalate ; awarded: Hallowcock. movie:
Featsaw ; awarded: Zorgion. movie: Premercy ; award: Chowwurst. movie:
Skirtsicine ; award: Hallowcock. award: Goatfly ; winner: Teeplemole.
person: Monsterscar ; award: Glodome. person: Muntaril ; award: Goatfly.
movie: Featsaw ; release year: 1973. movie: Zalate ; release year: 1964.
movie: Skirtsicine ; release year: 1973. movie: Premercy ; year: 1961.
Teeplemole was an actor in the movie Skirtsicine. Muntaril was an actor
in the movie Skirtsicine. Monsterscar was an actor in the movie
Premercy. Muntaril was an actor in the movie Featsaw. Teeplemole was an
actor in the movie Zalate. Muntaril was born in the year 1910.
Teeplemole was born in 1910. Monsterscar was born in 1942. Teeplemole is
from the country of Piperfish. Monsterscar is from the country of
Piperfish. Muntaril is from the country of Clony. Muntaril produced the
movie Skirtsicine with others. Monsterscar was one of the producers of
the movie Featsaw. Monsterscar produced the movie Premercy with others.
Monsterscar produced the movie Zalate with others. Teeplemole was one of
the producers of the movie Featsaw. Teeplemole produced the movie Zalate
with others. Muntaril produced the movie Premercy with others.
Monsterscar wrote for the movie Premercy. Muntaril was one of the
writers for the movie Zalate. Muntaril wrote for the movie Featsaw.
Teeplemole wrote for the movie Featsaw. Monsterscar was one of the
writers for the movie Zalate. Teeplemole was one of the writers for the
movie Skirtsicine.

Q: For which movies was Teeplemole the producer?

A: {[}``Featsaw'', ``Zalate''{]}

Q: For which movies was Muntaril the producer?

A: {[}''Premercy{]}

movie: Misgendery ; directed by: Wetherality. movie: Dewbar ; director:
Gigabut. movie: Caudacite ; director: Lougerière. movie: Tayenne ;
directed by: Lougerière. movie: Misgendery ; awarded: Microsouenesis.
movie: Dewbar ; awarded: Erowid. movie: Tayenne ; awarded: Cockspit.
movie: Caudacite ; award: Erowid. award: Aniconder ; winner:
Wetherality. award: Aniconder ; winner: Lougerière. person: Gigabut ;
award: Trifogation. movie: Dewbar ; release year: 1991. movie: Tayenne ;
year: 2013. movie: Caudacite ; release year: 2008. movie: Misgendery ;
year: 1991. Wetherality was an actor in the movie Dewbar. Gigabut was an
actor in the movie Tayenne. Lougerière was an actor in the movie
Tayenne. Lougerière acted in the movie Caudacite. Lougerière acted in
the movie Misgendery. Gigabut was an actor in the movie Caudacite.
Wetherality was an actor in the movie Misgendery. Wetherality was born
in the year 1917. Lougerière was born in 1926. Gigabut was born in the
year 1917. Gigabut grew up in the nation of Triclops. Lougerière is from
the country of Tatkin . Wetherality grew up in the nation of Tatkin.
Lougerière produced the movie Dewbar with others. Gigabut produced the
movie Tayenne with others. Gigabut produced the movie Dewbar with
others. Lougerière was one of the producers of the movie Misgendery.
Wetherality was one of the producers of the movie Caudacite. Gigabut was
one of the producers of the movie Caudacite. Wetherality produced the
movie Misgendery with others. Wetherality produced the movie Tayenne
with others. Wetherality wrote for the movie Tayenne. Gigabut wrote for
the movie Misgendery. Lougerière was one of the writers for the movie
Caudacite. Wetherality wrote for the movie Misgendery. Gigabut wrote for
the movie Tayenne. Gigabut wrote for the movie Dewbar. Lougerière wrote
for the movie Dewbar. Wetherality wrote for the movie Caudacite.

Q: Who are the actors in the movie Dewbar?

A: {[}``Wetherality''{]}

Q: Who are the actors in the movie Caudacite?

A: {[}``Lougerière,''Gigabut''{]}

movie: Pastillobox ; directed by: Firmline. movie: Clenestration ;
directed by: Carblock. movie: Pestok ; directed by: Bioperatology.
movie: Vitrilateral ; director: Bioperatology. movie: Vitrilateral ;
award: Antipositive. movie: Clenestration ; awarded: Handt. movie:
Pastillobox ; awarded: Handt. movie: Pestok ; awarded: Gutney. movie:
Pestok ; writer: Firmline. movie: Clenestration ; written by: Carblock.
movie: Pastillobox ; written by: Bioperatology. movie: Pestok ; writer:
Bioperatology. movie: Clenestration ; written by: Firmline. movie:
Vitrilateral ; writer: Bioperatology. movie: Pastillobox ; writer:
Carblock. movie: Vitrilateral ; written by : Carblock. movie: Pestok ;
release year: 1986. movie: Clenestration ; year: 1986. movie:
Vitrilateral ; year: 1999. movie: Pastillobox ; release year: 1984.
Carblock was an actor in the movie Pastillobox. Firmline was an actor in
the movie Vitrilateral. Bioperatology was an actor in the movie
Clenestration. Firmline acted in the movie Pastillobox. Carblock was an
actor in the movie Clenestration. Bioperatology was an actor in the
movie Pestok. Firmline was born in the year 1904. Bioperatology was born
in the year 1935. Carblock was born in 1935. Carblock grew up in the
nation of Knoppock. Firmline grew up in the nation of Tatkin.
Bioperatology grew up in the nation of Tatkin. Bioperatology won the
Modiparity award. Halfbill was awarded to Firmline. Halfbill was awarded
to Carblock. Bioperatology was one of the producers of the movie Pestok.
Bioperatology produced the movie Vitrilateral with others. Firmline
produced the movie Pastillobox with others. Firmline produced the movie
Clenestration with others. Carblock was one of the producers of the
movie Pastillobox. Carblock produced the movie Vitrilateral with others.
Carblock produced the movie Clenestration with others. Firmline was one
of the producers of the movie Pestok.

Q: Which movies has Bioperatology directed?

A: {[}``Pestok'', ``Vitrilateral''{]}

Q: Which movies has Carblock directed?

A: {[}``Clenestration''{]}

movie: Nohit ; director: Mimicocycle. movie: Noenometer ; director:
Mimicocycle. movie: Tayenne ; directed by: Zayage. movie: Pneumodendron
; director: Sclerocybin. movie: Tayenne ; awarded: Goosehead. movie:
Nohit ; awarded: Handt. movie: Pneumodendron ; award: Handt. movie:
Noenometer ; awarded: Brownbeard. movie: Nohit ; writer:

Mimicocycle. movie: Noenometer ; written by: Sclerocybin. movie: Tayenne
; writer: Sclerocybin. movie: Pneumodendron ; written by: Zayage. movie:
Tayenne ; writer: Zayage. movie: Pneumodendron ; written by:
Mimicocycle. movie: Noenometer ; release year: 1991. movie: Tayenne ;
year: 2013. movie: Nohit ; year: 2005. movie: Pneumodendron ; year:
2005. Mimicocycle was an actor in the movie Tayenne. Zayage acted in the
movie Pneumodendron. Zayage was an actor in the movie Nohit. Sclerocybin
was an actor in the movie Nohit. Sclerocybin was an actor in the movie
Tayenne. Mimicocycle was an actor in the movie Noenometer. Zayage was
born in 1935. Sclerocybin was born in 1935. Mimicocycle was born in
1930. Mimicocycle is from the country of Calderita. Sclerocybin grew up
in the nation of Calderita. Zayage is from the country of Obility.
Quinion was awarded to Zayage. Fannyxist was awarded to Sclerocybin.
Fannyxist was awarded to Mimicocycle. Mimicocycle produced the movie
Nohit with others. Zayage was one of the producers of the movie Nohit.
Sclerocybin was one of the producers of the movie Tayenne. Sclerocybin
produced the movie Pneumodendron with others. Zayage produced the movie
Pneumodendron with others. Mimicocycle was one of the producers of the
movie Tayenne. Sclerocybin was one of the producers of the movie
Noenometer. Zayage produced the movie Noenometer with others

Q: What movies has Sclerocybin written?

A: {[}``Noenometer'', ``Tayenne''{]}

Q: What movies has Zayage written?

A: {[}``Pneumodendron'', ``Tayenne''{]}

\section{simp\_qa}\label{simp_qa}

movie: Premercy ; directed by: Muntaril. movie: Skirtsicine ; director:
Teeplemole. movie: Featsaw ; directed by: Monsterscar. movie: Zalate ;
director: Monsterscar. movie: Zalate ; awarded: Hallowcock. movie:
Featsaw ; awarded: Zorgion. movie: Premercy ; award: Chowwurst. movie:
Skirtsicine ; award: Hallowcock. award: Goatfly ; winner: Teeplemole.
person: Monsterscar ; award: Glodome. person: Muntaril ; award: Goatfly.
movie: Featsaw ; release year: 1973. movie: Zalate ; release year: 1964.
movie: Skirtsicine ; release year: 1973. movie: Premercy ; year: 1961.
Teeplemole was an actor in the movie Skirtsicine. Muntaril was an actor
in the movie Skirtsicine. Monsterscar was an actor in the movie
Premercy. Muntaril was an actor in the movie Featsaw. Teeplemole was an
actor in the movie Zalate. Muntaril was born in the year 1910.
Teeplemole was born in 1910. Monsterscar was born in 1942. Teeplemole is
from the country of Piperfish. Monsterscar is from the country of
Piperfish. Muntaril is from the country of Clony. Muntaril produced the
movie Skirtsicine with others. Monsterscar was one of the producers of
the movie Featsaw. Monsterscar produced the movie Premercy with others.
Monsterscar produced the movie Zalate with others. Teeplemole was one of
the producers of the movie Featsaw. Teeplemole produced the movie Zalate
with others. Muntaril produced the movie Premercy with others.
Monsterscar wrote for the movie Premercy. Muntaril was one of the
writers for the movie Zalate. Muntaril wrote for the movie Featsaw.
Teeplemole wrote for the movie Featsaw. Monsterscar was one of the
writers for the movie Zalate. Teeplemole was one of the writers for the
movie Skirtsicine.

Q: Who were born in the year 1910?

A: {[}``Teeplemole'', ``Muntaril''{]}

Q: From which country is Monsterscar?

A: {[}``Piperfish''{]}

movie: Nilitude ; director: Monsterscar. movie: Dewbar ; directed by:
Metatoun. movie: Warpstone ; directed by: Gastrat. movie:
Partnershipmaker ; director: Metatoun. movie: Dewbar ; award:
Tachychronograph. movie: Partnershipmaker ; awarded: Tachychronograph.
movie: Nilitude ; award: Paleodactyl. movie: Warpstone ; award:
Sabonade. person: Gastrat ; award: Trifogation. award: Polyquadrase ;
winner: Monsterscar. award: Trifogation ; winner: Metatoun. movie:
Warpstone ; release year: 1956. movie: Dewbar ; release year: 1984.
movie: Nilitude ; year: 1984. movie: Partnershipmaker ; year: 1962.
Gastrat was an actor in the movie Partnershipmaker. Metatoun was an
actor in the movie Partnershipmaker. Metatoun was an actor in the movie
Nilitude. Gastrat acted in the movie Nilitude. Monsterscar was an actor
in the movie Dewbar. Gastrat acted in the movie Warpstone. Metatoun
acted in the movie Warpstone. Metatoun was born in 1939. Gastrat was
born in the year 1933.

Monsterscar was born in 1933. Metatoun grew up in the nation of Moulole.
Gastrat is from the country of Stridery. Monsterscar grew up in the
nation of Moulole. Monsterscar produced the movie Nilitude with others.
Monsterscar was one of the producers of the movie Warpstone. Metatoun
was one of the producers of the movie Warpstone. Gastrat was one of the
producers of the movie Nilitude. Metatoun produced the movie
Partnershipmaker with others. Metatoun produced the movie Dewbar with
others. Monsterscar was one of the producers of the movie
Partnershipmaker. Gastrat produced the movie Dewbar with others.
Metatoun wrote for the movie Partnershipmaker. Gastrat wrote for the
movie Warpstone. Gastrat was one of the writers for the movie Dewbar.
Monsterscar was one of the writers for the movie Nilitude. Metatoun
wrote for the movie Warpstone.

Q: Who is from the country Stridery?

A: {[}``Gastrat''{]}

Q: Which movies were released in 1984?

A: {[}``Dewbar'', ``Nilitude''{]}

movie: Nohit ; director: Mimicocycle. movie: Noenometer ; director:
Mimicocycle. movie: Tayenne ; directed by: Zayage. movie: Pneumodendron
; director: Sclerocybin. movie: Tayenne ; awarded: Goosehead. movie:
Nohit ; awarded: Handt. movie: Pneumodendron ; award: Handt. movie:
Noenometer ; awarded: Brownbeard. movie: Nohit ; writer: Mimicocycle.
movie: Noenometer ; written by: Sclerocybin. movie: Tayenne ; writer:
Sclerocybin. movie: Pneumodendron ; written by: Zayage. movie: Tayenne ;
writer: Zayage. movie: Pneumodendron ; written by: Mimicocycle. movie:
Noenometer ; release year: 1991. movie: Tayenne ; year: 2013. movie:
Nohit ; year: 2005. movie: Pneumodendron ; year: 2005. Mimicocycle was
an actor in the movie Tayenne. Zayage acted in the movie Pneumodendron.
Zayage was an actor in the movie Nohit. Sclerocybin was an actor in the
movie Nohit. Sclerocybin was an actor in the movie Tayenne. Mimicocycle
was an actor in the movie Noenometer. Zayage was born in 1935.
Sclerocybin was born in 1935. Mimicocycle was born in 1930. Mimicocycle
is from the country of Calderita. Sclerocybin grew up in the nation of
Calderita. Zayage is from the country of Obility. Quinion was awarded to
Zayage. Fannyxist was awarded to Sclerocybin. Fannyxist was awarded to
Mimicocycle. Mimicocycle produced the movie Nohit with others. Zayage
was one of the producers of the movie Nohit. Sclerocybin was one of the
producers of the movie Tayenne. Sclerocybin produced the movie
Pneumodendron with others. Zayage produced the movie Pneumodendron with
others. Mimicocycle was one of the producers of the movie Tayenne.
Sclerocybin was one of the producers of the movie Noenometer. Zayage
produced the movie Noenometer with others

Q: Who were born in the year 1935?

A: {[}``Sclerocybin'', ``Zayage''{]}

Q: When was the movie Nohit released?

A: {[}''2005''{]}

movie: Coacheship ; director: Metatoun. movie: Assamplifier ; director:
Kapod. movie: Misapportionment ; director: Sapien. movie: Quinsid ;
director: Kapod. movie: Assamplifier ; award: Zorgion. movie: Quinsid ;
awarded: Airpipe. movie: Coacheship ; award: Electrodesal. movie:
Misapportionment ; award: Airpipe. movie: Coacheship ; written by:
Metatoun. movie: Misapportionment ; written by: Kapod. movie: Coacheship
; written by: Kapod. movie : Quinsid ; writer: Sapien. movie:
Misapportionment ; written by: Metatoun. movie: Assamplifier ; written
by: Kapod. movie: Assamplifier ; written by: Sapien. movie: Assamplifier
; release year: 2000. movie: Coacheship ; year: 2001. movie: Quinsid ;
year: 2005. movie: Misapportionment ; year: 2005. Sapien was an actor in
the movie Misapportionment. Sapien acted in the movie Assamplifier.
Kapod acted in the movie Quinsid . Sapien acted in the movie Coacheship.
Metatoun was an actor in the movie Quinsid. Kapod acted in the movie
Misapportionment. Metatoun acted in the movie Coacheship. Kapod acted in
the movie Assamplifier. Sapien was born in the year 1910. Metatoun was
born in 1928. Kapod was born in the year 1910. Metatoun is from the
country of Legault. Kapod grew up in the nation of Tatkin. Sapien is
from the country of Legault. Malwarp was awarded to Sapien. Metatoun won
the Monkeynote award. Kapod won the Monkeynote award. Kapod was one of
the producers of the movie Quinsid. Metatoun was one of the producers of
the movie

Misapportionment. Metatoun was one of the producers of the movie
Quinsid. Sapien was one of the producers of the movie Assamplifier.
Sapien produced the movie Coacheship with others. Metatoun was one of
the producers of the movie Assamplifier. Kapod was one of the producers
of the movie Misapportionment. Kapod was one of the producers of the
movie Coacheship.\\
Q: Who is from the country Legault?\\
A: {[}``Metatoun'', ``Sapien''{]}\\
Q: When was Kapod born?\\
A: {[}''1910''{]}

\section{G.3.3 COT}\label{g.3.3-cot}

\section{COT}\label{cot-2}

movie: Premercy ; directed by: Muntaril. movie: Skirtsicine ; director:
Teeplemole. movie: Featsaw ; directed by: Monsterscar. movie: Zalate ;
director: Monsterscar. movie: Zalate ; awarded: Hallowcock. movie:
Featsaw ; awarded: Zorgion. movie: Premercy ; award: Chowwurst. movie:
Skirtsicine ; award: Hallowcock. award: Goatfly ; winner: Teeplemole.
person: Monsterscar ; award: Glodome. person: Muntaril ; award: Goatfly.
movie: Featsaw ; release year: 1973. movie: Zalate ; release year: 1964.
movie: Skirtsicine ; release year: 1973. movie: Premercy ; year: 1961.
Teeplemole was an actor in the movie Skirtsicine. Muntaril was an actor
in the movie Skirtsicine. Monsterscar was an actor in the movie
Premercy. Muntaril was an actor in the movie Featsaw. Teeplemole was an
actor in the movie Zalate. Muntaril was born in the year 1910.
Teeplemole was born in 1910. Monsterscar was born in 1942. Teeplemole is
from the country of Piperfish. Monsterscar is from the country of
Piperfish. Muntaril is from the country of Clony. Muntaril produced the
movie Skirtsicine with others. Monsterscar was one of the producers of
the movie Featsaw. Monsterscar produced the movie Premercy with others.
Monsterscar produced the movie Zalate with others. Teeplemole was one of
the producers of the movie Featsaw. Teeplemole produced the movie Zalate
with others. Muntaril produced the movie Premercy with others.
Monsterscar wrote for the movie Premercy. Muntaril was one of the
writers for the movie Zalate. Muntaril wrote for the movie Featsaw.
Teeplemole wrote for the movie Featsaw. Monsterscar was one of the
writers for the movie Zalate. Teeplemole was one of the writers for the
movie Skirtsicine.

QC: What awards have movies produced by people born in 1910 won?

QS: The people born in 1910 were Muntaril and Teeplemole. Teeplemole
produced the movies Featsaw and Zalate. Muntaril produced the movie
Premercy. Featsaw was awarded the Zorgion award. Premercy was awarded
the Chowwurst award. Zalate was awarded the Hallowcock award. So the
answer is {[}``Zorgion'', ``Chowwurst'', ``Hallowcock''{]}.

A: {[}``Zorgion'', ``Chowwurst'', ``Hallowcock''{]}

QS: {[}EOQ{]}

movie: Misgendery ; directed by: Wetherality. movie: Dewbar ; director:
Gigabut. movie: Caudacite ; director: Lougerière. movie: Tayenne ;
directed by: Lougerière. movie: Misgendery ; awarded: Microsouenesis.
movie: Dewbar ; awarded: Erowid. movie: Tayenne ; awarded: Cockspit.
movie: Caudacite ; award: Erowid. award: Aniconder ; winner:
Wetherality. award: Aniconder ; winner: Lougerière. person: Gigabut ;
award: Trifogation. movie: Dewbar ; release year: 1991. movie: Tayenne ;
year: 2013. movie: Caudacite ; release year: 2008. movie: Misgendery ;
year: 1991. Wetherality was an actor in the movie Dewbar. Gigabut was an
actor in the movie Tayenne. Lougerière was an actor in the movie
Tayenne. Lougerière acted in the movie Caudacite. Lougerière acted in
the movie Misgendery. Gigabut was an actor in the movie Caudacite.
Wetherality was an actor in the movie Misgendery. Wetherality was born
in the year 1917. Lougerière was born in 1926. Gigabut was born in the
year 1917. Gigabut grew up in the nation of Triclops. Lougerière is from
the country of Tatkin . Wetherality grew up in the nation of Tatkin.
Lougerière produced the movie Dewbar with others. Gigabut produced the
movie Tayenne with others. Gigabut produced the movie Dewbar with
others. Lougerière was one of the producers of the movie Misgendery.
Wetherality was one of the producers of the movie Caudacite. Gigabut was
one of the producers of the movie Caudacite. Wetherality produced the
movie Misgendery with others. Wetherality produced the movie Tayenne
with others. Wetherality wrote for the movie

Tayenne. Gigabut wrote for the movie Misgendery. Lougerière was one of
the writers for the movie Caudacite. Wetherality wrote for the movie
Misgendery. Gigabut wrote for the movie Tayenne. Gigabut wrote for the
movie Dewbar. Lougerière wrote for the movie Dewbar. Wetherality wrote
for the movie Caudacite.

QC: What awards have the actors of the Erowid winning movies received?

QS: The movies that won the Erowid award are Dewbar and Caudacite.
Wetherality was an actor in Dewbar. Lougerière and Gigabut acted in
Caudacite. Wetherality won the Aniconder award. Lougerière also won the
Aniconder award. Gigabut won the Trifogation award. So the answer is
{[}``Aniconder'', ``Trifogation''{]}.

A: {[}``Aniconder'', ``Trifogation''{]}

QS: {[}EOQ{]}

movie: Pastillobox ; directed by: Firmline. movie: Clenestration ;
directed by: Carblock. movie: Pestok ; directed by: Bioperatology.
movie: Vitrilateral ; director: Bioperatology. movie: Vitrilateral ;
award: Antipositive. movie: Clenestration ; awarded: Handt. movie:
Pastillobox ; awarded: Handt. movie: Pestok ; awarded: Gutney. movie:
Pestok ; writer: Firmline. movie: Clenestration ; written by: Carblock.
movie: Pastillobox ; written by: Bioperatology. movie: Pestok ; writer:
Bioperatology. movie: Clenestration ; written by: Firmline. movie:
Vitrilateral ; writer: Bioperatology. movie: Pastillobox ; writer:
Carblock. movie: Vitrilateral ; written by : Carblock. movie: Pestok ;
release year: 1986. movie: Clenestration ; year: 1986. movie:
Vitrilateral ; year: 1999. movie: Pastillobox ; release year: 1984.
Carblock was an actor in the movie Pastillobox. Firmline was an actor in
the movie Vitrilateral. Bioperatology was an actor in the movie
Clenestration. Firmline acted in the movie Pastillobox. Carblock was an
actor in the movie Clenestration. Bioperatology was an actor in the
movie Pestok. Firmline was born in the year 1904. Bioperatology was born
in the year 1935. Carblock was born in 1935. Carblock grew up in the
nation of Knoppock. Firmline grew up in the nation of Tatkin.
Bioperatology grew up in the nation of Tatkin. Bioperatology won the
Modiparity award. Halfbill was awarded to Firmline. Halfbill was awarded
to Carblock. Bioperatology was one of the producers of the movie Pestok.
Bioperatology produced the movie Vitrilateral with others. Firmline
produced the movie Pastillobox with others. Firmline produced the movie
Clenestration with others. Carblock was one of the producers of the
movie Pastillobox. Carblock produced the movie Vitrilateral with others.
Carblock produced the movie Clenestration with others. Firmline was one
of the producers of the movie Pestok.

QC: What awards did the movies directed by the Modiparity winners
receive?

QS: The Modiparity winners are Bioperatology. Bioperatology has directed
the movies Pestok and Vitrilateral. Pestok has won the Gutney award.
Vitrilateral has won the Antipositive award. So the answer is
{[}``Gutney'', ``Antipositive''{]}.

A: {[}``Gutney'', ``Antipositive''{]}

QS: {[}EOQ{]}

movie: Nohit ; director: Mimicocycle. movie: Noenometer ; director:
Mimicocycle. movie: Tayenne ; directed by: Zayage. movie: Pneumodendron
; director: Sclerocybin. movie: Tayenne ; awarded: Goosehead. movie:
Nohit ; awarded: Handt. movie: Pneumodendron ; award: Handt. movie:
Noenometer ; awarded: Brownbeard. movie: Nohit ; writer: Mimicocycle.
movie: Noenometer ; written by: Sclerocybin. movie: Tayenne ; writer:
Sclerocybin. movie: Pneumodendron ; written by: Zayage. movie: Tayenne ;
writer: Zayage. movie: Pneumodendron ; written by: Mimicocycle. movie:
Noenometer ; release year: 1991. movie: Tayenne ; year: 2013. movie:
Nohit ; year: 2005. movie: Pneumodendron ; year: 2005. Mimicocycle was
an actor in the movie Tayenne. Zayage acted in the movie Pneumodendron.
Zayage was an actor in the movie Nohit. Sclerocybin was an actor in the
movie Nohit. Sclerocybin was an actor in the movie Tayenne. Mimicocycle
was an actor in the movie Noenometer. Zayage was born in 1935.
Sclerocybin was born in 1935. Mimicocycle was born in 1930. Mimicocycle
is from the country of Calderita. Sclerocybin grew up in the nation of
Calderita. Zayage is from the country of Obility. Quinion was awarded to
Zayage. Fannyxist was awarded to Sclerocybin. Fannyxist was awarded to
Mimicocycle. Mimicocycle produced the movie Nohit with others. Zayage
was one of the producers of the movie Nohit. Sclerocybin was one of the
producers of the movie Tayenne. Sclerocybin produced the movie
Pneumodendron with others. Zayage produced the movie Pneumodendron with
others.

Mimicocycle was one of the producers of the movie Tayenne. Sclerocybin
was one of the producers of the movie Noenometer. Zayage produced the
movie Noenometer with others\\
QC: What awards have movies written by people born in 1935 won?\\
QS: The people born in 1935 were Sclerocybin and Zayage. Sclerocybin has
written the movies Noenometer and Tayenne. Zayage has written the movies
Pneumodendron and Tayenne.\\
Noenometer has won the Brownbeard award. Tayenne has won the Goosehead
award.\\
Pneumodendron has won the Handt award. So the answer is
{[}``Brownbeard'', ``Goosehead'', ``Handt''{]}.\\
A: {[}``Brownbeard'', ``Goosehead'', ``Handt''{]}\\
QS: {[}EOQ{]}

\section{G.4 SHORTER PROMPTS FOR SMALLER CONTEXT
WINDOWS}\label{g.4-shorter-prompts-for-smaller-context-windows}

\section{qa (small)}\label{qa-small}

movie: Premercy ; directed by: Muntaril. movie: Skirtsicine ; director:
Teeplemole. movie: Featsaw ; directed by: Monsterscar. movie: Zalate ;
director: Monsterscar. movie: Zalate ; awarded: Hallowcock. movie:
Featsaw ; awarded: Zorgion. movie: Premercy ; award: Chowwurst. movie:
Skirtsicine ; award: Hallowcock. award: Goatfly ; winner: Teeplemole.
person: Monsterscar ; award: Glodome. person: Muntaril ; award: Goatfly.
movie: Featsaw ; release year: 1973. movie: Zalate ; release year: 1964.
movie: Skirtsicine ; release year: 1973. movie: Premercy ; year: 1961.
Teeplemole was an actor in the movie Skirtsicine. Muntaril was an actor
in the movie Skirtsicine. Monsterscar was an actor in the movie
Premercy. Muntaril was an actor in the movie Featsaw. Teeplemole was an
actor in the movie Zalate. Muntaril was born in the year 1910.
Teeplemole was born in 1910. Monsterscar was born in 1942. Teeplemole is
from the country of Piperfish. Monsterscar is from the country of
Piperfish. Muntaril is from the country of Clony. Muntaril produced the
movie Skirtsicine with others. Monsterscar was one of the producers of
the movie Featsaw. Monsterscar produced the movie Premercy with others.
Monsterscar produced the movie Zalate with others. Teeplemole was one of
the producers of the movie Featsaw. Teeplemole produced the movie Zalate
with others. Muntaril produced the movie Premercy with others.
Monsterscar wrote for the movie Premercy. Muntaril was one of the
writers for the movie Zalate. Muntaril wrote for the movie Featsaw.
Teeplemole wrote for the movie Featsaw. Monsterscar was one of the
writers for the movie Zalate. Teeplemole was one of the writers for the
movie Skirtsicine.

Q: Which awards were given to Zalate?

A: {[}``Hallowcock''{]}

Q: For which movies was Muntaril the producer?

A: {[}``Premercy{]}

Q: Who are the actors in the movie Premercy?

A: {[}``Monsterscar''{]}

movie: Nilitude ; director: Monsterscar. movie: Dewbar ; directed by:
Metatoun. movie: Warpstone ; directed by: Gastrat. movie:
Partnershipmaker ; director: Metatoun. movie: Dewbar ; award:
Tachychronograph. movie: Partnershipmaker ; awarded: Tachychronograph.
movie: Nilitude ; award: Paleodactyl. movie: Warpstone ; award:
Sabonade. person: Gastrat ; award: Trifogation. award: Polyquadrase ;
winner: Monsterscar. award: Trifogation ; winner: Metatoun. movie:
Warpstone ; release year: 1956. movie: Dewbar ; release year: 1984.
movie: Nilitude ; year: 1984. movie: Partnershipmaker ; year: 1962.
Gastrat was an actor in the movie Partnershipmaker. Metatoun was an
actor in the movie Partnershipmaker. Metatoun was an actor in the movie
Nilitude. Gastrat acted in the movie Nilitude. Monsterscar was an actor
in the movie Dewbar. Gastrat acted in the movie Warpstone. Metatoun
acted in the movie Warpstone. Metatoun was born in 1939. Gastrat was
born in the year 1933. Monsterscar was born in 1933. Metatoun grew up in
the nation of Moulole. Gastrat is from the country of Stridery.
Monsterscar grew up in the nation of Moulole. Monsterscar produced the
movie Nilitude with others. Monsterscar was one of the producers of the
movie Warpstone . Metatoun was one of the producers of the movie
Warpstone. Gastrat was one of the producers of the movie Nilitude.
Metatoun produced the movie Partnershipmaker with others. Metatoun
produced the movie Dewbar with others. Monsterscar was one of the
producers of

the movie Partnershipmaker. Gastrat produced the movie Dewbar with
others. Metatoun wrote for the movie Partnershipmaker. Gastrat wrote for
the movie Warpstone. Gastrat was one of the writers for the movie
Dewbar. Monsterscar was one of the writers for the movie Nilitude.
Metatoun wrote for the movie Warpstone.

Q: Which movies has Gastrat been an actor in?

A: {[}``Partnershipmaker'', ``Nilitude'', ``Warpstone''{]}

Q: Who is from the country Stridery?

A: {[}``Gastrat''{]}

Q: Which movies were given the Tachychronograph award?

A: {[}``Dewbar'', ``Partnershipmaker''{]}

\section{aw\_qa(small)}\label{aw_qasmall}

movie: Premercy ; directed by: Muntaril. movie: Skirtsicine ; director:
Teeplemole. movie: Featsaw ; directed by: Monsterscar. movie: Zalate ;
director: Monsterscar. movie: Zalate ; awarded: Hallowcock. movie:
Featsaw ; awarded: Zorgion. movie: Premercy ; award: Chowwurst. movie:
Skirtsicine ; award: Hallowcock. award: Goatfly ; winner: Teeplemole.
person: Monsterscar ; award: Glodome. person: Muntaril ; award: Goatfly.
movie: Featsaw ; release year: 1973. movie: Zalate ; release year: 1964.
movie: Skirtsicine ; release year: 1973. movie: Premercy ; year: 1961.
Teeplemole was an actor in the movie Skirtsicine. Muntaril was an actor
in the movie Skirtsicine. Monsterscar was an actor in the movie
Premercy. Muntaril was an actor in the movie Featsaw. Teeplemole was an
actor in the movie Zalate. Muntaril was born in the year 1910.
Teeplemole was born in 1910. Monsterscar was born in 1942. Teeplemole is
from the country of Piperfish. Monsterscar is from the country of
Piperfish. Muntaril is from the country of Clony. Muntaril produced the
movie Skirtsicine with others. Monsterscar was one of the producers of
the movie Featsaw. Monsterscar produced the movie Premercy with others.
Monsterscar produced the movie Zalate with others. Teeplemole was one of
the producers of the movie Featsaw. Teeplemole produced the movie Zalate
with others. Muntaril produced the movie Premercy with others.
Monsterscar wrote for the movie Premercy. Muntaril was one of the
writers for the movie Zalate. Muntaril wrote for the movie Featsaw.
Teeplemole wrote for the movie Featsaw. Monsterscar was one of the
writers for the movie Zalate. Teeplemole was one of the writers for the
movie Skirtsicine.

Q: Which awards were given to Zalate?

A: {[}``Hallowcock''{]}

Q: Which awards were given to Premercy?

A: {[}``Chowwurst''{]}

movie: Misgendery ; directed by: Wetherality. movie: Dewbar ; director:
Gigabut. movie: Caudacite ; director: Lougerière. movie: Tayenne ;
directed by: Lougerière. movie: Misgendery ; awarded: Microsouenesis.
movie: Dewbar ; awarded: Erowid. movie: Tayenne ; awarded: Cockspit.
movie: Caudacite ; award: Erowid. award: Aniconder ; winner:
Wetherality. award: Aniconder ; winner: Lougerière. person: Gigabut ;
award: Trifogation. movie: Dewbar ; release year: 1991. movie: Tayenne ;
year: 2013. movie: Caudacite ; release year: 2008. movie: Misgendery ;
year: 1991. Wetherality was an actor in the movie Dewbar. Gigabut was an
actor in the movie Tayenne. Lougerière was an actor in the movie
Tayenne. Lougerière acted in the movie Caudacite. Lougerière acted in
the movie Misgendery. Gigabut was an actor in the movie Caudacite.
Wetherality was an actor in the movie Misgendery. Wetherality was born
in the year 1917. Lougerière was born in 1926. Gigabut was born in the
year 1917. Gigabut grew up in the nation of Triclops. Lougerière is from
the country of Tatkin . Wetherality grew up in the nation of Tatkin.
Lougerière produced the movie Dewbar with others. Gigabut produced the
movie Tayenne with others. Gigabut produced the movie Dewbar with
others. Lougerière was one of the producers of the movie Misgendery.
Wetherality was one of the producers of the movie Caudacite. Gigabut was
one of the producers of the movie Caudacite. Wetherality produced the
movie Misgendery with others. Wetherality produced the movie Tayenne
with others. Wetherality wrote for the movie Tayenne. Gigabut wrote for
the movie Misgendery. Lougerière was one of the writers for the movie
Caudacite. Wetherality wrote for the movie Misgendery. Gigabut wrote for
the movie

Tayenne. Gigabut wrote for the movie Dewbar. Lougerière wrote for the
movie Dewbar. Wetherality wrote for the movie Caudacite.

Q: Which movies were given the Erowid award?

A: {[}``Dewbar'', ``Caudacite''{]}

Q: Which awards were given to Wetherality?

A: {[}``Aniconder''{]}

\section{pos\_qa(small)}\label{pos_qasmall}

movie: Premercy ; directed by: Muntaril. movie: Skirtsicine ; director:
Teeplemole. movie: Featsaw ; directed by: Monsterscar. movie: Zalate ;
director: Monsterscar. movie: Zalate ; awarded: Hallowcock. movie:
Featsaw ; awarded: Zorgion. movie: Premercy ; award: Chowwurst. movie:
Skirtsicine ; award: Hallowcock. award: Goatfly ; winner: Teeplemole.
person: Monsterscar ; award: Glodome. person: Muntaril ; award: Goatfly.
movie: Featsaw ; release year: 1973. movie: Zalate ; release year: 1964.
movie: Skirtsicine ; release year: 1973. movie: Premercy ; year: 1961.
Teeplemole was an actor in the movie Skirtsicine. Muntaril was an actor
in the movie Skirtsicine. Monsterscar was an actor in the movie
Premercy. Muntaril was an actor in the movie Featsaw. Teeplemole was an
actor in the movie Zalate. Muntaril was born in the year 1910.
Teeplemole was born in 1910. Monsterscar was born in 1942. Teeplemole is
from the country of Piperfish. Monsterscar is from the country of
Piperfish. Muntaril is from the country of Clony. Muntaril produced the
movie Skirtsicine with others. Monsterscar was one of the producers of
the movie Featsaw. Monsterscar produced the movie Premercy with others.
Monsterscar produced the movie Zalate with others. Teeplemole was one of
the producers of the movie Featsaw. Teeplemole produced the movie Zalate
with others. Muntaril produced the movie Premercy with others.
Monsterscar wrote for the movie Premercy. Muntaril was one of the
writers for the movie Zalate. Muntaril wrote for the movie Featsaw.
Teeplemole wrote for the movie Featsaw. Monsterscar was one of the
writers for the movie Zalate. Teeplemole was one of the writers for the
movie Skirtsicine.

Q: For which movies was Teeplemole the producer?

A: {[}``Featsaw'', ``Zalate''{]}

Q: For which movies was Muntaril the producer?

A: {[}``Premercy{]}

movie: Nilitude ; director: Monsterscar. movie: Dewbar ; directed by:
Metatoun. movie: Warpstone ; directed by: Gastrat. movie:
Partnershipmaker ; director: Metatoun. movie: Dewbar ; award:
Tachychronograph. movie: Partnershipmaker ; awarded: Tachychronograph.
movie: Nilitude ; award: Paleodactyl. movie: Warpstone ; award:
Sabonade. person: Gastrat ; award: Trifogation. award: Polyquadrase ;
winner: Monsterscar. award: Trifogation ; winner: Metatoun. movie:
Warpstone ; release year: 1956. movie: Dewbar ; release year: 1984.
movie: Nilitude ; year: 1984. movie: Partnershipmaker ; year: 1962.
Gastrat was an actor in the movie Partnershipmaker. Metatoun was an
actor in the movie Partnershipmaker. Metatoun was an actor in the movie
Nilitude. Gastrat acted in the movie Nilitude. Monsterscar was an actor
in the movie Dewbar. Gastrat acted in the movie Warpstone. Metatoun
acted in the movie Warpstone. Metatoun was born in 1939. Gastrat was
born in the year 1933. Monsterscar was born in 1933. Metatoun grew up in
the nation of Moulole. Gastrat is from the country of Stridery.
Monsterscar grew up in the nation of Moulole. Monsterscar produced the
movie Nilitude with others. Monsterscar was one of the producers of the
movie Warpstone . Metatoun was one of the producers of the movie
Warpstone. Gastrat was one of the producers of the movie Nilitude.
Metatoun produced the movie Partnershipmaker with others. Metatoun
produced the movie Dewbar with others. Monsterscar was one of the
producers of the movie Partnershipmaker. Gastrat produced the movie
Dewbar with others. Metatoun wrote for the movie Partnershipmaker.
Gastrat wrote for the movie Warpstone. Gastrat was one of the writers
for the movie Dewbar. Monsterscar was one of the writers for the movie
Nilitude. Metatoun wrote for the movie Warpstone.

Q: Which movies has Gastrat been an actor in?

A: {[}``Partnershipmaker'', ``Nilitude'', ``Warpstone''{]}

Q: Which movies has Monsterscar been an actor in?

A: {[}``Dewbar''{]}

\section{simp\_qa(small)}\label{simp_qasmall}

movie: Premercy ; directed by: Muntaril. movie: Skirtsicine ; director:
Teeplemole. movie: Featsaw ; directed by: Monsterscar. movie: Zalate ;
director: Monsterscar. movie: Zalate ; awarded: Hallowcock. movie:
Featsaw ; awarded: Zorgion. movie: Premercy ; award: Chowwurst. movie:
Skirtsicine ; award: Hallowcock. award: Goatfly ; winner: Teeplemole.
person: Monsterscar ; award: Glodome. person: Muntaril ; award: Goatfly.
movie: Featsaw ; release year: 1973. movie: Zalate ; release year: 1964.
movie: Skirtsicine ; release year: 1973. movie: Premercy ; year: 1961.
Teeplemole was an actor in the movie Skirtsicine. Muntaril was an actor
in the movie Skirtsicine. Monsterscar was an actor in the movie
Premercy. Muntaril was an actor in the movie Featsaw. Teeplemole was an
actor in the movie Zalate. Muntaril was born in the year 1910.
Teeplemole was born in 1910. Monsterscar was born in 1942. Teeplemole is
from the country of Piperfish. Monsterscar is from the country of
Piperfish. Muntaril is from the country of Clony. Muntaril produced the
movie Skirtsicine with others. Monsterscar was one of the producers of
the movie Featsaw. Monsterscar produced the movie Premercy with others.
Monsterscar produced the movie Zalate with others. Teeplemole was one of
the producers of the movie Featsaw. Teeplemole produced the movie Zalate
with others. Muntaril produced the movie Premercy with others.
Monsterscar wrote for the movie Premercy. Muntaril was one of the
writers for the movie Zalate. Muntaril wrote for the movie Featsaw.
Teeplemole wrote for the movie Featsaw. Monsterscar was one of the
writers for the movie Zalate. Teeplemole was one of the writers for the
movie Skirtsicine.

Q: Who were born in the year 1910?

A: {[}``Teeplemole'', ``Muntaril''{]}

Q: Who is from the country Piperfish?

A: {[}``Teeplemole'', ``Monsterscar''{]}

movie: Nilitude ; director: Monsterscar. movie: Dewbar ; directed by:
Metatoun. movie: Warpstone ; directed by: Gastrat. movie:
Partnershipmaker ; director: Metatoun. movie: Dewbar ; award:
Tachychronograph. movie: Partnershipmaker ; awarded: Tachychronograph.
movie: Nilitude ; award: Paleodactyl. movie: Warpstone ; award:
Sabonade. person: Gastrat ; award: Trifogation. award: Polyquadrase ;
winner: Monsterscar. award: Trifogation ; winner: Metatoun. movie:
Warpstone ; release year: 1956. movie: Dewbar ; release year: 1984.
movie: Nilitude ; year: 1984. movie: Partnershipmaker ; year: 1962.
Gastrat was an actor in the movie Partnershipmaker. Metatoun was an
actor in the movie Partnershipmaker. Metatoun was an actor in the movie
Nilitude. Gastrat acted in the movie Nilitude. Monsterscar was an actor
in the movie Dewbar. Gastrat acted in the movie Warpstone. Metatoun
acted in the movie Warpstone. Metatoun was born in 1939. Gastrat was
born in the year 1933. Monsterscar was born in 1933. Metatoun grew up in
the nation of Moulole. Gastrat is from the country of Stridery.
Monsterscar grew up in the nation of Moulole. Monsterscar produced the
movie Nilitude with others. Monsterscar was one of the producers of the
movie Warpstone . Metatoun was one of the producers of the movie
Warpstone. Gastrat was one of the producers of the movie Nilitude.
Metatoun produced the movie Partnershipmaker with others. Metatoun
produced the movie Dewbar with others. Monsterscar was one of the
producers of the movie Partnershipmaker. Gastrat produced the movie
Dewbar with others. Metatoun wrote for the movie Partnershipmaker.
Gastrat wrote for the movie Warpstone. Gastrat was one of the writers
for the movie Dewbar. Monsterscar was one of the writers for the movie
Nilitude. Metatoun wrote for the movie Warpstone.

Q: Who is from the country Stridery?

A: {[}''Gastrat''{]}

Q: Who were born in the year 1939?

A: {[}``Metatoun''{]}

\section{COT(small)}\label{cotsmall}

movie: Premercy ; directed by: Muntaril. movie: Skirtsicine ; director:
Teeplemole. movie: Featsaw ; directed by: Monsterscar. movie: Zalate ;
director: Monsterscar. movie: Zalate ; awarded:

Hallowcock. movie: Featsaw ; awarded: Zorgion. movie: Premercy ; award:
Chowwurst. movie: Skirtsicine ; award: Hallowcock. award: Goatfly ;
winner: Teeplemole. person: Monsterscar ; award: Glodome. person:
Muntaril ; award: Goatfly. movie: Featsaw ; release year: 1973. movie:
Zalate ; release year: 1964. movie: Skirtsicine ; release year: 1973.
movie: Premercy ; year: 1961. Teeplemole was an actor in the movie
Skirtsicine. Muntaril was an actor in the movie Skirtsicine. Monsterscar
was an actor in the movie Premercy. Muntaril was an actor in the movie
Featsaw. Teeplemole was an actor in the movie Zalate. Muntaril was born
in the year 1910. Teeplemole was born in 1910. Monsterscar was born in
1942. Teeplemole is from the country of Piperfish. Monsterscar is from
the country of Piperfish. Muntaril is from the country of Clony.
Muntaril produced the movie Skirtsicine with others. Monsterscar was one
of the producers of the movie Featsaw. Monsterscar produced the movie
Premercy with others. Monsterscar produced the movie Zalate with others.
Teeplemole was one of the producers of the movie Featsaw. Teeplemole
produced the movie Zalate with others. Muntaril produced the movie
Premercy with others. Monsterscar wrote for the movie Premercy. Muntaril
was one of the writers for the movie Zalate. Muntaril wrote for the
movie Featsaw. Teeplemole wrote for the movie Featsaw. Monsterscar was
one of the writers for the movie Zalate. Teeplemole was one of the
writers for the movie Skirtsicine.

QC: What awards have movies produced by people born in 1910 won?

QS: The people born in 1910 were Muntaril and Teeplemole. Teeplemole
produced the movies Featsaw and Zalate. Muntaril produced the movie
Premercy. Featsaw was awarded the Zorgion award. Premercy was awarded
the Chowwurst award. Zalate was awarded the Hallowcock award. So the
answer is {[}``Zorgion'', ``Chowwurst'', ``Hallowcock''{]}.

A: {[}``Zorgion'', ``Chowwurst'', ``Hallowcock''{]}

QS: {[}EOQ{]}

movie: Nilitude ; director: Monsterscar. movie: Dewbar ; directed by:
Metatoun. movie: Warpstone ; directed by: Gastrat. movie:
Partnershipmaker ; director: Metatoun. movie: Dewbar ; award:
Tachychronograph. movie: Partnershipmaker ; awarded: Tachychronograph.
movie: Nilitude ; award: Paleodactyl. movie: Warpstone ; award:
Sabonade. person: Gastrat ; award: Trifogation. award: Polyquadrase ;
winner: Monsterscar. award: Trifogation ; winner: Metatoun. movie:
Warpstone ; release year: 1956. movie: Dewbar ; release year: 1984.
movie: Nilitude ; year: 1984. movie: Partnershipmaker ; year: 1962.
Gastrat was an actor in the movie Partnershipmaker. Metatoun was an
actor in the movie Partnershipmaker. Metatoun was an actor in the movie
Nilitude. Gastrat acted in the movie Nilitude. Monsterscar was an actor
in the movie Dewbar. Gastrat acted in the movie Warpstone. Metatoun
acted in the movie Warpstone. Metatoun was born in 1939. Gastrat was
born in the year 1933. Monsterscar was born in 1933. Metatoun grew up in
the nation of Moulole. Gastrat is from the country of Stridery.
Monsterscar grew up in the nation of Moulole. Monsterscar produced the
movie Nilitude with others. Monsterscar was one of the producers of the
movie Warpstone . Metatoun was one of the producers of the movie
Warpstone. Gastrat was one of the producers of the movie Nilitude.
Metatoun produced the movie Partnershipmaker with others. Metatoun
produced the movie Dewbar with others. Monsterscar was one of the
producers of the movie Partnershipmaker. Gastrat produced the movie
Dewbar with others. Metatoun wrote for the movie Partnershipmaker.
Gastrat wrote for the movie Warpstone. Gastrat was one of the writers
for the movie Dewbar. Monsterscar was one of the writers for the movie
Nilitude. Metatoun wrote for the movie Warpstone.

QC: What movies have people from the country Stridery acted in?

QS: The people born in Stridery are Gastrat. Gastrat acted in
Partnershipmaker, Nilitude and Warpstone. So the answer is
{[}``Partnershipmaker'', ``Nilitude'', ``Warpstone''{]}.

A: {[}``Partnershipmaker'', ``Nilitude'', ``Warpstone''{]}

QS: {[}EOQ{]}

\section{G.5 MATH QA}\label{g.5-math-qa}

The decomposer here deterministically calls cot to generate the CoT and
then calls gpt\_ans to extract the answer.

Q: There are 15 trees in the grove. Grove workers will plant trees in
the grove today. After they are done, there will be 21 trees. How many
trees did the grove workers plant today?

A: There are 15 trees originally. Then there were 21 trees after some
more were planted. So there must have been 21 - 15 = 6. The answer is 6.

Q: If there are 3 cars in the parking lot and 2 more cars arrive, how
many cars are in the parking lot?

A: There are originally 3 cars. 2 more cars arrive. \(3 + 2 = 5\) . The
answer is 5.

Q: Leah had 32 chocolates and her sister had 42. If they ate 35, how
many pieces do they have left in total?

A: Originally, Leah had 32 chocolates. Her sister had 42. So in total
they had \(32 + 42 = 74\) . After eating 35, they had 74 - 35 = 39. The
answer is 39.

Q: Jason had 20 lollipops. He gave Denny some lollipops. Now Jason has
12 lollipops. How many lollipops did Jason give to Denny?

A: Jason started with 20 lollipops. Then he had 12 after giving some to
Denny. So he gave Denny \(20 - 12 = 8\) . The answer is 8.

Q: Shawn has five toys. For Christmas, he got two toys each from his mom
and dad. How many toys does he have now?

A: Shawn started with 5 toys. If he got 2 toys each from his mom and
dad, then that is 4 more toys. \(5 + 4 = 9\) . The answer is 9.

Q: There were nine computers in the server room. Five more computers
were installed each day, from monday to thursday. How many computers are
now in the server room?

A: There were originally 9 computers. For each of 4 days, 5 more
computers were added. So 5 * 4 = 20 computers were added. 9 + 20 is 29.
The answer is 29.

Q: Michael had 58 golf balls. On tuesday, he lost 23 golf balls. On
wednesday, he lost 2 more. How many golf balls did he have at the end of
wednesday?

A: Michael started with 58 golf balls. After losing 23 on tuesday, he
had 58 - 23 = 35. After losing 2 more, he had 35 - 2 = 33 golf balls.
The answer is 33.

Q: Olivia has \$23. She bought five bagels for \$3 each. How much money
does she have left?

A: Olivia had 23 dollars. 5 bagels for 3 dollars each will be 5 x 3 = 15
dollars. So she has 23 - 15 dollars left. 23 - 15 is 8. The answer is 8.

\section{gpt\_ans}\label{gpt_ans}

Q: There are 15 trees originally. Then there were 21 trees after some
more were planted. So there must have been 21 - 15 = 6 trees planted.

A: 6

Q: There are originally 3 cars. 2 more cars arrive. \(3 + 2 = 5\) . The
answer is 5.

A: 5

Q: Originally, Leah had 32 chocolates. Her sister had 42. So in total
they had \(32 + 42 = 74\) . After eating 35, they had 74 - 35 = 39. The
answer is 39.

A: 39

Q: Jason started with 20 lollipops. Then he had 12 after giving some to
Denny. So he gave Denny 20 - 12 = 8 apples.

A: 8

Q: Shawn started with 5 toys. If he got 2 toys each from his mom and
dad, then that is 4 more toys. \(5 + 4 = 9\) . The answer is 9.

A: 9

Q: There were originally 9 computers. For each of 4 days, 5 more
computers were added. So \(5 \times 4 = 20\) computers were added.
\(9 + 20\) is 29. The answer is 29.\\
A: 29\\
Q: Michael started with 58 golf balls. After losing 23 on tuesday, he
had 58 - 23 = 35. After losing 2 more, he had 35 - 2 = 33 golf balls.\\
A: 33\\
Q: Olivia had 23 dollars. 5 bagels for 3 dollars each will be 5 x 3 = 15
dollars. So she has 23 - 15 = 8 dollars left.\\
A: 8

\section{G.6 OPEN DOMAIN QA}\label{g.6-open-domain-qa}

The prompts in this section implement Decomposed Prompting approach to
open-domain multihop QA. For brevity we've included prompts for 5 of 20
randomly sampled questions. The full prompts are attached with the
submission and will also be released with the code. Note that we
selected a set of 100 questions from the development set to tune the
hyperparameter (number of paragraphs to retrieve for all of the
retrieval-based approaches).

\section{hotpotqa: decomp}\label{hotpotqa-decomp}

QC: In which country did this Australian who was detained in Guantanamo
Bay detention camp and published ``Guantanamo: My Journey'' receive
para--military training?

QS: (select) {[}retrieve\_odqa{]} Who is the Australian who was detained
in Guantanamo Bay detention camp and published ``Guantanamo: My
Journey''?

A: \{ ``titles'': {[}``Guananamo: My Journey'', ``Bismullah v. Gates'',
``Guananamo Bay detention camp''{]}, ``answer'': {[}``David Hicks''{]}\}

QS: (select) {[}retrieve\_odqa{]} In which country did David Hicks
receive his para--military training?

A: \{''titles'': {[}``John Adams Project'', ``Camp Echo (Guantanamo
Bay)'', ``Guantanamo Bay Museum of Art and History'', ``David
Hicks''{]}, ``answer'': {[}``Afghanistan''{]}\}

QS: (select) {[}multihop\_titleqa{]} Titles: {[}``Guananamo: My
Journey'', ``Bismullah v. Gates'', ``Guananamo Bay detention camp'',
``John Adams Project'', ``Camp Echo (Guananamo Bay)'', ``Guananamo Bay
Museum of Art and History'', ``David Hicks''{]}. Question: In which
country did this Australian who was detained in Guantanamo Bay detention
camp and published ``Guananamo: My Journey'' receive para--military
training?

A: {[}``Afghanistan''{]}

QS: {[}EOQ{]}

QC: How many awards did the ``A Girl Like Me'' singer win at the
American Music Awards of 2012?

QS: (select) {[}retrieve\_odqa{]} Who is the singer of ``A Girl Like
Me''?

A: \{``titles'': {[}``American Music Award'', ``Rihanna''{]},
``answer'': {[}``Rihanna''{]}\}

QS: (select) {[}retrieve\_odqa{]} How many awards did Rihanna won in the
American Music Awards of 2012?

A: \{''titles'': {[}``American Music Awards of 2012'', ``APRA Music
Awards of 2012'', ``List of awards and nominations received by TLC'',
``Native American Music Awards''{]}, ``answer'': {[}``one''{]}\}

QS: (select) {[}multihop\_titleqa{]} Titles: {[}``American Music
Award'', ``Rihanna'', ``American Music Awards of 2012'', ``APRA Music
Awards of 2012'', ``List of awards and nominations received by TLC'',
``Native American Music Awards''{]}. Question: How many awards did the
``A Girl Like Me'' singer win at the American Music Awards of 2012?

A: {[}``one''{]}

QS: {[}EOQ{]}

QC: who is older Jeremy Horn or Renato Sobral?

QS: (select) {[}retrieve\_odqa{]} When was Jeremy Horn born?

A: \{``titles'': {[}``Zaza Tkeshelashvili'', ``Jeremy Horn'', ``Jeremy
Horn (singer)'', ``Ricardo Arona''{]}, ``answer'': {[}``August 25,
1975''{]}\}

QS: (select) {[}retrieve\_odqa{]} When was Renato Sobral born?

A: \{``titles'': {[}``Brian Warren'', ``Renato Sobral''{]}, ``answer'':
{[}``September 7, 1975''{]}\}

QS: (select) {[}multihop\_titleqa{]} Titles: {[}``Zaza Tkeshelashvili'',
``Jeremy Horn'', ``Jeremy Horn (singer)'', ``Ricardo Arona'', ``Brian
Warren'', ``Renato Sobral''{]}. Question: who is older Jeremy Horn or
Renato Sobral?

A: {[}``Jeremy Horn''{]}

QS: {[}EOQ{]}

QC: What was the 2014 population of the city where Lake Wales Medical
Center is located?

QS: (select) {[}retrieve\_odqa{]} Lake Wales Medical Center is located
in what city?

A: \{''titles'': {[}``Baylor College of Medicine'', ``Lake Wales Medical
Center'', ``Tufts University School of Medicine'', ``Hanford Community
Medical Center''{]}, ``answer'': {[}``Polk County, Florida''{]}\}

QS: (select) {[}retrieve\_odqa{]} What was the population of Polk County
in 2014?

A: \{''titles'': {[}''Banner University Medical Center Tucson'', ''Lake
Wales, Florida''{]}, ''answer'': {[}''15,140''{]}\}

QS: (select) {[}multihop\_titleqa{]} Titles: {[}``Baylor College of
Medicine'', ``Lake Wales Medical Center'', ``Tufts University School of
Medicine'', ``Hanford Community Medical Center'', ``Banner University
Medical Center Tucson'', ``Lake Wales, Florida''{]}. Question: What was
the 2014 population of the city where Lake Wales Medical Center is
located?

A: {[}``15,140''{]}

QS: {[}EOQ{]}

QC: Nobody Loves You was written by John Lennon and released on what
album that was issued by Apple Records, and was written, recorded, and
released during his 18 month separation from Yoko Ono?

QS: (select) {[}retrieve\_odqa{]} What album was issued by Apple
Records, and written, recorded, and released during John Lennon's 18
month separation from Yoko Ono?

A: \{''titles'': {[}``John Lennon/Plastic Ono Band'', ``Milk and Honey
(album)'', ``Walls and Bridges''{]}, ``answer'': {[}``Walls and
Bridges''{]}\}

QS: (select) {[}retrieve\_odqa{]} Nobody Loves You was written by John
Lennon on what album?

A: \{''titles'': {[}''John Lennon Museum'', ``Nobody Loves You (When
You're Down and Out)'', ''answer'': {[}''Walls and Bridges''{]}\}

QS: (select) {[}multihop\_titleqa{]} Titles: {[}``John Lennon/Plastic
Ono Band'', ``Milk and Honey (album)'', ``Walls and Bridges'', ``John
Lennon Museum'', ``Nobody Loves You (When You're Down and Out)''{]}.
Question: Nobody Loves You was written by John Lennon and released on
what album that was issued by Apple Records, and was written, recorded,
and released during his 18 month separation from Yoko Ono?

A: {[}``Walls and Bridges''{]}

QS: {[}EOQ{]}

\section{hotpotqa: retrieve\_odqa}\label{hotpotqa-retrieve_odqa}

QC: Who is the Australian who was detained in Guantanamo Bay detention
camp and published ``Guantanamo: My Journey''?

QS: (select) {[}retrieve{]} Who is the Australian who was detained in
Guantanamo Bay detention camp and published ``Guantanamo: My Journey''?

A: {[}``Guantanamo: My Journey'', ``Bismullah v. Gates'', ``Guantanamo
Bay detention camp''{]}

QS: (select) {[}singlehop\_titleqa{]} Titles: {[}``Guantanamo: My
Journey'', ``Bismullah v. Gates'', ``Guantanamo Bay detention camp''{]}.
Question: Who is the Australian who was detained in Guantanamo Bay
detention camp and published ``Guantanamo: My Journey''?

A: \{''titles'': {[}``Guananamo: My Journey'', ``Bismullah v. Gates'',
``Guananamo Bay detention camp''{]}, ``answer'': {[}``David Hicks''{]}\}

QS: {[}EOQ{]}

QC: In which country did David Hicks receive his para-military training?

QS: (select) {[}retrieve{]} In which country did David Hicks receive his
para--military training?

A: {[}``John Adams Project'', ``Camp Echo (Guantanamo Bay)'',
``Guantanamo Bay Museum of Art and History'', ``David Hicks''{]}

QS: (select) {[}singlehop\_titleqa{]} Titles: {[}``John Adams Project'',
``Camp Echo (Guantanamo Bay)'', ``Guantanamo Bay Museum of Art and
History'', ``David Hicks''{]}. Question: In which country did David
Hicks receive his para--military training?

A: \{''titles'': {[}``John Adams Project'', ``Camp Echo (Guantanamo
Bay)'', ``Guantanamo Bay Museum of Art and History'', ``David
Hicks''{]}, ``answer'': {[}``Afghanistan''{]}\}

QS: {[}EOQ{]}

QC: Who is the singer of ``A Girl Like Me''?

QS: (select) {[}retrieve{]} Who is the singer of ``A Girl Like Me''?

A: {[}``American Music Award'', ``Rihanna''{]}

QS: (select) {[}singlehop\_titleqa{]} Titles: {[}``American Music
Award'', ``Rihanna''{]}. Question: Who is the singer of ``A Girl Like
Me''?

A: \{``titles'': {[}``American Music Award'', ``Rihanna''{]},
``answer'': {[}``Rihanna''{]}\}

QS: {[}EOQ{]}

QC: How many awards did Rihanna won in the American Music Awards of
2012?

QS: (select) {[}retrieve{]} How many awards did Rihanna won in the
American Music Awards of 2012?

A: {[}``American Music Awards of 2012'', ``APRA Music Awards of 2012'',
``List of awards and nominations received by TLC'', ``Native American
Music Awards''{]}

QS: (select) {[}singlehop\_titleqa{]} Titles: {[}``American Music Awards
of 2012'', ``APRA Music Awards of 2012'', ``List of awards and
nominations received by TLC'', ``Native American Music Awards''{]}.
Question: How many awards did Rihanna won in the American Music Awards
of 2012?

A: \{''titles'': {[}``American Music Awards of 2012'', ``APRA Music
Awards of 2012'', ``List of awards and nominations received by TLC'',
``Native American Music Awards''{]}, ``answer'': {[}``one''{]}\}

QS: {[}EOQ{]}

QC: When was Jeremy Horn born?

QS: (select) {[}retrieve{]} When was Jeremy Horn born?

A: {[}``Zaza Tkeshelashvili'', ``Jeremy Horn'', ``Jeremy Horn
(singer)'', ``Ricardo Arona''{]}

QS: (select) {[}singlehop\_titleqa{]} Titles: {[}``Zaza
Tkeshelashvili'', ``Jeremy Horn'', ``Jeremy Horn (singer)'', ``Ricardo
Arona''{]}. Question: When was Jeremy Horn born?

A: \{''titles'': {[}``Zaza Tkeshelashvili'', ``Jeremy Horn'', ``Jeremy
Horn (singer)'', ``Ricardo Arona''{]}, ``answer'': {[}``August 25,
1975''{]}\}

QS: {[}EOQ{]}

QC: When was Renato Sobral born?

QS: (select) {[}retrieve{]} When was Renato Sobral born?

A: {[}``Brian Warren'', ``Renato Sobral''{]}

QS: (select) {[}singlehop\_titleqa{]} Titles: {[}''Brian Warren'',
''Renato Sobral''{]}. Question: When was Renato Sobral born?

A: \{``titles'': {[}``Brian Warren'', ``Renato Sobral''{]}, ``answer'':
{[}``September 7, 1975''{]}\}

QS: {[}EOQ{]}

QC: Lake Wales Medical Center is located in what city?

QS: (select) {[}retrieve{]} Lake Wales Medical Center is located in what
city?

A: {[}``Baylor College of Medicine'', ``Lake Wales Medical Center'',
``Tufts University School of Medicine'', ``Hanford Community Medical
Center''{]}

QS: (select) {[}singlehop\_titleqa{]} Titles: {[}``Baylor College of
Medicine'', ``Lake Wales Medical Center'', ``Tufts University School of
Medicine'', ``Hanford Community Medical Center''{]}. Question: Lake
Wales Medical Center is located in what city?

A: \{``titles'': {[}``Baylor College of Medicine'', ``Lake Wales Medical
Center'', ``Tufts University School of Medicine'', ``Hanford Community
Medical Center''{]}, ``answer'': {[}``Polk County, Florida''{]}\}

QS: {[}EOQ{]}

QC: What was the population of Polk County in 2014?

QS: (select) {[}retrieve{]} What was the population of Polk County in
2014?

A: {[}``Banner University Medical Center Tucson'', ``Lake Wales,
Florida''{]}

QS: (select) {[}singlehop\_titleqa{]} Titles: {[}``Banner University
Medical Center Tucson'', ``Lake Wales, Florida''{]}. Question: What was
the population of Polk County in 2014?

A: \{``titles'': {[}``Banner University Medical Center Tucson'', ``Lake
Wales, Florida''{]}, ``answer'': {[}``15,140''{]}\}

QS: {[}EOQ{]}

QC: What album was issued by Apple Records, and written, recorded, and
released during John Lennon's 18 month separation from Yoko Ono?

QS: (select) {[}retrieve{]} What album was issued by Apple Records, and
written, recorded, and released during John Lennon's 18 month separation
from Yoko Ono?

A: {[}``John Lennon/Plastic Ono Band'', ``Milk and Honey (album)'',
``Walls and Bridges''{]}

QS: (select) {[}singlehop\_titleqa{]} Titles: {[}``John Lennon/Plastic
Ono Band'', ``Milk and Honey (album)'', ``Walls and Bridges''{]}.
Question: What album was issued by Apple Records, and written, recorded,
and released during John Lennon's 18 month separation from Yoko Ono?

A: \{ ``titles'': {[}``John Lennon/Plastic Ono Band'', ``Milk and Honey
(album)'', ``Walls and Bridges''{]}, ``answer'': {[}``Walls and
Bridges''{]} \}

QS: {[}EOQ{]}

QC: Nobody Loves You was written by John Lennon on what album?

QS: (select) {[}retrieve{]} Nobody Loves You was written by John Lennon
on what album?

A: {[}``John Lennon Museum'', ``Nobody Loves You (When You're Down and
Out)''{]}

QS: (select) {[}singlehop\_titleqa{]} Titles: {[}'John Lennon Museum'',
''Nobody Loves You (When You're Down and Out)''. Question: Nobody Loves
You was written by John Lennon on what album?

A: \{``titles'': {[}``John Lennon Museum'', ``Nobody Loves You (When
You're Down and Out){]}``,''answer'': {[}``Walls and Bridges''{]}\}

QS: {[}EOQ{]}

hotpotqa: singlehop\_titleqa

Wikipedia Title: David Hicks

Wikipedia Title: Guantanamo: My Journey

Wikipedia Title: Murat Kurnaz

Q: Who is the Australian who was detained in Guantanamo Bay detention
camp and published ``Guantanamo: My Journey''?

A: {[}``David Hicks''{]}

Q: In which country did David Hicks receive his para--military training?

A: {[}``Afghanistan''{]}

Wikipedia Title: American Music Award

Wikipedia Title: Rihanna

Wikipedia Title: American Music Awards of 2012

Q: Who is the singer of ``A Girl Like Me''? A: {[}``Rihanna''{]}

Q: How many awards did Rihanna won in the American Music Awards of 2012?
A: {[}``one''{]}

Wikipedia Title: Renato Sobral

Wikipedia Title: Zaza Tkeshelashvili

Wikipedia Title: Jeremy Horn

Q: When was Jeremy Horn born? A: {[}``August 25, 1975''{]}

Q: When was Renato Sobral born? A: {[}``September 7, 1975''{]}

Wikipedia Title: Lake Wales Medical Center

Wikipedia Title: Lake Wales, Florida

Wikipedia Title: Tufts University School of Medicine

Q: Lake Wales Medical Center is located in what city? A: {[}``Polk
County, Florida''{]}

Q: What was the population of Polk County in 2014? A: {[}``15,140''{]}

Wikipedia Title: Nobody Loves You (When You're Down and Out)

Wikipedia Title: Walls and Bridges

Wikipedia Title: Mother (John Lennon song)

Q: What album was issued by Apple Records, and written, recorded, and
released during John Lennon's 18 month separation from Yoko Ono?

A: {[}``Walls and Bridges''{]}

Q: Nobody Loves You was written by John Lennon on what album?

A: {[}``Walls and Bridges''{]}

hotpotqa: multihop\_titleqa (direct)

Wikipedia Title: David Hicks

Wikipedia Title: Guantanamo: My Journey

Wikipedia Title: Murat Kurnaz

Q: In which country did this Australian who was detained in Guantanamo
Bay detention camp and published ``Guantanamo: My Journey'' receive
para--military training?

A: {[}``Afghanistan''{]}

Wikipedia Title: American Music Award

Wikipedia Title: Rihanna

Wikipedia Title: American Music Awards of 2012

Q: How many awards did the ``A Girl Like Me'' singer win at the American
Music Awards of 2012?

A: {[}``one''{]}

Wikipedia Title: Renato Sobral

Wikipedia Title: Zaza Tkeshelashvili

Wikipedia Title: Jeremy Horn

Q: who is older Jeremy Horn or Renato Sobral?

A: {[}``Jeremy Horn''{]}

Wikipedia Title: Lake Wales Medical Center

Wikipedia Title: Lake Wales, Florida

Wikipedia Title: Tufts University School of Medicine

Q: What was the 2014 population of the city where Lake Wales Medical
Center is located?

A: {[}''15,140''{]}

Wikipedia Title: Nobody Loves You (When You're Down and Out)

Wikipedia Title: Walls and Bridges

Wikipedia Title: Mother (John Lennon song)

Q: Nobody Loves You was written by John Lennon and released on what
album that was issued by Apple Records, and was written, recorded, and
released during his 18 month separation from Yoko Ono?

A: {[}``Walls and Bridges''{]}

hotpotqa: multihop\_titleqa (cot)

Wikipedia Title: David Hicks

Wikipedia Title: Guantanamo: My Journey

Wikipedia Title: Murat Kurnaz

Q: In which country did this Australian who was detained in Guantanamo
Bay detention camp and published ``Guantanamo: My Journey'' receive
para--military training?

A: The Australian who was detained in Guantanamo Bay detention camp and
published ``Guantanamo: My Journey'' is David Hicks. David Hicks
received his para--military training in Afghanistan. So the answer is:
Afghanistan.

Wikipedia Title: American Music Award

Wikipedia Title: Rihanna

Wikipedia Title: American Music Awards of 2012

Q: How many awards did the ``A Girl Like Me'' singer win at the American
Music Awards of 2012?

A: The singer of ``A Girl Like Me'' singer is Rihanna. In the American
Music Awards of 2012, Rihanna won one award. So the answer is: one.

Wikipedia Title: Renato Sobral

Wikipedia Title: Zaza Tkeshelashvili

Wikipedia Title: Jeremy Horn

Q: who is older Jeremy Horn or Renato Sobral?

A: Jeremy Horn was born on August 25, 1975. Renato Sobral was born on
September 7, 1975. Thus, Jeremy Horn is older. So the answer is: Jeremy
Horn.

Wikipedia Title: Lake Wales Medical Center

Wikipedia Title: Lake Wales, Florida

Wikipedia Title: Tufts University School of Medicine

Q: What was the 2014 population of the city where Lake Wales Medical
Center is located?

A: Lake Wales Medical Center is located in the city of Polk County,
Florida. The population of Polk County in 2014 was 15,140. So the answer
is: 15,140.

Wikipedia Title: Nobody Loves You (When You're Down and Out)

Wikipedia Title: Walls and Bridges

Wikipedia Title: Mother (John Lennon song)

Q: Nobody Loves You was written by John Lennon and released on what
album that was issued by Apple Records, and was written, recorded, and
released during his 18 month separation from Yoko Ono?

A: The album issued by Apple Records, and written, recorded, and
released during John Lennon's 18 month separation from Yoko Ono is Walls
and Bridges. Nobody Loves You was written by John Lennon on Walls and
Bridges album. So the answer is: Walls and Bridges.

\section{2wikimultihopqa: decomp}\label{wikimultihopqa-decomp}

QC: Do director of film Coolie No.~1 (1995 Film) and director of film
The Sensational Trial have the same nationality?

QS: (select) {[}retrieve\_odqa{]} Who is the director of the film Coolie
No.~1 (1995 film)?

A: \{``titles'': {[}``Hanro Smitsman'', ``Coolie No.~1 (1995 film)''{]},
``answer'': {[}``David Dhawan''{]}\}

QS: (select) {[}retrieve\_odqa{]} Who is the director of the film The
Sensational Trial?

A: \{``titles'': {[}``Rachel Feldman'', ``The Sensational Trial''{]},
``answer'': {[}``Karl Freund''{]}\}

QS: (select) {[}retrieve\_odqa{]} What is David Dhawan's nationality?

A: \{``titles'': {[}``David Dhawan'', ``Brian Johnson (special effects
artist){]}``,''answer'': {[}``India''{]}\}

QS: (select) {[}retrieve\_odqa{]} What is Karl Freund's nationality?

A: \{``titles'': {[}``Karl Freund'', ``Ian Barry (director)''{]},
``answer'': {[}``Germany''{]}\}

QS: (select) {[}multihop\_titleqa{]} Titles: {[}``Hanro Smitsman'',
``Coolie No.~1 (1995 film)'', ``Rachel Feldman'', ``The Sensational
Trial'', ``David Dhawan'', ``Brian Johnson (special effects artist)'',
``Karl Freund'', ``Ian Barry (director)''{]}. Question: Do director of
film Coolie No.~1 (1995 Film) and director of film The Sensational Trial
have the same nationality?

A: {[}``no''{]} QS: {[}EOQ{]}

QC: Which film has the director died later, The Gal Who Took the West or
Twenty Plus Two?

QS: (select) {[}retrieve\_odqa{]} Who is the director of the film Twenty
Plus?

A: \{``titles'': {[}``Riki Gal'', ``Twenty Plus Two''{]}, ``answer'':
{[}``Joseph M. Newman''{]}\}

QS: (select) {[}retrieve\_odqa{]} Who is the director of the film The
Gal Who Took the West?

A: \{``titles'': {[}``Querelle'', ``The Gal Who Took the West''{]},
``answer'': {[}``Frederick de Cordova''{]}\}

QS: (select) {[}retrieve\_odqa{]} When did Joseph M. Newman die?

A: \{``titles'': {[}``Joseph M. Newman'', ``Thulasi (1987 film)''{]},
``answer'': {[}``January 23, 2006''{]}\}\\
QS: (select) {[}retrieve\_odqa{]} When did Fred de Cordova die?\\
A: \{``titles'': {[}``Fred de Cordova'', ``Le Masque de la Meduse''{]},
``answer'': {[}``September 15, 2001''{]}\}\\
QS: (select) {[}multihop\_titleqa{]} Titles: {[}``Riki Gal'', ``Twenty
Plus Two'', ``Querelle'', ``The Gal Who Took the West'', ``Joseph M.
Newman'', ``Thulasi (1987 film)'', ``Fred de Cordova'', ``Le Masque de
la Meduse''{]}. Question: Which film has the director died later, The
Gal Who Took the West or Twenty Plus Two?\\
A: {[}``Twenty Plus Two''{]}\\
QS: {[}EOQ{]}\\
QC: Who is the grandchild of Krishna Shah (Nepalese Royal)?\\
QS: (select) {[}retrieve\_odqa{]} Who is the child of Krishna Shah?\\
A: \{''titles'': {[}``Ana Gruzinsky--Golitsyn'', ``Krishna Shah
(Nepalese royal)'', ``Diana Weston'', ``Albina du Boisrouvray''{]},
``answer'': {[}``Rudra Shah''{]}\}\\
QS: (select) {[}retrieve\_odqa{]} Who is the child of Rudra Shah?\\
A: \{''titles'': {[}''Jim Ramel Kjellgren'', ''Constance Anne
Herschel'', ''Rudra Shah''{]}, ''answer'': {[}''Prithvipati
Shah''{]}\}\\
QS: (select) {[}multihop\_titleqa{]} Titles: {[}``Ana
Gruzinsky-Golitsyn'', ``Krishna Shah (Nepalese royal)'', ``Diana
Weston'', ``Albina du Boisrouvray'', ``Jim Ramel Kjellgren'',
``Constance Anne Herschel'', ``Rudra Shah''{]}. Question: Who is the
grandchild of Krishna Shah (Nepalese Royal)?\\
A: {[}``Prithvipati Shah''{]}\\
QS: {[}EOQ{]}\\
QC: Where did the director of film Maddalena (1954 Film) die?\\
QS: (select) {[}retrieve\_odqa{]} Who is the director of the film
Maddalena?\\
A: \{''titles'': {[}''Ian Barry (director)'', ''Maddalena (1954 film)'',
''Brian Kennedy (gallery director)'', ''answer'': {[}''Augusto
Genina''{]}\}\\
QS: (select) {[}retrieve\_odqa{]} Where did Augusto Genina die?\\
A: \{ ``titles'': {[}``Dana Blankstein'', ``Augusto Genina'', ``Peter
Levin'', ``Olav Aaraas''{]}, ``answer'': {[}'' Rome''{]}\}\\
QS: (select) {[}multihop\_titleqa{]} Titles: {[}``Ian Barry
(director)'', ``Maddalena (1954 film)'', ``Brian Kennedy (gallery
director)'', ``Dana Blankstein'', ``Augusto Genina'', ``Peter Levin'',
``Olav Aaraas''{]}. Question: Where did the director of film Maddalena
(1954 Film) die?\\
A: {[}``Rome''{]}\\
QS: {[}EOQ{]}\\
QC: Who was born first out of Martin Hodge and Ivania Martinich?\\
QS: (select) {[}retrieve\_odqa{]} When was Martin Hodge born?\\
A: \{''titles'': {[}``Greg A. Hill (artist)'', ``Martin Hodge'', ``John
Allen (Oxford University cricketer){]}``,''answer'': {[}``4 February
1959''{]}\}\\
QS: (select) {[}retrieve\_odqa{]} When was Ivania Martinich born?\\
A: \{''titles'': {[}``Tom Dickinson'', ``Ivania Martinich'', ``Wesley
Barresi'', ``Henry Moore (cricketer)''{]}, ``answer'': {[}``25 July
1995''{]}\}\\
QS: (select) {[}multihop\_titleqa{]} Titles: {[}``Greg A. Hill
(artist)'', ``Martin Hodge'', ``John Allen (Oxford University
cricketer)'', ``Tom Dickinson'', ``Ivania Martinich'', ``Wesley
Barresi'', ``Henry Moore (cricketer)''{]}. Question: Who was born first
out of Martin Hodge and Ivania Martinich?\\
A: {[}``Martin Hodge''{]}\\
QS: {[}EOQ{]}

2wikimultihopqa: retrieve\_odqa

QC: Who is the director of the film Coolie No.~1 (1995 film)?

QS: (select) {[}retrieve{]} Who is the director of the film Coolie No.~1
(1995 film)?

A: {[}``Hanro Smitsman'', ``Coolie No.~1 (1995 film)''{]}

QS: (select) {[}singlehop\_titleqa{]} Titles: {[}``Hanro Smitsman'',
``Coolie No.~1 (1995 film)''{]}. Question: Who is the director of the
film Coolie No.~1 (1995 film)?

A: \{``titles'': {[}``Hanro Smitsman'', ``Coolie No.~1 (1995 film)''{]},
``answer'': {[}``David Dhawan''{]}\}

QS: {[}EOQ{]}

QC: Who is the director of the film The Sensational Trial?

QS: (select) {[}retrieve{]} Who is the director of the film The
Sensational Trial?

A: {[}``Rachel Feldman'', ``The Sensational Trial''{]}

QS: (select) {[}singlehop\_titleqa{]} Titles: {[}``Rachel Feldman'',
``The Sensational Trial''{]}. Question: Who is the director of the film
The Sensational Trial?

A: \{``titles'': {[}``Rachel Feldman'', ``The Sensational Trial''{]},
``answer'': {[}``Karl Freund''{]}\}

QS: {[}EOQ{]}

QC: What is David Dhawan's nationality?

QS: (select) {[}retrieve{]} What is David Dhawan's nationality?

A: {[}``David Dhawan'', ``Brian Johnson (special effects artist)''{]}

QS: (select) {[}singlehop\_titleqa{]} Titles: {[}``David Dhawan'',
``Brian Johnson (special effects artist)''{]}. Question: What is David
Dhawan's nationality?

A: \{``titles'': {[}``David Dhawan'', ``Brian Johnson (special effects
artist)''{]}, ``answer'': {[}``India''{]}\}

QS: {[}EOQ{]}

QC: What is Karl Freund's nationality?

QS: (select) {[}retrieve{]} What is Karl Freund's nationality?

A: {[}``Karl Freund'', ``Ian Barry (director)''{]}

QS: (select) {[}singlehop\_titleqa{]} Titles: {[}``Karl Freund'', ``Ian
Barry (director)''{]}. Question: What is Karl Freund's nationality?

A: \{``titles'': {[}``Karl Freund'', ``Ian Barry (director)''{]},
``answer'': {[}``Germany''{]}\}

QS: {[}EOQ{]}

QC: Who is the director of the film Twenty Plus?

QS: (select) {[}retrieve{]} Who is the director of the film Twenty Plus?

A: {[}``Riki Gal'', ``Twenty Plus Two''{]}

QS: (select) {[}singlehop\_titleqa{]} Titles: {[}``Riki Gal'', ``Twenty
Plus Two''{]}. Question: Who is the director of the film Twenty Plus?

A: \{``titles'': {[}``Riki Gal'', ``Twenty Plus Two''{]}, ``answer'':
{[}``Joseph M. Newman''{]}\}

QS: {[}EOQ{]}

QC: Who is the director of the film The Gal Who Took the West?

QS: (select) {[}retrieve{]} Who is the director of the film The Gal Who
Took the West?

A: {[}``Querelle'', ``The Gal Who Took the West''{]}

QS: (select) {[}singlehop\_titleqa{]} Titles: {[}``Querelle'', ``The Gal
Who Took the West''{]}. Question: Who is the director of the film The
Gal Who Took the West?

A: \{``titles'': {[}``Querelle'', ``The Gal Who Took the West''{]},
``answer'': {[}``Frederick de Cordova''{]}\}

QS: {[}EOQ{]}

QC: When did Joseph M. Newman die?

QS: (select) {[}retrieve{]} When did Joseph M. Newman die?

A: {[}``Joseph M. Newman'', ``Thulasi (1987 film)''{]}

QS: (select) {[}singlehop\_titleqa{]} Titles: {[}``Joseph M. Newman'',
``Thulasi (1987 film)''{]}. Question: When did Joseph M. Newman die?

A: \{``titles'': {[}``Joseph M. Newman'', ``Thulasi (1987 film)''{]},
``answer'': {[}``January 23, 2006''{]}\}

QS: {[}EOQ{]}

QC: When did Fred de Cordova die?

QS: (select) {[}retrieve{]} When did Fred de Cordova die?

A: {[}``Fred de Cordova'', ``Le Masque de la Meduse''{]}

QS: (select) {[}singlehop\_titleqa{]} Titles: {[}``Fred de Cordova'',
``Le Masque de la Meduse''{]}. Question: When did Fred de Cordova die?

A: \{``titles'': {[}``Fred de Cordova'', ``Le Masque de la Meduse''{]},
``answer'': {[}``September 15, 2001''{]}\}

QS: {[}EOQ{]}

QC: Who is the child of Krishna Shah?

QS: (select) {[}retrieve{]} Who is the child of Krishna Shah?

A: {[}``Ana Gruzinsky--Golitsyn'', ``Krishna Shah (Nepalese royal)'',
``Diana Weston'', ``Albina du Boisrouvray''{]}

QS: (select) {[}singlehop\_titleqa{]} Titles: {[}``Ana
Gruzinsky--Golitsyn'', ``Krishna Shah (Nepalese royal)'', ``Diana
Weston'', ``Albina du Boisrouvray''{]}. Question: Who is the child of
Krishna Shah?

A: \{''titles'': {[}``Ana Gruzinsky--Golitsyn'',''Krishna Shah (Nepalese
royal)'', ``Diana Weston'',''Albina du Boisrouvray''{]}, ``answer'':
{[}``Rudra Shah''{]}\}

QS: {[}EOQ{]}

QC: Who is the child of Rudra Shah?

QS: (select) {[}retrieve{]} Who is the child of Rudra Shah?

A: {[}``Jim Ramel Kjellgren'', ``Constance Anne Herschel'', ``Rudra
Shah''{]}

QS: (select) {[}singlehop\_titleqa{]} Titles: {[}``Jim Ramel
Kjellgren'', ``Constance Anne Herschel'', ``Rudra Shah''{]}. Question:
Who is the child of Rudra Shah?

A: \{''titles'': {[}``Jim Ramel Kjellgren'', ``Constance Anne
Herschel'', ``Rudra Shah''{]}, ``answer'': {[}``Prithvipati Shah''{]}\}

QS: {[}EOQ{]}

QC: Who is the director of the film Maddalena?

QS: (select) {[}retrieve{]} Who is the director of the film Maddalena?

A: {[}``Ian Barry (director)'', ``Maddalena (1954 film)'', ``Brian
Kennedy (gallery director)''{]}

QS: (select) {[}singlehop\_titleqa{]} Titles: {[}'Ian Barry
(director)'', ''Maddalena (1954 film)'', ''Brian Kennedy (gallery
director)''. Question: Who is the director of the film Maddalena?

A: \{''titles'': {[}``Ian Barry (director)'',''Maddalena (1954 film)'',
``Brian Kennedy (gallery director)''{]}, ``answer'': {[}``Augusto
Genina''{]}\}

QS: {[}EOQ{]}

QC: Where did Augusto Genina die?

QS: (select) {[}retrieve{]} Where did Augusto Genina die?

A: {[}``Dana Blankstein'', ``Augusto Genina'', ``Peter Levin'', ``Olav
Aaraas''{]}

QS: (select) {[}singlehop\_titleqa{]} Titles: {[}``Dana Blankstein'',
``Augusto Genina'', ``Peter Levin'', ``Olav Aaraas''{]}. Question: Where
did Augusto Genina die?

A: \{''titles'': {[}``Dana Blankstein'', ``Augusto Genina'', ``Peter
Levin'', ``Olav Aaraas''{]}, ``answer'': {[}``Rome''{]}\}

QS: {[}EOQ{]}

QC: When was Martin Hodge born?

QS: (select) {[}retrieve{]} When was Martin Hodge born?

A: {[}``Greg A. Hill (artist)'', ``Martin Hodge'', ``John Allen (Oxford
University cricketer)''{]}

QS: (select) {[}singlehop\_titleqa{]} Titles: {[}``Greg A. Hill
(artist)'', ``Martin Hodge'', ``John Allen (Oxford University
cricketer)''{]}. Question: When was Martin Hodge born?

A: \{''titles'': {[}``Greg A. Hill (artist)'',''Martin Hodge'', ``John
Allen (Oxford University cricketer){]}``,''answer'': {[}``4 February
1959''{]}\}

QS: {[}EOQ{]}

QC: When was Imania Martinich born?

QS: (select) {[}retrieve{]} When was Ivania Martinich born?

A: {[}``Tom Dickinson'', ``Ivania Martinich'', ``Wesley Barresi'',
``Henry Moore (cricketer)''{]}

QS: (select) {[}singlehop\_titleqa{]} Titles: {[}``Tom Dickinson'',
``Ivania Martinich'', ``Wesley Barresi'', ``Henry Moore
(cricketer)''{]}. Question: When was Ivania Martinich born?

A: \{''titles'': {[}''Tom Dickinson'', ``Ivania Martinich'', ``Wesley
Barresi'', ``Henry Moore (cricketer)'', ``answer'': {[}''25 July
1995''{]}\}

QS: {[}EOQ{]}

2wikimultihopqa: singlehop\_titleqa

Wikipedia Title: David Dhawan

Wikipedia Title: Howard W. Koch

Wikipedia Title: The Sensational Trial

Wikipedia Title: Coolie No.~1 (1995 film)

Wikipedia Title: Karl Freund

Q: Who is the director of the film Coolie No.~1 (1995 film)?

A: {[}``David Dhawan''{]}

Q: Who is the director of the film The Sensational Trial?

A: {[}``Karl Freund''{]}

Q: What is David Dhawan's nationality?

A: {[}``India''{]}

Q: What is Karl Freund's nationality?

A: {[}``Germany''{]}

Wikipedia Title: Fred de Cordova

Wikipedia Title: Thulasi (1987 film)

Wikipedia Title: Joseph M. Newman

Wikipedia Title: The Gal Who Took the West

Wikipedia Title: Twenty Plus Two

Q: Who is the director of the film Twenty Plus?

A: {[}``Joseph M. Newman''{]}

Q: Who is the director of the film The Gal Who Took the West?

A: {[}``Frederick de Cordova''{]}

Q: When did Joseph M. Newman die?

A: {[}``January 23, 2006''{]}

Q: When did Fred de Cordova die?

A: {[}``September 15, 2001''{]}

Wikipedia Title: Ana Gruzinsky--Golitsyn

Wikipedia Title: Rudra Shah

Wikipedia Title: Krishna Shah (Nepalese royal)

Q: Who is the child of Krishna Shah?

A: {[}``Rudra Shah''{]}

Q: Who is the child of Rudra Shah?

A: {[}``Prithvipati Shah''{]}

Wikipedia Title: Augusto Genina

Wikipedia Title: Ian Barry (director)

Wikipedia Title: Maddalena (1954 film)

Q: Who is the director of the film Maddalena?

A: {[}``Augusto Genina''{]}

Q: Where did Augusto Genina die?

A: {[}``Rome''{]}

Wikipedia Title: Martin Hodge

Wikipedia Title: Ivania Martinich

Wikipedia Title: Tom Dickinson

Q: When was Martin Hodge born?

A: {[}``4 February 1959''{]}

Q: When was Ivania Martinich born?

A: {[}``25 July 1995''{]}

2wikimultihopqa: multihop\_titleqa (direct)

Wikipedia Title: David Dhawan

\begin{Shaded}
\begin{Highlighting}[]
\NormalTok{Wikipedia Title: Howard W. Koch}
\NormalTok{\textless{}hidden for brevity\textgreater{}}
\NormalTok{Wikipedia Title: The Sensational Trial}
\NormalTok{\textless{}hidden for brevity\textgreater{}}
\NormalTok{Wikipedia Title: Coolie No. 1 (1995 film)}
\NormalTok{\textless{}hidden for brevity\textgreater{}}
\NormalTok{Wikipedia Title: Karl Freund}
\NormalTok{\textless{}hidden for brevity\textgreater{}}
\NormalTok{Q: Do director of film Coolie No. 1 (1995 Film) and director of film The Sensational Trial have the same nationality?}
\NormalTok{A: ["no"]}
\NormalTok{Wikipedia Title: Fred de Cordova}
\NormalTok{\textless{}hidden for brevity\textgreater{}}
\NormalTok{Wikipedia Title: Thulasi (1987 film)}
\NormalTok{\textless{}hidden for brevity\textgreater{}}
\NormalTok{Wikipedia Title: Joseph M. Newman}
\NormalTok{\textless{}hidden for brevity\textgreater{}}
\NormalTok{Wikipedia Title: The Gal Who Took the West}
\NormalTok{\textless{}hidden for brevity\textgreater{}}
\NormalTok{Wikipedia Title: Twenty Plus Two}
\NormalTok{\textless{}hidden for brevity\textgreater{}}
\NormalTok{Q: Which film has the director died later, The Gal Who Took the West or Twenty Plus Two?}
\NormalTok{A: ["Twenty Plus Two"]}
\NormalTok{Wikipedia Title: Ana Gruzinsky–Golitsyn}
\NormalTok{\textless{}hidden for brevity\textgreater{}}
\NormalTok{Wikipedia Title: Rudra Shah}
\NormalTok{\textless{}hidden for brevity\textgreater{}}
\NormalTok{Wikipedia Title: Krishna Shah (Nepalese royal)}
\NormalTok{\textless{}hidden for brevity\textgreater{}}
\NormalTok{Q: Who is the grandchild of Krishna Shah (Nepalese Royal)?}
\NormalTok{A: ["Prithvipati Shah"]}
\NormalTok{Wikipedia Title: Augusto Genina}
\NormalTok{\textless{}hidden for brevity\textgreater{}}
\NormalTok{Wikipedia Title: Ian Barry (director)}
\NormalTok{\textless{}hidden for brevity\textgreater{}}
\NormalTok{Wikipedia Title: Maddalena (1954 film)}
\NormalTok{\textless{}hidden for brevity\textgreater{} }
\end{Highlighting}
\end{Shaded}

Q: Where did the director of film Maddalena (1954 Film) die?

A: {[}``Rome''{]}

Wikipedia Title: Martin Hodge

Wikipedia Title: Ivania Martinich

Wikipedia Title: Tom Dickinson

Q: Who was born first out of Martin Hodge and Ivania Martinich?

A: {[}``Martin Hodge''{]}

2wikimultihopqa: multihop\_titleqa (cot)

Wikipedia Title: David Dhawan

Wikipedia Title: Howard W. Koch

Wikipedia Title: The Sensational Trial

Wikipedia Title: Coolie No.~1 (1995 film)

Wikipedia Title: Karl Freund

Q: Do director of film Coolie No.~1 (1995 Film) and director of film The
Sensational Trial have the same nationality?

A: Coolie No.~1 (1995 film) was directed by David Dhawan. The
Sensational Trial was directed by Karl Freund. David Dhawan's
nationality is India. Karl Freund's nationality is Germany.

Thus, they do not have the same nationality. So the answer is: no.

Wikipedia Title: Fred de Cordova

Wikipedia Title: Thulasi (1987 film)

Wikipedia Title: Joseph M. Newman

Wikipedia Title: The Gal Who Took the West

Wikipedia Title: Twenty Plus Two

Q: Which film has the director died later, The Gal Who Took the West or
Twenty Plus Two?

A: The film Twenty Plus Two was directed by Joseph M. Newman. The Gal
Who Took the West was directed by Frederick de Cordova. Joseph M. Newman
died on January 23, 2006. Fred de Cordova died on September 15, 2001.
Thus, the person to die later from the two is Twenty Plus Two. So the
answer is: Twenty Plus Two.

Wikipedia Title: Ana Gruzinsky--Golitsyn

Wikipedia Title: Rudra Shah

Wikipedia Title: Krishna Shah (Nepalese royal)

Q: Who is the grandchild of Krishna Shah (Nepalese Royal)?

A: Krishna Shah has a child named Rudra Shah. Rudra Shah has a child
named Prithvipati Shah. Thus, Krishna Shah has a grandchild named
Prithvipati Shah. So the answer is: Prithvipati Shah.

Wikipedia Title: Augusto Genina

Wikipedia Title: Ian Barry (director)

Wikipedia Title: Maddalena (1954 film)

Q: Where did the director of film Maddalena (1954 Film) die?

A: The film Maddalena is directed by Augusto Genina. Augusto Genina died
in Rome. So the answer is: Rome.

Wikipedia Title: Martin Hodge

Wikipedia Title: Ivania Martinich

Wikipedia Title: Tom Dickinson

Q: Who was born first out of Martin Hodge and Ivania Martinich?

A: Martin Hodge was born on 4 February 1959. Ivania Martinich was born
on 25 July 1995. Thus, Martin Hodge was born first. So the answer is:
Martin Hodge.

\section{musique\_ans: decomp}\label{musique_ans-decomp}

QC: When did the first large winter carnival take place in the city
where CIMI--FM is licensed to broadcast?

QS: (select) {[}retrieve\_odqa{]} What city is CIMI--FM is licensed to
broadcast in?

A: \{``titles'': {[}``KWMZ--FM'', ``CIMI--FM'', ``KSAO (FM)'',
``WTTL--FM'', ``KOLU'', ``WORW'', ``WLRX (FM){]}``,''answer'':
{[}``Quebec City''{]}\}

QS: (select) {[}retrieve\_odqa{]} When did the first large winter
carnival in Quebec City take place?

A: \{``titles'': {[}``WRQY'', ``Quebec Winter Carnival'',
``KBCR-FM''{]}, ``answer'': {[}``1894''{]}\}

QS: (select) {[}multihop\_titleqa{]} Titles: {[}``KWMZ--FM'',
``CIMI--FM'', ``KSAO (FM)'', ``WTTL--FM'', ``KOLU'', ``WORW'', ``WLRX
(FM)'', ``WRQY'', ``Quebec Winter Carnival'', ``KBCR--FM''{]}. Question:
When did the first large winter carnival take place in the city where
CIMI--FM is licensed to broadcast?

A: {[}''1894''{]}

QS: {[}EOQ{]}

QC: What is the headquarters for the organization who sets the standards
for ISO 21500?

QS: (select) {[}retrieve\_odqa{]} Who sets the standards for ISO 21500?

A: \{``titles'': {[}``ISO 3166-2:BM'', ``ISO 3166-2:IS'', ``ISO
3166-2:IT'', ``ISO 3166-2:BB'', ``ISO 21500'', ``ISO 3166-2:ET'', ``ISO
3166-2:AT'', ``ISO 3166-2:CG'', ``ISO 22000''{]}, ``answer'':
{[}``International Organization for Standardization''{]}\}

QS: (select) {[}retrieve\_odqa{]} Where is the headquarters for
International Organization for Standardization?

A: \{``titles'': {[}``International Organization for Standardization'',
``ISO 4031'', ``ISO 3166-2:CN'', ``Unicode'', ``ISO 3166-2:GH'', ``ISO
3166-2:AO'', ``ISO/TC 68'', ``ISO 7001'', ``ISO 3307''{]}, ``answer'':
{[}``Geneva''{]}\}

QS: (select) {[}multihop\_titleqa{]} Titles: {[}``ISO 3166-2:BM'', ``ISO
3166-2:IS'', ``ISO 3166-2:IT'', ``ISO 3166-2:BB'', ``ISO 21500'', ``ISO
3166-2:ET'', ``ISO 3166-2:AT'', ``ISO 3166-2:CG'', ``ISO 22000'',
``International Organization for Standardization'', ``ISO 4031'', ``ISO
3166-2:CN'', ``Unicode'', ``ISO 3166-2:GH'', ``ISO 3166-2:AO'', ``ISO/TC
68'', ``ISO 7001'', ``ISO 3307''{]}. Question: What is the headquarters
for the organization who sets the standards for ISO 21500?

A: {[}``Geneva''{]}

QS: {[}EOQ{]}

QC: How long is the US border with the country that borders the state
where Finding Dory takes place?

QS: (select) {[}retrieve\_odqa{]} In which state does Finding Dory take
place?

A: \{``titles'': {[}``Finding Dory'', ``Latvia'', ``List of countries
that border only one other country'', ``Removal of Hungary's border
fence with Austria''{]}, ``answer'': {[}``California''{]}\}

QS: (select) {[}retrieve\_odqa{]} Which country shares a border with
California?

A: \{``titles'': {[}``Mexico--United States border'', ``Kingdom of
Gera'', ``Pesticide'', ``Currie Cup''{]}, ``answer'': {[}``Mexico''{]}\}

QS: (select) {[}retrieve\_odqa{]} What is the length of the US border
with Mexico?

A: \{``titles'': {[}``Piscataqua River border dispute'', ``Share a
Coke'', ``Mexico--United States border'', ``Hotel Arbez''{]},
``answer'': {[}``1,989 mi''{]}\}

QS: (select) {[}multihop\_titleqa{]} Titles: {[}``Finding Dory'',
``Latvia'', ``List of countries that border only one other country'',
``Removal of Hungary's border fence with Austria'', ``Mexico--United
States border'', ``Kingdom of Gera'', ``Pesticide'', ``Currie Cup'',
``Piscataqua River border dispute'', ``Share a Coke'', ``Mexico--United
States border'', ``Hotel Arbez''{]}. Question: How long is the US border
with the country that borders the state where Finding Dory takes place?

A: {[}``1,989 mi''{]}

QS: {[}EOQ{]}

QC: When was Neville A. Stanton's employer founded?

QS: (select) {[}retrieve\_odqa{]} Who is the employer of Neville A.
Stanton?

A: \{``titles'': {[}``Harriot Stanton Blatch'', ``Robichaud v Canada
(Treasury Board)'', ``The Peggy Neville Show'', ``Stanton Township,
Champaign County, Illinois'', ``Women's suffrage in the United States'',
``Neville A. Stanton''{]}, ``answer'': {[}``University of
Southampton''{]}\}

QS: (select) {[}retrieve\_odqa{]} When was the University of Southampton
founded?

A: \{``titles'': {[}``Southampton'', ``Presley Neville''{]}, ``answer'':
{[}``1862''{]}\}

QS: (select) {[}multihop\_titleqa{]} Titles: {[}``Harriot Stanton
Blatch'', ``Robichaud v Canada (Treasury Board)'', ``The Peggy Neville
Show'', ``Stanton Township, Champaign County, Illinois'', ``Women's
suffrage in the United States'', ``Neville A. Stanton'',
``Southampton'', ``Presley Neville''{]}. Question: When was Neville A.
Stanton's employer founded?

A: {[}``1862''{]}

QS: {[}EOQ{]}

QC: In which county was the birthplace of the Smoke in tha City
performer?

QS: (select) {[}retrieve\_odqa{]} Who is the performer of Smoke in the
City?

A: \{``titles'': {[}``Nicotine marketing'', ``Khabarovsky District'',
``Beyonce'', ``Asthma'', ``Smoke in tha City''{]}, ``answer'': {[}``MC
Eiht''{]}\}\\
QS: (select) {[}retrieve\_odqa{]} What is the birthplace of MC Eiht?\\
A: \{``titles'': {[}``Cherokee City, Arkansas'', ``MC Eiht'', ``Biysky
District''{]}, ``answer'': {[}``Compton''{]}\}\\
QS: (select) {[}retrieve\_odqa{]} Compton is located in which county?\\
A: \{''titles'': {[}``Vilnius County'', ``Biblioteca Ayacucho'',
``Compton, California''{]}, ``answer'': {[}``Los Angeles County''{]}\}\\
QS: (select) {[}multihop\_titleqa{]} Titles: {[}``Nicotine marketing'',
``Khabarovsky District'', ``Beyonce'', ``Asthma'', ``Smoke in tha
City'', ``Cherokee City, Arkansas'', ``MC Eiht'', ``Biysky District'',
``Vilnius County'', ``Biblioteca Ayacucho'', ``Compton, California''{]}.
Question: In which county was the birthplace of the Smoke in tha City
performer?\\
A: {[}``Los Angeles County''{]}\\
QS: {[}EOQ{]}

\section{musique\_ans: retrieve\_odqa}\label{musique_ans-retrieve_odqa}

QC: What city is CIMI--FM is licensed to broadcast in?

QS: (select) {[}retrieve{]} What city is CIMI--FM is licensed to
broadcast in?

A: {[}``KWMZ--FM'', ``CIMI--FM'', ``KSAO (FM)'', ``WTTL--FM'', ``KOLU'',
``WORW'', ``WLRX (FM)''{]}

QS: (select) {[}singlehop\_titleqa{]} Titles: {[}``KWMZ--FM'',
``CIMI--FM'', ``KSAO (FM)'', ``WTTL--FM'', ``KOLU'', ``WORW'', ``WLRX
(FM)''{]}. Question: What city is CIMI--FM is licensed to broadcast in?

A: \{ ``titles'': {[}``KWMZ--FM'', ``CIMI--FM'', ``KSAO (FM)'',
``WTTL--FM'', ``KOLU'', ``WORW'', ``WLRX (FM)''{]}, ``answer'':
{[}``Quebec City''{]}\}

QS: {[}EOQ{]}

QC: When did the first large winter carnival in Quebec City take place?

QS: (select) {[}retrieve{]} When did the first large winter carnival in
Quebec City take place?

A: {[}``WRQY'', ``Quebec Winter Carnival'', ``KBCR-FM''{]}

QS: (select) {[}singlehop\_titleqa{]} Titles: {[}``WRQY'', ``Quebec
Winter Carnival'', ``KBCR--FM''{]}. Question: When did the first large
winter carnival in Quebec City take place?

A: \{``titles'': {[}``WRQY'', ``Quebec Winter Carnival'',
``KBCR-FM''{]}, ``answer'': {[}``1894''{]}\}

QS: {[}EOQ{]}

QC: Who sets the standards for ISO 21500?

QS: (select) {[}retrieve{]} Who sets the standards for ISO 21500?

A: {[}``ISO 3166-2:BM'', ``ISO 3166-2:IS'', ``ISO 3166-2:IT'', ``ISO
3166-2:BB'', ``ISO 21500'', ``ISO 3166-2:ET'', ``ISO 3166-2:AT'', ``ISO
3166-2:CG'', ``ISO 22000''{]}

QS: (select) {[}singlehop\_titleqa{]} Titles: {[}``ISO 3166-2:BM'',
``ISO 3166-2:IS'', ``ISO 3166-2:IT'', ``ISO 3166-2:BB'', ``ISO 21500'',
``ISO 3166-2:ET'', ``ISO 3166-2:AT'', ``ISO 3166-2:CG'', ``ISO
22000''{]}. Question: Who sets the standards for ISO 21500?

A: \{''titles'': {[}``ISO 3166--2:BM'',''ISO 3166--2:IS'', ``ISO
3166--2:IT'', ``ISO 3166--2:BB'', ``ISO 21500'', ``ISO 3166--2:ET'',
``ISO 3166--2:AT'', ``ISO 3166--2:CG'', ``ISO 22000''{]}, ``answer'':
{[}``International Organization for Standardization''{]}\}

QS: {[}EOQ{]}

QC: Where is the headquarters for International Organization for
Standardization?

QS: (select) {[}retrieve{]} Where is the headquarters for International
Organization for Standardization?

A: {[}``International Organization for Standardization'', ``ISO 4031'',
``ISO 3166--2:CN'', ``Unicode'', ``ISO 3166--2:GH'', ``ISO 3166--2:AO'',
``ISO/TC 68'', ``ISO 7001'', ``ISO 3307''{]}

QS: (select) {[}singlehop\_titleqa{]} Titles: {[}``International
Organization for Standardization'', ``ISO 4031'', ``ISO 3166-2:CN'',
``Unicode'', ``ISO 3166-2:GH'', ``ISO 3166-2:AO'', ``ISO/TC 68'', ``ISO
7001'', ``ISO 3307''{]}. Question: Where is the headquarters for
International Organization for Standardization?

A: \{``titles'': {[}``International Organization for Standardization'',
``ISO 4031'', ``ISO 3166-2:CN'', ``Unicode'', ``ISO 3166-2:GH'', ``ISO
3166-2:AO'', ``ISO/TC 68'', ``ISO 7001'', ``ISO 3307''{]}, ``answer'':
{[}``Geneva''{]}\}

QS: {[}EOQ{]}

QC: In which state does Finding Dory take place?

QS: (select) {[}retrieve{]} In which state does Finding Dory take place?

A: {[}``Finding Dory'', ``Latvia'', ``List of countries that border only
one other country'', ``Removal of Hungary's border fence with
Austria''{]}

QS: (select) {[}singlehop\_titleqa{]} Titles: {[}``Finding Dory'',
``Latvia'', ``List of countries that border only one other country'',
``Removal of Hungary's border fence with Austria''{]}. Question: In
which state does Finding Dory take place?

A: \{``titles'': {[}``Finding Dory'', ``Latvia'', ``List of countries
that border only one other country'', ``Removal of Hungary's border
fence with Austria''{]}, ``answer'': {[}``California''{]}\}

QS: {[}EOQ{]}

QC: Which country shares a border with California?

QS: (select) {[}retrieve{]} Which country shares a border with
California?

A: {[}``Mexico--United States border'', ``Kingdom of Gera'',
``Pesticide'', ``Currie Cup''{]}

QS: (select) {[}singlehop\_titleqa{]} Titles: {[}``Mexico--United States
border'', ``Kingdom of Gera'', ``Pesticide'', ``Currie Cup''{]}.
Question: Which country shares a border with California?

A: \{``titles'': {[}``Mexico--United States border'', ``Kingdom of
Gera'', ``Pesticide'', ``Currie Cup''{]}, ``answer'': {[}``Mexico''{]}\}

QS: {[}EOQ{]}

QC: What is the length of the US border with Mexico?

QS: (select) {[}retrieve{]} What is the length of the US border with
Mexico?

A: {[}``Piscataqua River border dispute'', ``Share a Coke'',
``Mexico--United States border'', ``Hotel Arbez''{]}

QS: (select) {[}singlehop\_titleqa{]} Titles: {[}``Piscataqua River
border dispute'', ``Share a Coke'', ``Mexico--United States border'',
``Hotel Arbez''{]}. Question: What is the length of the US border with
Mexico?

A: \{``titles'': {[}``Piscataqua River border dispute'', ``Share a
Coke'', ``Mexico--United States border'', ``Hotel Arbez''{]},
``answer'': {[}``1,989 mi''{]}\}

QS: {[}EOQ{]}

QC: Who is the employer of Neville A. Stanton?

QS: (select) {[}retrieve{]} Who is the employer of Neville A. Stanton?

A: {[}''Harriot Stanton Blatch'', ''Robichaud v Canada (Treasury
Board)'', ''The Peggy Neville Show'', ''Stanton Township, Champaign
County, Illinois'', ''Women's suffrage in the United States'', ''Neville
A. Stanton''{]}

QS: (select) {[}singlehop\_titleqa{]} Titles: {[}``Harriot Stanton
Blatch'', ``Robichaud v Canada (Treasury Board)'', ``The Peggy Neville
Show'', ``Stanton Township, Champaign County, Illinois'', ``Women's
suffrage in the United States'', ``Neville A. Stanton''{]}. Question:
Who is the employer of Neville A. Stanton?

A: \{``titles'': {[}``Harriot Stanton Blatch'', ``Robichaud v Canada
(Treasury Board)'', ``The Peggy Neville Show'', ``Stanton Township,
Champaign County, Illinois'', ``Women's suffrage in the United States'',
``Neville A. Stanton''{]}, ``answer'': {[}``University of
Southampton''{]}\}

QS: {[}EOQ{]}

QC: When was the University of Southampton founded?

QS: (select) {[}retrieve{]} When was the University of Southampton
founded?

A: {[}``Southampton'', ``Presley Neville''{]}

QS: (select) {[}singlehop\_titleqa{]} Titles: {[}``Southampton'',
``Presley Neville''{]}. Question: When was the University of Southampton
founded?

A: \{``titles'': {[}``Southampton'', ``Presley Neville''{]}, ``answer'':
{[}``1862''{]}\}

QS: {[}EOQ{]}

QC: Who is the performer of Smoke in the City?

QS: (select) {[}retrieve{]} Who is the performer of Smoke in the City?

A: {[}``Nicotine marketing'', ``Khabarovsky District'', ``Beyonce'',
``Asthma'', ``Smoke in tha City''{]}

QS: (select) {[}singlehop\_titleqa{]} Titles: {[}``Nicotine marketing'',
``Khabarovsky District'', ``Beyonce'', ``Asthma'', ``Smoke in tha
City''{]}. Question: Who is the performer of Smoke in the City?

A: \{''titles'': {[}``Nicotine marketing'', ``Khabarovsky District'',
``Beyonce'', ``Asthma'', ``Smoke in tha City''{]}, ``answer'': {[}``MC
Eiht''{]}\}

QS: {[}EOQ{]}

QC: What is the birthplace of MC Eiht?

QS: (select) {[}retrieve{]} What is the birthplace of MC Eiht?

A: {[}``Cherokee City, Arkansas'', ``MC Eiht'', ``Biysky District''{]}

QS: (select) {[}singlehop\_titleqa{]} Titles: {[}``Cherokee City,
Arkansas'', ``MC Eiht'', ``Biysky District''{]}. Question: What is the
birthplace of MC Eiht?

A: \{``titles'': {[}``Cherokee City, Arkansas'', ``MC Eiht'', ``Biysky
District''{]}, ``answer'': {[}``Compton''{]}\}

QS: {[}EOQ{]}

QC: Compton is located in which county?

QS: (select) {[}retrieve{]} Compton is located in which county?

A: {[}``Vilnius County'', ``Biblioteca Ayacucho'', ``Compton,
California''{]}

QS: (select) {[}singlehop\_titleqa{]} Titles: {[}''Vilnius County'',
''Biblioteca Ayacucho'', ''Compton, California''{]}. Question: Compton
is located in which county?

A: \{''titles'': {[}``Vilnius County'', ``Biblioteca Ayacucho'',
``Compton, California''{]}, ``answer'': {[}``Los Angeles County''{]}\}

QS: {[}EOQ{]}

musique\_ans: singlehop\_titleqa

Wikipedia Title: CIMI--FM

Wikipedia Title: Quebec Winter Carnival

Wikipedia Title: KELD--FM

Q: What city is CIMI--FM is licensed to broadcast in?

A: {[}``Quebec City''{]}

Q: When did the first large winter carnival in Quebec City take place?

A: {[}``1894''{]}

Wikipedia Title: ISO/TC 68

Wikipedia Title: ISO 21500

Wikipedia Title: ISO 3166-2:GH

Q: Who sets the standards for ISO 21500?

A: {[}``International Organization for Standardization''{]}

Q: Where is the headquarters for International Organization for
Standardization?

A: {[}``Geneva''{]}

Wikipedia Title: Mexico--United States border

Wikipedia Title: Mexico--United States border

Wikipedia Title: Finding Dory

Wikipedia Title: India

Q: In which state does Finding Dory take place?

A: {[}``California''{]}

Q: Which country shares a border with California?

A: {[}``Mexico''{]}

Q: What is the length of the US border with Mexico?

A: {[}``1,989 mi''{]}

Wikipedia Title: Southampton

Wikipedia Title: Jonathan Stanton

Wikipedia Title: Neville A. Stanton

Q: Who is the employer of Neville A. Stanton?

A: {[}``University of Southampton''{]}

Q: When was the University of Southampton founded?

A: {[}''1862''{]}

Wikipedia Title: Compton, California

Wikipedia Title: Smoke in tha City

Wikipedia Title: MC Eiht

Wikipedia Title: Eastern Bengal and Assam

Q: Who is the performer of Smoke in the City?

\begin{Shaded}
\begin{Highlighting}[]
\NormalTok{A: ["MC Eiht"]}
\NormalTok{Q: What is the birthplace of MC Eiht?}
\NormalTok{A: ["Compton"]}
\NormalTok{Q: Compton is located in which county?}
\NormalTok{A: ["Los Angeles County"]}
\NormalTok{musique\_ans: multihop\_titleqa (direct)}
\NormalTok{Wikipedia Title: CIMI{-}FM}
\NormalTok{\textless{}hidden for brevity\textgreater{}}
\NormalTok{Wikipedia Title: Quebec Winter Carnival}
\NormalTok{\textless{}hidden for brevity\textgreater{}}
\NormalTok{Wikipedia Title: KELD{-}FM}
\NormalTok{\textless{}hidden for brevity\textgreater{}}
\NormalTok{Q: When did the first large winter carnival take place in the city where CIMI{-}FM is licensed to broadcast?}
\NormalTok{A: ["1894"]}
\NormalTok{Wikipedia Title: ISO/TC 68}
\NormalTok{\textless{}hidden for brevity\textgreater{}}
\NormalTok{Wikipedia Title: ISO 21500}
\NormalTok{\textless{}hidden for brevity\textgreater{}}
\NormalTok{Wikipedia Title: ISO 3166{-}2:GH}
\NormalTok{\textless{}hidden for brevity\textgreater{}}
\NormalTok{Q: What is the headquarters for the organization who sets the standards for ISO 21500?}
\NormalTok{A: ["Geneva"]}
\NormalTok{Wikipedia Title: Mexico{-}United States border}
\NormalTok{\textless{}hidden for brevity\textgreater{}}
\NormalTok{Wikipedia Title: Mexico{-}United States border}
\NormalTok{\textless{}hidden for brevity\textgreater{}}
\NormalTok{Wikipedia Title: Finding Dory}
\NormalTok{\textless{}hidden for brevity\textgreater{}}
\NormalTok{Wikipedia Title: India}
\NormalTok{\textless{}hidden for brevity\textgreater{}}
\NormalTok{Q: How long is the US border with the country that borders the state where Finding Dory takes place?}
\NormalTok{A: ["1,989 mi"]}
\NormalTok{Wikipedia Title: Southampton}
\NormalTok{\textless{}hidden for brevity\textgreater{}}
\NormalTok{Wikipedia Title: Jonathan Stanton}
\NormalTok{\textless{}hidden for brevity\textgreater{} }
\end{Highlighting}
\end{Shaded}

\begin{Shaded}
\begin{Highlighting}[]
\NormalTok{Wikipedia Title: Neville A. Stanton}
\NormalTok{\textless{}hidden for brevity\textgreater{}}
\NormalTok{Q: When was Neville A. Stanton\textquotesingle{}s employer founded?}
\NormalTok{A: ["1862"]}
\NormalTok{Wikipedia Title: Compton, California}
\NormalTok{\textless{}hidden for brevity\textgreater{}}
\NormalTok{Wikipedia Title: Smoke in tha City}
\NormalTok{\textless{}hidden for brevity\textgreater{}}
\NormalTok{Wikipedia Title: MC Eiht}
\NormalTok{\textless{}hidden for brevity\textgreater{}}
\NormalTok{Wikipedia Title: Eastern Bengal and Assam}
\NormalTok{\textless{}hidden for brevity\textgreater{}}
\NormalTok{Q: In which county was the birthplace of the Smoke in tha City performer?}
\NormalTok{A: ["Los Angeles County"] }
\end{Highlighting}
\end{Shaded}

\begin{Shaded}
\begin{Highlighting}[]
\NormalTok{musique\_ans: multihop\_titleqa (cot) }
\end{Highlighting}
\end{Shaded}

\begin{Shaded}
\begin{Highlighting}[]
\NormalTok{Wikipedia Title: CIMI–FM}
\NormalTok{\textless{}hidden for brevity\textgreater{} }
\end{Highlighting}
\end{Shaded}

\begin{Shaded}
\begin{Highlighting}[]
\NormalTok{Wikipedia Title: Quebec Winter Carnival}
\NormalTok{\textless{}hidden for brevity\textgreater{} }
\end{Highlighting}
\end{Shaded}

\begin{Shaded}
\begin{Highlighting}[]
\NormalTok{Wikipedia Title: KELD–FM}
\NormalTok{\textless{}hidden for brevity\textgreater{} }
\end{Highlighting}
\end{Shaded}

\begin{Shaded}
\begin{Highlighting}[]
\NormalTok{Q: When did the first large winter carnival take place in the city where CIMI–FM is licensed to broadcast?}
\NormalTok{A: CIMI–FM is licensed to broadcast in Quebec City. The first large winter carnival in Quebec City took place in 1894. So the answer is: 1894. }
\end{Highlighting}
\end{Shaded}

\begin{Shaded}
\begin{Highlighting}[]
\NormalTok{Wikipedia Title: ISO/TC 68}
\NormalTok{\textless{}hidden for brevity\textgreater{} }
\end{Highlighting}
\end{Shaded}

\begin{Shaded}
\begin{Highlighting}[]
\NormalTok{Wikipedia Title: ISO 21500}
\NormalTok{\textless{}hidden for brevity\textgreater{} }
\end{Highlighting}
\end{Shaded}

\begin{Shaded}
\begin{Highlighting}[]
\NormalTok{Wikipedia Title: ISO 3166{-}2:GH}
\NormalTok{\textless{}hidden for brevity\textgreater{} }
\end{Highlighting}
\end{Shaded}

\begin{Shaded}
\begin{Highlighting}[]
\NormalTok{Q: What is the headquarters for the organization who sets the standards for ISO 21500?}
\NormalTok{A: The standards for ISO 21500 were set by International Organization for Standardization. The International Organization for Standardization has headquarters in Geneva. So the answer is: Geneva. }
\end{Highlighting}
\end{Shaded}

\begin{Shaded}
\begin{Highlighting}[]
\NormalTok{Wikipedia Title: Mexico–United States border \textless{}hidden for brevity\textgreater{} }
\end{Highlighting}
\end{Shaded}

\begin{Shaded}
\begin{Highlighting}[]
\NormalTok{Wikipedia Title: Mexico–United States border \textless{}hidden for brevity\textgreater{} }
\end{Highlighting}
\end{Shaded}

Wikipedia Title: Finding Dory

Wikipedia Title: India

Q: How long is the US border with the country that borders the state
where Finding Dory takes place?

A: Finding Dory is supposed to take place in California. The country
that shares a border with California is Mexico. The length of the us
border with Mexico is 1,989 mi. So the answer is: 1,989 mi.

Wikipedia Title: Southampton

Wikipedia Title: Jonathan Stanton

Wikipedia Title: Neville A. Stanton

Q: When was Neville A. Stanton's employer founded?

A: The employer of Neville A. Stanton is University of Southampton. The
University of Southampton was founded in 1862. So the answer is: 1862.

Wikipedia Title: Compton, California

Wikipedia Title: Smoke in tha City

Wikipedia Title: MC Eiht

Wikipedia Title: Eastern Bengal and Assam

Q: In which county was the birthplace of the Smoke in tha City
performer?

A: The performer of Smoke in the City is MC Eiht. MC Eiht's birthplace
is Compton. Compton is located in the county of Los Angeles County. So
the answer is: Los Angeles County.

\end{document}
